\input{preamble}

% OK, start here.
%
\begin{document}

\title{Cohomologie pro-\'etale}


\maketitle

\phantomsection
\label{section-phantom}

\tableofcontents

\section{Introduction}
\label{section-introduction}

\noindent
Le contenu de ce chapitre, et davantage encore, se trouve dans la pr\'epublication \cite{BS}.

\medskip\noindent
Le but de ce chapitre est d'introduire la topologie pro-\'etale et de
d\'evelopper les fondements de la cohomologie des faisceaux ab\'eliens dans cette topologie.
Un objectif secondaire est de montrer comment l'emploi de la topologie pro-\'etale simplifie
l'introduction de la cohomologie $\ell$-adique en g\'eom\'etrie alg\'ebrique.

\medskip\noindent
Voici un bref aper\c{c}u de l'histoire de la cohomologie \'etale $\ell$-adique
telle que nous l'avons comprise.
Dans \cite[Expos\'es V et VI]{SGA5}, Grothendieck et ses collaborateurs ont d\'evelopp\'e une th\'eorie
qui traite les faisceaux $\ell$-adiques comme des syst\`emes projectifs de faisceaux de
$\mathbf{Z}/\ell^n\mathbf{Z}$-modules.
Dans son second article sur les conjectures de Weil (\cite{WeilII}), Deligne a introduit
une cat\'egorie d\'eriv\'ee de faisceaux $\ell$-adiques comme une certaine 2-limite de cat\'egories
de complexes de faisceaux de $\mathbf{Z}/\ell^n\mathbf{Z}$-modules sur le
site \'etale d'un sch\'ema $X$. Cette approche sert, dans l'article de
Beilinson, Bernstein et Deligne (\cite{BBD}), de fondement \`a leur
belle th\'eorie des faisceaux pervers. Dans un article intitul\'e ``Cohomologie
\'etale continue'' (\cite{Jannsen}), Uwe Jannsen examine une variante importante
de la cohomologie d'un faisceau $\ell$-adique sur une vari\'et\'e sur un corps.
Son article est prolong\'e par un article de Torsten Ekedahl (\cite{Ekedahl}),
qui \`etudie le formalisme adique n\'ecessaire pour manier commod\'ement les cat\'egories
d\'eriv\'ees d\'efinies comme des limites.

\medskip\noindent
Il se trouve qu'en travaillant avec le site pro-\'etale d'un sch\'ema,
on peut \`eviter certaines des difficult\'es techniques rencontr\'ees par ces auteurs.
Cela oblige en contrepartie \`a travailler avec des sch\'emas non noeth\'eriens,
m\^eme lorsque l'on ne s'int\'eresse qu'aux faisceaux $\ell$-adiques
et \`a leur cohomologie sur des vari\'et\'es d\'efinies sur un corps alg\'ebriquement clos.

\medskip\noindent
Une propri\'et\'e tr\`es importante et remarquable du
(petit) site pro-\'etale d'un sch\'ema est
qu'il poss\`ede assez d'objets quasi-compacts w-contractiles. L'existence
de ces objets entra\^ine
plusieurs cons\'equences utiles et (peut-\^etre) inhabituelles pour la cat\'egorie d\'eriv\'ee
des faisceaux ab\'eliens et pour les syst\`emes projectifs de faisceaux.
C'est pr\'ecis\'ement cette propri\'et\'e qui nous permettra de traiter
les subtilit\'es propres aux faisceaux $\ell$-adiques; mais, comme nous
le verrons, elle pr\'esente aussi plusieurs autres avantages.




\section{Quelques \`el\'ements de topologie}
\label{section-topology}

\noindent
Quelques pr\'eliminaires. Nous avons d\'efini les {\it espaces spectraux} et les
{\it applications spectrales} entre espaces spectraux dans
Topologie, section \ref{topology-section-spectral}.
Le spectre d'un anneau est un espace spectral; voir
Alg\`ebre, lemme \ref{algebra-lemma-spec-spectral}.

\begin{lemma}
\label{lemma-spectral-split}
Soit $X$ un espace spectral. Soit $X_0 \subset X$ l'ensemble des points ferm\'es.
Les conditions suivantes sont \`equivalentes:
\begin{enumerate}
\item Tout recouvrement ouvert de $X$ peut \^etre raffin\'e par une d\'ecomposition finie
en union disjointe $X = \coprod U_i$, o\`u les $U_i$ sont
ouverts et ferm\'es dans $X$.
\item L'application compos\'ee $X_0 \to X \to \pi_0(X)$ est bijective.
\end{enumerate}
De plus, si $X_0$ est ferm\'e dans $X$ et si tout point de $X$ se sp\'ecialise
en un unique point de $X_0$, alors ces conditions sont satisfaites.
\end{lemma}

\begin{proof}
Nous utiliserons sans autre mention le fait que
$X_0$ est quasi-compact
(Topologie, lemme \ref{topology-lemma-closed-points-quasi-compact})
et que $\pi_0(X)$ est profini
(Topologie, lemme \ref{topology-lemma-spectral-pi0}).
Consid\'erons le diagramme
$$
\xymatrix{
X_0 \ar[rd]_f \ar[r] & X \ar[d]^\pi \\
& \pi_0(X)
}
$$
Si (2) est satisfaite, l'application continue bijective $f : X_0 \to \pi_0(X)$ est
un hom\'eomorphisme d'apr\`es
Topologie, lemme \ref{topology-lemma-bijective-map}.
\'Etant donn\'e un recouvrement ouvert $X = \bigcup U_i$, on obtient le recouvrement ouvert
$\pi_0(X) = \bigcup f(X_0 \cap U_i)$. D'apr\`es
Topologie, lemme \ref{topology-lemma-profinite-refine-open-covering},
il existe un recouvrement ouvert fini de la forme $\pi_0(X) = \coprod V_j$
qui raffine ce recouvrement.
Puisque $X_0 \to \pi_0(X)$ est bijective, chaque composante connexe de
$X$ poss\`ede un unique point ferm\'e et est donc l'ensemble des points qui se
sp\'ecialisent en ce point ferm\'e. Ainsi $\pi^{-1}(V_j)$ est
l'ensemble des points qui se sp\'ecialisent en des points de $f^{-1}(V_j)$.
Or, si $f^{-1}(V_j) \subset X_0 \cap U_i \subset U_i$, alors
$\pi^{-1}(V_j) \subset U_i$ (car l'ouvert
$U_i$ est stable par g\'en\'erisation). On voit donc
que le recouvrement ouvert $X = \coprod \pi^{-1}(V_j)$ raffine
le recouvrement de d\'epart. Cela montre que
(2) implique (1).

\medskip\noindent
Supposons (1). Soient $x, y \in X$ des points ferm\'es. On a alors le recouvrement ouvert
$X = (X \setminus \{x\}) \cup (X \setminus \{y\})$.
Il r\'esulte de (1) qu'il existe une d\'ecomposition en union disjointe
$X = U \amalg V$, avec $U$ et $V$ ouverts (et ferm\'es), $x \in U$ et
$y \in V$. En particulier, toute composante connexe de $X$
poss\`ede au plus un point ferm\'e. D'apr\`es
Topologie, lemme \ref{topology-lemma-quasi-compact-closed-point},
toute composante connexe (qui est ferm\'ee) poss\`ede aussi un point ferm\'e.
Ainsi $X_0 \to \pi_0(X)$ est bijective. Cela montre que (1) implique (2).

\medskip\noindent
Supposons que $X_0$ soit ferm\'e dans $X$ et que tout point se sp\'ecialise en un unique
point de $X_0$. Alors $X_0$ est un espace spectral
(Topologie, lemme \ref{topology-lemma-spectral-sub})
constitu\'e de points ferm\'es, donc profini
(Topologie, lemme \ref{topology-lemma-characterize-profinite-spectral}).
Soient $x, y \in X_0$ distincts. D'apr\`es
Topologie, lemme \ref{topology-lemma-profinite-refine-open-covering},
il existe une d\'ecomposition en union disjointe
$X_0 = U_0 \amalg V_0$, o\`u $U_0$ et $V_0$ sont ouverts et ferm\'es,
avec $x \in U_0$ et $y \in V_0$.
Soient $U \subset X$, resp.\ $V \subset X$, les ensembles form\'es des points
qui se sp\'ecialisent en un point de $U_0$, resp.\ de $V_0$.
On a $X = U \amalg V$.
D'apr\`es Topologie, lemme \ref{topology-lemma-make-spectral-space},
$U$ est une intersection d'ouverts quasi-compacts.
Ainsi $U$ est ferm\'e pour la topologie constructible.
Puisque $U$ est stable par sp\'ecialisation,
$U$ est ferm\'e d'apr\`es Topologie, lemme
\ref{topology-lemma-constructible-stable-specialization-closed}.
Par sym\'etrie, $V$ est ferm\'e; ainsi $U$ et $V$ sont tous deux
ouverts et ferm\'es.
Cela prouve que $x, y$ n'appartiennent pas \`a la m\^eme composante connexe de $X$.
Autrement dit, $X_0 \to \pi_0(X)$ est injective. Elle est aussi
surjective d'apr\`es
Topologie, lemme \ref{topology-lemma-quasi-compact-closed-point}
et le fait que les composantes connexes sont ferm\'ees.
On voit ainsi que la derni\`ere condition implique (2).
\end{proof}

\begin{example}
\label{example-not-w-local}
Soit $T$ un espace profini. Soit $t \in T$ un point et supposons
que $T \setminus \{t\}$ ne soit pas quasi-compact.
Posons $X = T \times \{0, 1\}$. Munissons $X$ de la topologie
dont une pr\'ebase est donn\'ee par les ensembles
$U \times \{0, 1\}$ pour $U \subset T$ ouvert, $X \setminus \{(t, 0)\}$,
et $U \times \{1\}$ pour $U \subset T$ ouvert tel que $t \not \in U$.
L'ensemble des points ferm\'es de $X$ est $X_0 = T \times \{0\}$, et
$(t, 1)$ appartient \`a l'adh\'erence de $X_0$.
De plus, $X_0 \to \pi_0(X)$ est bijective.
Cet exemple montre que les conditions (1) et (2) du
lemme \ref{lemma-spectral-split} n'impliquent pas que l'ensemble des points ferm\'es
soit ferm\'e.
\end{example}

\noindent
Il est en fait plus commode de travailler avec des espaces
spectraux qui poss\`edent la propri\'et\'e un peu plus forte mentionn\'ee dans
le dernier \`enonc\'e du lemme \ref{lemma-spectral-split}.
Nous donnons un nom \`a cette propri\'et\'e.

\begin{definition}
\label{definition-w-local}
Un espace spectral $X$ est dit {\it w-local} si l'ensemble $X_0$ de ses points ferm\'es
est ferm\'e et si tout point de $X$ se sp\'ecialise en un unique point ferm\'e.
Une application continue $f : X \to Y$ entre espaces w-locaux est dite {\it w-locale}
si elle est spectrale et envoie tout point ferm\'e de $X$ sur un point ferm\'e de $Y$.
\end{definition}

\noindent
Nous avons vu dans la d\'emonstration du lemme \ref{lemma-spectral-split}
que, dans ce cas, $X_0 \to \pi_0(X)$ est un hom\'eomorphisme et que
$X_0 \cong \pi_0(X)$ est un espace profini. De plus, une composante
connexe de $X$ est exactement l'ensemble des points qui se sp\'ecialisent en
un $x \in X_0$ donn\'e.

\begin{lemma}
\label{lemma-closed-subspace-w-local}
Soit $X$ un espace spectral w-local. Si $Y \subset X$ est ferm\'e,
alors $Y$ est w-local.
\end{lemma}

\begin{proof}
Le sous-ensemble $Y_0 \subset Y$ des points ferm\'es est ferm\'e puisque
$Y_0 = X_0 \cap Y$. Comme $X$ est $w$-local, tout $y \in Y$ se sp\'ecialise
en un unique point de $X_0$. Ce point appartient \`a $Y$, et donc
aussi \`a $Y_0$, puisque $\overline{\{y\}}\subset Y$. En conclusion, $Y$
est $w$-local.
\end{proof}

\begin{lemma}
\label{lemma-silly}
Soit $X$ un espace spectral. Supposons que
$$
\xymatrix{
Y \ar[r] \ar[d] & T \ar[d] \\
X \ar[r] & \pi_0(X)
}
$$
soit un diagramme cart\'esien dans la cat\'egorie des espaces topologiques,
o\`u $T$ est profini. Alors $Y$ est spectral et $T = \pi_0(Y)$.
Si de plus $X$ est w-local, alors $Y$ est w-local, $Y \to X$ est w-locale,
et l'ensemble des points ferm\'es de $Y$ est l'image r\'eciproque de
l'ensemble des points ferm\'es de $X$.
\end{lemma}

\begin{proof}
Notons que $Y$ est un sous-espace ferm\'e de $X \times T$, car $\pi_0(X)$
est un espace profini, donc s\'epar\'e
(utiliser Topologie, lemmes \ref{topology-lemma-spectral-pi0} et
\ref{topology-lemma-fibre-product-closed}).
Comme $X \times T$ est spectral
(Topologie, lemme \ref{topology-lemma-product-spectral-spaces}),
il s'ensuit que $Y$ est spectral
(Topologie, lemme \ref{topology-lemma-spectral-sub}).
Consid\'erons la factorisation canonique $Y \to \pi_0(Y) \to T$
(Topologie, lemme \ref{topology-lemma-space-connected-components}).
Il est clair que $\pi_0(Y) \to T$ est surjective.
Les fibres de $Y \to T$ sont hom\'eomorphes aux fibres de
$X \to \pi_0(X)$. Elles sont donc connexes. Il s'ensuit
que $\pi_0(Y) \to T$ est injective. On en conclut que $\pi_0(Y) \to T$
est un hom\'eomorphisme d'apr\`es
Topologie, lemme \ref{topology-lemma-bijective-map}.

\medskip\noindent
Supposons ensuite $X$ w-local et notons $X_0 \subset X$
l'ensemble des points ferm\'es. L'image r\'eciproque $Y_0 \subset Y$ de $X_0$ dans
$Y$ s'identifie à $T$ par une bijection, car $X_0 \to \pi_0(X)$ est bijective
d'apr\`es le lemme \ref{lemma-spectral-split}. De plus, $Y_0$ est quasi-compact
comme sous-ensemble ferm\'e de l'espace spectral $Y$. Par cons\'equent,
$Y_0 \to \pi_0(Y) = T$ est un hom\'eomorphisme d'apr\`es
Topologie, lemme \ref{topology-lemma-bijective-map}.
Il s'ensuit que tous les points de $Y_0$ sont ferm\'es dans $Y$.
R\'eciproquement, si $y \in Y$ est un point ferm\'e, alors
il est ferm\'e dans la fibre de $Y \to \pi_0(Y) = T$;
son image $x$ dans $X$ est donc ferm\'e dans la fibre (hom\'eomorphe) de
$X \to \pi_0(X)$. Cela implique $x \in X_0$, puis $y \in Y_0$.
Ainsi $Y_0$ est l'ensemble des points ferm\'es de $Y$,
et, pour tout $y \in Y_0$, l'ensemble des g\'en\'erisations de $y$ est
la fibre de $Y \to \pi_0(Y)$. Le lemme en r\'esulte.
\end{proof}




\section{Isomorphismes locaux}
\label{section-local-isomorphism}

\noindent
Commen\c{c}ons par une d\'efinition.

\begin{definition}
\label{definition-local-isomorphism}
Soit $\varphi : A \to B$ un homomorphisme d'anneaux.
\begin{enumerate}
\item On dit que $A \to B$ est un {\it isomorphisme local} si, pour tout id\'eal premier
$\mathfrak q \subset B$, il existe $g \in B$, $g \not \in \mathfrak q$,
tel que $A \to B_g$ induise une immersion ouverte $\Spec(B_g) \to \Spec(A)$.
\item On dit que $A \to B$ {\it identifie les anneaux locaux} si, pour tout id\'eal premier
$\mathfrak q \subset B$, l'application canonique
$A_{\varphi^{-1}(\mathfrak q)} \to B_\mathfrak q$ est un isomorphisme.
\end{enumerate}
\end{definition}

\noindent
\'Enum\'erons quelques propri\'et\'es \`el\'ementaires.

\begin{lemma}
\label{lemma-base-change-local-isomorphism}
Soient $A \to B$ et $A \to A'$ des homomorphismes d'anneaux. Notons $B' = B \otimes_A A'$,
l'alg\`ebre obtenue \`a partir de $B$ par changement de base.
\begin{enumerate}
\item Si $A \to B$ est un isomorphisme local, alors $A' \to B'$ est un
isomorphisme local.
\item Si $A \to B$ identifie les anneaux locaux, alors $A' \to B'$
identifie les anneaux locaux.
\end{enumerate}
\end{lemma}

\begin{proof}
La d\'emonstration est omise.
\end{proof}

\begin{lemma}
\label{lemma-composition-local-isomorphism}
Soient $A \to B$ et $B \to C$ des homomorphismes d'anneaux.
\begin{enumerate}
\item Si $A \to B$ et $B \to C$ sont des isomorphismes locaux, alors $A \to C$
est un isomorphisme local.
\item Si $A \to B$ et $B \to C$ identifient les anneaux locaux, alors $A \to C$
identifie les anneaux locaux.
\end{enumerate}
\end{lemma}

\begin{proof}
La d\'emonstration est omise.
\end{proof}

\begin{lemma}
\label{lemma-local-isomorphism-permanence}
Soit $A$ un anneau. Soit $B \to C$ un homomorphisme de $A$-alg\`ebres.
\begin{enumerate}
\item Si $A \to B$ et $A \to C$ sont des isomorphismes locaux, alors $B \to C$
est un isomorphisme local.
\item Si $A \to B$ et $A \to C$ identifient les anneaux locaux, alors $B \to C$
identifie les anneaux locaux.
\end{enumerate}
\end{lemma}

\begin{proof}
La d\'emonstration est omise.
\end{proof}

\begin{lemma}
\label{lemma-local-isomorphism-implies}
Soit $A \to B$ un isomorphisme local. Alors
\begin{enumerate}
\item $A \to B$ est \'etale,
\item $A \to B$ identifie les anneaux locaux,
\item $A \to B$ est quasi-fini.
\end{enumerate}
\end{lemma}

\begin{proof}
La d\'emonstration est omise.
\end{proof}

\begin{lemma}
\label{lemma-structure-local-isomorphism}
Soit $A \to B$ un isomorphisme local. Il existe alors $n \geq 0$,
$g_1, \ldots, g_n \in B$ et $f_1, \ldots, f_n \in A$ tels que
$(g_1, \ldots, g_n) = B$ et $A_{f_i} \cong B_{g_i}$.
\end{lemma}

\begin{proof}
La d\'emonstration est omise.
\end{proof}

\begin{lemma}
\label{lemma-fully-faithful-spaces-over-X}
Soient $p : (Y, \mathcal{O}_Y) \to (X, \mathcal{O}_X)$ et
$q : (Z, \mathcal{O}_Z) \to (X, \mathcal{O}_X)$
des morphismes d'espaces localement annel\'es.
Si $\mathcal{O}_Y = p^{-1}\mathcal{O}_X$, alors
$$
\Mor_{\text{LRS}/(X, \mathcal{O}_X)}((Z, \mathcal{O}_Z), (Y, \mathcal{O}_Y))
\longrightarrow
\Mor_{\textit{Top}/X}(Z, Y),\quad
(f, f^\sharp) \longmapsto f
$$
est bijective. Ici $\text{LRS}/(X, \mathcal{O}_X)$ est la cat\'egorie des
espaces localement annel\'es au-dessus de $X$, et $\textit{Top}/X$ la cat\'egorie
des espaces topologiques au-dessus de $X$.
\end{lemma}

\begin{proof}
Cela r\'esulte imm\'ediatement des d\'efinitions.
\end{proof}

\begin{lemma}
\label{lemma-local-isomorphism-fully-faithful}
Soit $A$ un anneau. Posons $X = \Spec(A)$. Le foncteur
$$
B \longmapsto \Spec(B)
$$
de la cat\'egorie des $A$-alg\`ebres $B$ telles que $A \to B$ identifie
les anneaux locaux vers la cat\'egorie des
espaces topologiques au-dessus de $X$ est pleinement fid\`ele.
\end{lemma}

\begin{proof}
Cela r\'esulte du lemme \ref{lemma-fully-faithful-spaces-over-X}
et du fait que, si $A \to B$ identifie les anneaux locaux, l'image inverse
du faisceau structural de $\Spec(A)$ par $p : \Spec(B) \to \Spec(A)$
est \`egale au faisceau structural de $\Spec(B)$.
\end{proof}




\section{Alg\`ebres ind-Zariski}
\label{section-ind-zariski}

\noindent
Commen\c{c}ons par une d\'efinition; voir la remarque \ref{remark-slightly-stronger}
pour une comparaison avec la d\'efinition correspondante donn\'ee dans l'article \cite{BS}.

\begin{definition}
\label{definition-ind-zariski}
Un homomorphisme d'anneaux $A \to B$ est dit {\it ind-Zariski} si $B$ peut s'\'ecrire
comme une limite inductive filtrante $B = \colim B_i$, chaque $A \to B_i$ \`etant un
isomorphisme local.
\end{definition}

\noindent
Une localisation $A \to S^{-1}A$ est un exemple d'homomorphisme ind-Zariski; voir
Alg\`ebre, lemme \ref{algebra-lemma-localization-colimit}.
La cat\'egorie des alg\`ebres ind-Zariski est stable par plusieurs op\'erations
naturelles.

\begin{lemma}
\label{lemma-base-change-ind-zariski}
Soient $A \to B$ et $A \to A'$ des homomorphismes d'anneaux. Notons $B' = B \otimes_A A'$,
l'alg\`ebre obtenue \`a partir de $B$ par changement de base.
Si $A \to B$ est ind-Zariski, alors $A' \to B'$ est ind-Zariski.
\end{lemma}

\begin{proof}
La d\'emonstration est omise.
\end{proof}

\begin{lemma}
\label{lemma-composition-ind-zariski}
Soient $A \to B$ et $B \to C$ des homomorphismes d'anneaux. Si $A \to B$ et $B \to C$
sont ind-Zariski, alors $A \to C$ est ind-Zariski.
\end{lemma}

\begin{proof}
La d\'emonstration est omise.
\end{proof}

\begin{lemma}
\label{lemma-ind-zariski-permanence}
Soit $A$ un anneau. Soit $B \to C$ un homomorphisme de $A$-alg\`ebres.
Si $A \to B$ et $A \to C$ sont ind-Zariski, alors $B \to C$
est ind-Zariski.
\end{lemma}

\begin{proof}
La d\'emonstration est omise.
\end{proof}

\begin{lemma}
\label{lemma-ind-ind-zariski}
Une limite inductive filtrante de $A$-alg\`ebres ind-Zariski est ind-Zariski sur $A$.
\end{lemma}

\begin{proof}
La d\'emonstration est omise.
\end{proof}

\begin{lemma}
\label{lemma-ind-zariski-implies}
Soit $A \to B$ ind-Zariski. Alors $A \to B$ identifie les anneaux locaux.
\end{lemma}

\begin{proof}
La d\'emonstration est omise.
\end{proof}







\section{Construction de sch\'emas affines w-locaux}
\label{section-construction}

\noindent
Un sch\'ema affine $X$ est dit {\it w-local} si son espace topologique
sous-jacent est w-local (d\'efinition \ref{definition-w-local}).
Pour tout anneau $A$, il existe en fait un homomorphisme d'anneaux canonique, ind-Zariski
et fid\`element plat, $A \to A_w$, tel que $\Spec(A_w)$ soit
w-local. La construction de $A_w$ repose sur le lemme simple suivant.

\begin{lemma}
\label{lemma-localization}
Soit $A$ un anneau. Posons $X = \Spec(A)$. Soit $Z \subset X$ un sous-sch\'ema localement
ferm\'e de la forme $D(f) \cap V(I)$ pour un $f \in A$ et un
id\'eal $I \subset A$. Alors
\begin{enumerate}
\item il existe une partie multiplicative $S \subset A$ telle que
$\Spec(S^{-1}A)$ s'identifie par un hom\'eomorphisme \`a l'ensemble des points de $X$
qui se sp\'ecialisent en un point de $Z$,
\item la $A$-alg\`ebre $A_Z^\sim = S^{-1}A$ ne d\'epend que du
sous-ensemble localement ferm\'e sous-jacent $Z \subset X$,
\item $Z$ est un sous-sch\'ema ferm\'e de $\Spec(A_Z^\sim)$,
\end{enumerate}
Si $A \to A'$ est un homomorphisme d'anneaux et si $Z' \subset X' = \Spec(A')$ est un
sous-sch\'ema localement ferm\'e de la m\^eme forme dont l'image est contenue dans $Z$,
alors il existe un unique homomorphisme de $A$-alg\`ebres
$A_Z^\sim \to (A')_{Z'}^\sim$.
\end{lemma}

\begin{proof}
Soit $S \subset A$ la partie multiplicative form\'ee des \`el\'ements qui deviennent
inversibles dans $\Gamma(Z, \mathcal{O}_Z) = (A/I)_f$.
Si $\mathfrak p$ est un id\'eal premier de $A$ qui ne se sp\'ecialise pas en un point de $Z$,
alors $\mathfrak p$ engendre l'id\'eal unit\'e dans $(A/I)_f$. On peut donc
\'ecrire $f^n =  g + h$ pour un $n \geq 0$, avec $g \in \mathfrak p$ et
$h \in I$. Alors $g \in S$, et $\mathfrak p$ n'appartient pas
au spectre de $S^{-1}A$. R\'eciproquement, si $\mathfrak p$ se sp\'ecialise
en un point de $Z$, disons $\mathfrak p \subset \mathfrak q \supset I$ avec
$f \not \in \mathfrak q$, alors il existe un homomorphisme de $S^{-1}A$ vers
$A_\mathfrak q$; par cons\'equent, $\mathfrak p$ appartient au spectre de $S^{-1}A$.
Ceci prouve (1).

\medskip\noindent
La classe d'isomorphisme de la localisation $S^{-1}A$ ne d\'epend que
du sous-ensemble correspondant $\Spec(S^{-1}A) \subset \Spec(A)$, d'o\`u
(2). Par construction, l'application surjective de $S^{-1}A$ vers
$(A/I)_f$, d'o\`u (3). La derni\`ere assertion r\'esulte de ce que la partie multiplicative
$S' \subset A'$ correspondant \`a $Z'$ contient l'image de la
partie multiplicative $S$.
\end{proof}

\noindent
Soit $A$ un anneau. Soit $E \subset A$ une partie finie. On obtient une
stratification de $X = \Spec(A)$ en sous-sch\'emas localement ferm\'es en
distinguant, pour chaque \`el\'ement de $E$, s'il s'annule ou non. Plus pr\'ecis\'ement,
pour toute d\'ecomposition en union disjointe $E = E' \amalg E''$, posons
\begin{equation}
\label{equation-stratum}
Z(E', E'') =
\bigcap\nolimits_{f \in E'} D(f) \cap \bigcap\nolimits_{f \in E''} V(f) =
D(\prod\nolimits_{f \in E'} f) \cap V( \sum\nolimits_{f \in E''} fA)
\end{equation}
Les points de $Z(E', E'')$ sont exactement les $x \in X$ tels que
tout $f \in E'$ ait une image non nulle dans $\kappa(x)$ et que tout $f \in E''$
ait une image nulle dans $\kappa(x)$. Il est donc clair que
\begin{equation}
\label{equation-stratify}
X = \coprod\nolimits_{E = E' \amalg E''} Z(E', E'')
\end{equation}
en tant qu'ensembles. Notons que chaque strate est constructible.

\begin{lemma}
\label{lemma-refine}
Soit $X = \Spec(A)$ comme ci-dessus. Pour toute stratification finie
$X = \coprod T_i$ par des sous-ensembles constructibles, il existe une partie finie
$E \subset A$ telle que la stratification (\ref{equation-stratify})
raffine $X = \coprod T_i$.
\end{lemma}

\begin{proof}
On peut \`ecrire $T_i = \bigcup_j U_{i, j} \cap V_{i, j}^c$ comme une union finie,
o\`u les $U_{i, j}$ et $V_{i, j}$ sont des ouverts quasi-compacts de $X$.
Puis on peut \`ecrire $U_{i, j} = \bigcup D(f_{i, j, k})$ et
$V_{i, j} = \bigcup D(g_{i, j, l})$. Posons alors
$E = \{f_{i, j, k}\} \cup \{g_{i, j, l}\}$. Cela convient, car
la stratification (\ref{equation-stratify}) est celle dont les strates sont
index\'ees par les conditions d'annulation ou de non-annulation des \`el\'ements de $E$; elle
raffine manifestement la stratification donn\'ee.
\end{proof}

\noindent
Poursuivons la discussion.
Pour toute partie finie $E \subset A$, posons
\begin{equation}
\label{equation-ring}
A_E = \prod\nolimits_{E = E' \amalg E''} A_{Z(E', E'')}^\sim
\end{equation}
avec les notations du lemme \ref{lemma-localization}. Cette expression est bien d\'efinie, car
(\ref{equation-stratum}) montre que chaque $Z(E', E'')$ est de la forme voulue.
Prenons le spectre de cet anneau et notons-le
\begin{equation}
\label{equation-spectrum}
X_E = \Spec(A_E) = \coprod\nolimits_{E = E' \amalg E''} X_{E', E''}
\end{equation}
o\`u $X_{E', E''} = \Spec(A_{Z(E', E'')}^\sim)$. Notons que
\begin{equation}
\label{equation-closed}
Z_E = \coprod\nolimits_{E = E' \amalg E''} Z(E', E'')
\longrightarrow
X_E
\end{equation}
est un sous-sch\'ema ferm\'e. Par construction, le sous-sch\'ema ferm\'e $Z_E$
contient tous les points ferm\'es du sch\'ema affine $X_E$, puisque tout point
de $X_{E', E''}$ se sp\'ecialise en un point de $Z(E', E'')$.

\medskip\noindent
Soit $I(A)$ l'ensemble ordonn\'e de toutes les parties finies de $A$.
C'est un ensemble ordonn\'e filtrant. Pour $E_1 \subset E_2$, il existe
un homomorphisme de transition canonique $A_{E_1} \to A_{E_2}$ de $A$-alg\`ebres.
Plus pr\'ecis\'ement, pour une d\'ecomposition $E_2 = E'_2 \amalg E''_2$, posons
$E'_1 = E_1 \cap E'_2$ et $E''_1 = E_1 \cap E''_2$. On a alors
$Z(E'_2, E''_2) \subset Z(E'_1, E''_1)$, d'o\`u un unique homomorphisme de $A$-alg\`ebres
$A_{Z(E'_1, E''_1)}^\sim \to A_{Z(E'_2, E''_2)}^\sim$ d'apr\`es le
lemme \ref{lemma-localization}. En rassemblant ces homomorphismes, on obtient
l'homomorphisme d'anneaux voulu $A_{E_1} \to A_{E_2}$. Notons que le morphisme correspondant
de sch\'emas affines
\begin{equation}
\label{equation-transition}
X_{E_2} \longrightarrow X_{E_1}
\end{equation}
envoie $Z_{E_2}$ dans $Z_{E_1}$. Par unicit\'e, on obtient un syst\`eme de
$A$-alg\`ebres index\'e par $I(A)$, et l'on pose
\begin{equation}
\label{equation-colimit-ring}
A_w = \colim_{E \in I(A)} A_E
\end{equation}
Cette $A$-alg\`ebre est ind-Zariski et fid\`element plate sur $A$.
Enfin, posons $X_w = \Spec(A_w)$ et consid\'erons-y le sous-sch\'ema ferm\'e
$Z = \lim_{E \in I(A)} Z_E$. Cela s'\'ecrit
\begin{equation}
\label{equation-final}
X_w = \lim_{E \in I(A)} X_E \supset Z = \lim_{E \in I(A)} Z_E
\end{equation}

\begin{lemma}
\label{lemma-make-w-local}
Soit $X = \Spec(A)$ un sch\'ema affine. Consid\'erons $A \to A_w$, $X_w = \Spec(A_w)$,
ainsi que $Z \subset X_w$, d\'efinis comme ci-dessus.
\begin{enumerate}
\item $A \to A_w$ est ind-Zariski et fid\`element plat,
\item $X_w \to X$ induit une bijection $Z \to X$,
\item $Z$ est l'ensemble des points ferm\'es de $X_w$,
\item $Z$ est un sch\'ema r\'eduit, et
\item tout point de $X_w$ se sp\'ecialise en un unique point de $Z$.
\end{enumerate}
En particulier, $X_w$ est w-local (d\'efinition \ref{definition-w-local}).
\end{lemma}

\begin{proof}
L'homomorphisme $A \to A_w$ est ind-Zariski par construction.
Pour tout $E$, le morphisme $Z_E \to X$ est bijectif, d'o\`u (2).
Comme $Z \subset X_w$, on en conclut que $X_w \to X$ est surjectif et que
$A \to A_w$ est fid\`element plat d'apr\`es
Alg\`ebre, lemme \ref{algebra-lemma-ff-rings}. Ceci prouve (1).

\medskip\noindent
Supposons que $y \in X_w$ et que $y \not \in Z$. Il existe alors
$E$ tel que l'image de $y$ dans $X_E$ n'appartienne pas \`a
$Z_E$. D\`es que $E \subset E'$, l'image de $y$ dans $X_{E'}$
n'appartient donc pas \`a $Z_{E'}$. Soit $T_{E'} \subset X_{E'}$ le sous-sch\'ema
ferm\'e r\'eduit qui est l'adh\'erence de l'image de $y$. Il est
clair que $T = \lim_{E \subset E'} T_{E'}$ est l'adh\'erence de $y$ dans $X_w$.
D\`es que $E \subset E'$, le sch\'ema $T_{E'} \cap Z_{E'}$ est non vide
par construction de $X_{E'}$. Par suite, $\lim T_{E'} \cap Z_{E'}$ est non vide,
et l'on en conclut que $T \cap Z$ est non vide. Ainsi $y$ n'est pas un point ferm\'e.
Il s'ensuit que tout point ferm\'e de $X_w$ appartient \`a $Z$.

\medskip\noindent
Supposons que $y \in X_w$ se sp\'ecialise en $z, z' \in Z$. Montrons que
$z = z'$, ce qui ach\`evera la preuve de (3) et impliquera (5).
Soient $x, x' \in X$ les images de $z$ et $z'$. Puisque $Z \to X$ est
bijective, il suffit de montrer que $x = x'$. Si $x \not = x'$, alors
il existe $f \in A$ tel que $x \in D(f)$ et $x' \in V(f)$
(ou inversement). Posons $E = \{f\}$; alors
$$
X_E = \Spec(A_f) \amalg \Spec(A_{V(f)}^\sim)
$$
On voit que $z$ et $z'$ ont pour images $x_E$ et $x'_E$, qui appartiennent \`a deux
parties diff\'erentes de la d\'ecomposition ci-dessus de $X_E$. Il est alors impossible
que $x_E$ et $x'_E$ soient des sp\'ecialisations d'un m\^eme point.
C'est la contradiction recherch\'ee.

\medskip\noindent
Rappelons que, pour toute partie finie $E \subset A$, $Z_E$
est une union disjointe des sous-sch\'emas localement ferm\'es $Z(E', E'')$,
chacun isomorphe au spectre de $(A/I)_f$, o\`u $I$ est l'id\'eal
engendr\'e par $E''$ et $f$ le produit des \`el\'ements de $E'$.
Tout \`el\'ement nilpotent $b$ de $(A/I)_f$ est la classe de $g/f^n$
pour un certain $g \in A$. En posant alors $E' = E \cup \{g\}$, le lecteur
v\'erifie que l'image inverse de $b$ est nulle par le morphisme de transition
$Z_{E'} \to Z_E$ du syst\`eme. Ceci prouve (4).
\end{proof}

\begin{remark}
\label{remark-size-w}
Soit $A$ un anneau. Soit $\kappa$ un cardinal infini sup\'erieur ou
\'egal au cardinal de $A$. Alors le cardinal de $A_w$
(lemme \ref{lemma-make-w-local})
est au plus $\kappa$. En effet, chaque $A_E$ est de cardinal au plus
$\kappa$, et l'ensemble des parties finies de $A$ est lui aussi de cardinal au plus $\kappa$.
Le r\'esultat en d\'ecoule puisque $\kappa \otimes \kappa = \kappa$; voir
Ensembles, section \ref{sets-section-cardinals}.
\end{remark}

\begin{lemma}[Propri\'et\'e universelle de la construction]
\label{lemma-universal}
Soit $A$ un anneau. Soit $A \to A_w$ l'homomorphisme d'anneaux construit au
lemme \ref{lemma-make-w-local}. Pour tout homomorphisme d'anneaux $A \to B$ tel que
$\Spec(B)$ soit w-local, il existe une unique factorisation $A \to A_w \to B$
telle que $\Spec(B) \to \Spec(A_w)$ soit w-locale.
\end{lemma}

\begin{proof}
Notons $Y = \Spec(B)$ et $Y_0 \subset Y$ l'ensemble des points ferm\'es.
Notons $f : Y \to X$ le morphisme donn\'e.
Rappelons que $Y_0$ est profini; en particulier, tout sous-ensemble constructible
de $Y_0$ est ouvert et ferm\'e. Soit $E \subset A$ une partie finie.
Rappelons que $A_w = \colim A_E$ et que l'ensemble des points ferm\'es de
$\Spec(A_w)$ est la limite des sous-ensembles ferm\'es $Z_E \subset X_E = \Spec(A_E)$.
Il suffit donc de montrer qu'il existe une unique factorisation $A \to A_E \to B$
telle que $Y \to X_E$ envoie $Y_0$ dans $Z_E$.
Puisque $Z_E \to X = \Spec(A)$ est bijective et que les strates
$Z(E', E'')$ sont constructibles, on voit que
$$
Y_0 = \coprod f^{-1}(Z(E', E'')) \cap Y_0
$$
est une d\'ecomposition en sous-ensembles ouverts et ferm\'es disjoints.
Comme $Y_0 = \pi_0(Y)$, on obtient une d\'ecomposition correspondante de
$Y$ en parties ouvertes et ferm\'ees. Il suffit donc de construire
la factorisation lorsque $f(Y_0) \subset Z(E', E'')$ pour
une d\'ecomposition $E = E' \amalg E''$.
Dans ce cas, $f(Y)$ est contenu dans l'ensemble des points de $X$
qui se sp\'ecialisent en un point de $Z(E', E'')$, ensemble hom\'eomorphe \`a $X_{E', E''}$.
On obtient donc une unique application continue $Y \to X_{E', E''}$ au-dessus de $X$. D'apr\`es le
lemme \ref{lemma-fully-faithful-spaces-over-X},
elle correspond \`a un unique morphisme de sch\'emas
$Y \to X_{E', E''}$ au-dessus de $X$. Ceci ach\`eve la d\'emonstration.
\end{proof}

\noindent
Rappelons que le spectre d'un anneau est profini si et seulement si
tout point est ferm\'e. Il existe en fait toute une s\'erie de conditions
\'equivalentes qui entra\^inent ce fait. Voir
Alg\`ebre, lemme \ref{algebra-lemma-ring-with-only-minimal-primes}, ou
Topologie, lemme \ref{topology-lemma-characterize-profinite-spectral}.

\begin{lemma}
\label{lemma-profinite-goes-up}
Soit $A$ un anneau tel que $\Spec(A)$ soit profini. Soit $A \to B$ un
homomorphisme d'anneaux. Alors $\Spec(B)$ est profini dans chacun des cas suivants:
\begin{enumerate}
\item si $\mathfrak q,\mathfrak q' \subset B$ sont au-dessus d'un m\^eme
id\'eal premier de $A$, on n'a ni $\mathfrak q \subset \mathfrak q'$, ni
$\mathfrak q' \subset \mathfrak q$,
\item $A \to B$ induit des extensions alg\'ebriques des corps r\'esiduels,
\item $A \to B$ est un isomorphisme local,
\item $A \to B$ identifie les anneaux locaux,
\item $A \to B$ est faiblement \'etale,
\item $A \to B$ est quasi-fini,
\item $A \to B$ est non ramifi\'e,
\item $A \to B$ est \'etale,
\item $B$ est une limite inductive filtrante de $A$-alg\`ebres satisfaisant \`a l'une des conditions (1) -- (8),
\item etc.
\end{enumerate}
\end{lemma}

\begin{proof}
D'apr\`es les r\'ef\'erences mentionn\'ees ci-dessus
(Alg\`ebre, lemme \ref{algebra-lemma-ring-with-only-minimal-primes}, ou
Topologie, lemme \ref{topology-lemma-characterize-profinite-spectral}),
il n'existe pas de sp\'ecialisation entre deux points distincts de $\Spec(A)$, et
$\Spec(B)$ est profini si et seulement s'il n'existe pas de sp\'ecialisation
entre deux points distincts de $\Spec(B)$. De telles sp\'ecialisations ne peuvent
avoir lieu que dans les fibres de $\Spec(B) \to \Spec(A)$. On voit ainsi
que (1) est vraie.

\medskip\noindent
L'hypoth\`ese de (2) implique que tous les id\'eaux premiers de $B$ sont maximaux, d'apr\`es
Alg\`ebre, lemme \ref{algebra-lemma-finite-residue-extension-closed}.
Ainsi (2) est vraie.
Si $A \to B$ est un isomorphisme local ou identifie les anneaux locaux,
alors les extensions des corps r\'esiduels sont triviales; (3) et (4)
r\'esultent donc de (2).
Si $A \to B$ est faiblement \'etale, alors Compl\'ements sur l'alg\`ebre, lemme
\ref{more-algebra-lemma-weakly-etale-residue-field-extensions}
montre qu'il induit des extensions alg\'ebriques s\'eparables des corps r\'esiduels, donc
(5) r\'esulte de (2).
Si $A \to B$ est quasi-fini, alors les fibres sont des espaces topologiques
finis discrets. Ainsi (6) r\'esulte de (1).
Ainsi (3) r\'esulte de (1). Les cas (7) et (8)
en d\'ecoulent, car les homomorphismes d'anneaux non ramifi\'es et \'etales sont quasi-finis
(Alg\`ebre, lemmes
\ref{algebra-lemma-unramified-quasi-finite} et
\ref{algebra-lemma-etale-quasi-finite}).
Si $B = \colim B_i$ est une limite inductive filtrante de $A$-alg\`ebres, alors
$\Spec(B) = \lim \Spec(B_i)$ dans la cat\'egorie des espaces topologiques d'apr\`es
Limites, lemme \ref{limits-lemma-inverse-limit-top}.
Ainsi, si chaque $\Spec(B_i)$ est profini, $\Spec(B)$ l'est aussi d'apr\`es
Topologie, lemme \ref{topology-lemma-directed-inverse-limit-profinite}.
Ceci prouve (9).
\end{proof}

\begin{lemma}
\label{lemma-localize-along-closed-profinite}
Soit $A$ un anneau. Soit $V(I) \subset \Spec(A)$ un sous-ensemble ferm\'e
qui est un espace topologique profini. Il existe alors un
homomorphisme d'anneaux ind-Zariski $A \to B$ tel que $\Spec(B)$ soit w-local,
que l'ensemble de ses points ferm\'es soit $V(IB)$ et que $A/I \cong B/IB$.
\end{lemma}

\begin{proof}
Soient $A \to A_w$ et $Z \subset Y = \Spec(A_w)$ comme au
lemme \ref{lemma-make-w-local}.
Soit $T \subset Z$ l'image r\'eciproque de $V(I)$.
Alors $T \to V(I)$ est un hom\'eomorphisme d'apr\`es
Topologie, lemme \ref{topology-lemma-bijective-map}.
Posons $B = (A_w)_T^\sim$; voir le lemme \ref{lemma-localization}.
Il est clair que $B$ est w-local et a pour ensemble de points ferm\'es $V(IB)$.
L'homomorphisme d'anneaux $A/I \to B/IB$ est ind-Zariski
et induit un hom\'eomorphisme entre les espaces topologiques
sous-jacents. C'est donc un isomorphisme d'apr\`es le
lemme \ref{lemma-local-isomorphism-fully-faithful}.
\end{proof}

\begin{lemma}
\label{lemma-w-local-algebraic-residue-field-extensions}
Soit $A$ un anneau tel que $X = \Spec(A)$ soit w-local. Soit $I \subset A$
l'id\'eal radical qui d\'efinit l'ensemble $X_0$ des points ferm\'es de $X$.
Soit $A \to B$ un homomorphisme d'anneaux qui induit des extensions alg\'ebriques
des corps r\'esiduels en chaque id\'eal premier de l'anneau d'arriv\'ee. Alors
\begin{enumerate}
\item tout point de $Z = V(IB)$ est un point ferm\'e de $\Spec(B)$,
\item il existe un homomorphisme d'anneaux ind-Zariski $B \to C$ tel que
\begin{enumerate}
\item $B/IB \to C/IC$ soit un isomorphisme,
\item l'espace $Y = \Spec(C)$ soit w-local,
\item l'application induite $p : Y \to X$ soit w-locale, et
\item $p^{-1}(X_0)$ soit l'ensemble des points ferm\'es de $Y$.
\end{enumerate}
\end{enumerate}
\end{lemma}

\begin{proof}
D'apr\`es le lemme \ref{lemma-profinite-goes-up} appliqu\'e \`a $A/I \to B/IB$,
tous les points de $Z = V(IB) = \Spec(B/IB)$ sont ferm\'es; en fait, $\Spec(B/IB)$
est un espace profini.
Pour achever la d\'emonstration, appliquons le lemme \ref{lemma-localize-along-closed-profinite}
\`a $IB \subset B$.
\end{proof}





\section{Identification des anneaux locaux et condition ind-Zariski}
\label{section-connected-components}

\noindent
Un homomorphisme d'anneaux ind-Zariski $A \to B$ identifie les anneaux locaux
(lemme \ref{lemma-ind-zariski-implies}). La r\'eciproque est fausse
(Exemples, section \ref{examples-section-not-ind-etale}).
Il existe toutefois une sorte de th\'eor\`eme de structure qui d\'ecrit les
homomorphismes d'anneaux identifiant les anneaux locaux au moyen d'homomorphismes
ind-Zariski; voir la proposition \ref{proposition-maps-wich-identify-local-rings}.

\medskip\noindent
Soit $A$ un anneau et posons $X = \Spec(A)$. L'espace des composantes
connexes $\pi_0(X)$ est un espace profini d'apr\`es
Topologie, lemme \ref{topology-lemma-spectral-pi0}
(et Alg\`ebre, lemme \ref{algebra-lemma-spec-spectral}).

\begin{lemma}
\label{lemma-construct}
Soit $A$ un anneau et posons $X = \Spec(A)$. Soit $T \subset \pi_0(X)$ un
sous-ensemble ferm\'e. Il existe un homomorphisme d'anneaux surjectif et ind-Zariski $A \to B$
tel que $\Spec(B) \to \Spec(A)$ induise un hom\'eomorphisme de $\Spec(B)$
sur l'image r\'eciproque de $T$ dans $X$.
\end{lemma}

\begin{proof}
Soit $Z \subset X$ l'image r\'eciproque de $T$. Alors $Z$ est l'intersection
$Z = \bigcap Z_\alpha$ des sous-ensembles ouverts et ferm\'es de $X$ qui contiennent $Z$;
voir Topologie, lemme \ref{topology-lemma-closed-union-connected-components}.
Pour tout $\alpha$, on a $Z_\alpha = \Spec(A_\alpha)$, o\`u
$A \to A_\alpha$ est un isomorphisme local (une localisation en un idempotent).
En posant $B = \colim A_\alpha$, on d\'emontre le lemme.
\end{proof}

\begin{lemma}
\label{lemma-construct-profinite}
Soit $A$ un anneau et posons $X = \Spec(A)$. Soit $T$ un espace profini et
soit $T \to \pi_0(X)$ une application continue. Il existe un
homomorphisme d'anneaux ind-Zariski $A \to B$ tel que, en posant $Y = \Spec(B)$, le diagramme
$$
\xymatrix{
Y \ar[r] \ar[d] & \pi_0(Y) \ar[d] \\
X \ar[r] & \pi_0(X)
}
$$
soit cart\'esien dans la cat\'egorie des espaces topologiques et que
$\pi_0(Y) = T$ comme espaces au-dessus de $\pi_0(X)$.
\end{lemma}

\begin{proof}
\'Ecrivons $T = \lim T_i$ comme la limite d'un syst\`eme projectif d'espaces
finis discrets index\'e par un ensemble ordonn\'e filtrant (voir
Topologie, lemme \ref{topology-lemma-profinite}). Pour tout $i$, posons
$Z_i = \Im(T \to \pi_0(X) \times T_i)$. C'est un sous-ensemble ferm\'e.
Notons que $X \times T_i$ est le spectre de $A_i = \prod_{t \in T_i} A$
et que $A \to A_i$ est un isomorphisme local. D'apr\`es le lemme \ref{lemma-construct},
le sous-ensemble $Z_i \subset \pi_0(X \times T_i) = \pi_0(X) \times T_i$
correspond \`a un homomorphisme d'anneaux surjectif et ind-Zariski $A_i \to B_i$
tel que $\Spec(B_i) = X \times_{\pi_0(X)} Z_i$ comme sous-ensemble de
$X \times T_i$. Les applications de transition $T_i \to T_{i'}$ induisent des applications
$Z_i \to Z_{i'}$ et $X \times_{\pi_0(X)} Z_i \to X \times_{\pi_0(X)} Z_{i'}$,
donc des homomorphismes d'anneaux $B_{i'} \to B_i$
(lemmes \ref{lemma-local-isomorphism-fully-faithful} et
\ref{lemma-ind-zariski-implies}).
Posons $B = \colim B_i$. Puisque $T = \lim Z_i$, on a
$X \times_{\pi_0(X)} T = \lim  X \times_{\pi_0(X)} Z_i$
et donc $Y = \Spec(B) = \lim \Spec(B_i)$
s'ins\`ere dans le diagramme cart\'esien
$$
\xymatrix{
Y \ar[r] \ar[d] & T \ar[d] \\
X \ar[r] & \pi_0(X)
}
$$
d'espaces topologiques. D'apr\`es le lemme \ref{lemma-silly},
on en conclut que $T = \pi_0(Y)$.
\end{proof}

\begin{example}
\label{example-construct-space}
Soit $k$ un corps et soit $T$ un espace topologique profini.
Il existe un homomorphisme d'anneaux ind-Zariski $k \to A$ tel que
$\Spec(A)$ soit hom\'eomorphe \`a $T$. Il suffit en effet d'appliquer le
lemme \ref{lemma-construct-profinite} \`a $T \to \pi_0(\Spec(k)) = \{*\}$.
Dans ce cas, on a m\^eme
$$
A = \colim \text{Map}(T_i, k)
$$
pour toute \`ecriture $T = \lim T_i$ comme limite projective filtrante o\`u chaque $T_i$ est fini.
\end{example}

\begin{lemma}
\label{lemma-w-local-morphism-equal-points-stalks-is-iso}
Soit $A \to B$ un homomorphisme d'anneaux tel que
\begin{enumerate}
\item $A \to B$ identifie les anneaux locaux,
\item les espaces topologiques $\Spec(B)$ et $\Spec(A)$ soient w-locaux,
\item $\Spec(B) \to \Spec(A)$ soit w-locale, et
\item $\pi_0(\Spec(B)) \to \pi_0(\Spec(A))$ soit bijective.
\end{enumerate}
Alors $A \to B$ est un isomorphisme.
\end{lemma}

\begin{proof}
Soient $X_0 \subset X = \Spec(A)$ et $Y_0 \subset Y = \Spec(B)$ les
ensembles de points ferm\'es. Par hypoth\`ese, l'image de $Y_0$ est contenue dans $X_0$ et
l'application induite $Y_0 \to X_0$ est bijective.
Topologiquement, $\Spec(A)$ est l'union disjointe des spectres
des anneaux locaux de $A$ en ses points ferm\'es.
Il en va de m\^eme pour $B$. Ainsi $X \to Y$ est bijective.
Puisque $A \to B$ est plat, le th\'eor\`eme de descente des id\'eaux premiers s'applique
(Alg\`ebre, lemme \ref{algebra-lemma-flat-going-down}).
Ainsi, le lemme \ref{algebra-lemma-unique-prime-over-localize-below} d'Alg\`ebre
montre que, pour tout id\'eal premier $\mathfrak q \subset B$ au-dessus de
$\mathfrak p \subset A$, on a $B_\mathfrak q = B_\mathfrak p$.
Comme $B_\mathfrak q = A_\mathfrak p$ par hypoth\`ese, on
obtient $A_\mathfrak p = B_\mathfrak p$ pour tout id\'eal premier $\mathfrak p$
de $A$. Par cons\'equent, $A = B$ d'apr\`es
Alg\`ebre, lemme \ref{algebra-lemma-characterize-zero-local}.
\end{proof}

\begin{lemma}
\label{lemma-w-local-morphism-equal-stalks-is-ind-zariski}
Soit $A \to B$ un homomorphisme d'anneaux tel que
\begin{enumerate}
\item $A \to B$ identifie les anneaux locaux,
\item les espaces topologiques $\Spec(B)$ et $\Spec(A)$ soient w-locaux, et
\item $\Spec(B) \to \Spec(A)$ soit w-locale.
\end{enumerate}
Alors $A \to B$ est ind-Zariski.
\end{lemma}

\begin{proof}
Posons $X = \Spec(A)$ et $Y = \Spec(B)$. Soient $X_0 \subset X$ et
$Y_0 \subset Y$ les ensembles de points ferm\'es. Soit $A \to A'$ l'homomorphisme d'anneaux ind-Zariski
tel que, en posant $X' = \Spec(A')$, le diagramme
$$
\xymatrix{
X' \ar[r] \ar[d] & \pi_0(X') \ar[d] \\
X \ar[r] & \pi_0(X)
}
$$
soit cart\'esien dans la cat\'egorie des espaces topologiques et que
$\pi_0(X') = \pi_0(Y)$ comme espaces au-dessus de $\pi_0(X)$; voir le
lemme \ref{lemma-construct-profinite}. D'apr\`es le
lemme \ref{lemma-silly}, $X'$ est w-local et
l'ensemble $X'_0 \subset X'$ de ses points ferm\'es est l'image r\'eciproque de $X_0$.

\medskip\noindent
On obtient une application continue $Y \to X'$ entre les espaces topologiques sous-jacents
au-dessus de $X$, laquelle induit l'identification de $\pi_0(Y)$ avec $\pi_0(X')$. D'apr\`es le
lemme \ref{lemma-local-isomorphism-fully-faithful}
(et le lemme \ref{lemma-ind-zariski-implies}),
elle correspond \`a un morphisme de sch\'emas affines $Y \to X'$
au-dessus de $X$. Puisque $Y \to X$ envoie $Y_0$ dans $X_0$, on voit que
$Y \to X'$ envoie $Y_0$ dans $X'_0$, c'est-à-dire que $Y \to X'$ est w-locale.
D'après le lemme \ref{lemma-w-local-morphism-equal-points-stalks-is-iso},
nous voyons que $Y \cong X'$ et le résultat s'ensuit.
\end{proof}

\noindent
La proposition suivante est un résultat préparatoire du type de ceux
que nous démontrerons plus loin.

\begin{proposition}
\label{proposition-maps-wich-identify-local-rings}
Soit $A \to B$ un homomorphisme d'anneaux qui identifie les anneaux locaux.
Alors il existe un homomorphisme d'anneaux fidèlement plat et ind-Zariski
$B \to B'$ tel que $A \to B'$ soit ind-Zariski.
\end{proposition}

\begin{proof}
Soient $A \to A_w$, resp. $B \to B_w$, les homomorphismes d'anneaux fidèlement plats et ind-Zariski
construits dans le lemme \ref{lemma-make-w-local} pour $A$, resp.\ $B$.
Puisque $\Spec(B_w)$ est w-local, il existe une unique factorisation
$A \to A_w \to B_w$ telle que $\Spec(B_w) \to \Spec(A_w)$ soit w-locale
d'après le lemme \ref{lemma-universal}. Remarquons que $A_w \to B_w$ identifie
les anneaux locaux; voir
le lemme \ref{lemma-local-isomorphism-permanence}.
D'après le lemme \ref{lemma-w-local-morphism-equal-stalks-is-ind-zariski},
cela signifie que $A_w \to B_w$ est ind-Zariski. Puisque $B \to B_w$ est
fidèlement plat et ind-Zariski (lemme \ref{lemma-make-w-local})
et que la composée $A \to B \to B_w$ est ind-Zariski
(lemme \ref{lemma-composition-ind-zariski}),
la proposition est démontrée.
\end{proof}

\noindent
La proposition ci-dessus nous permet de caractériser les objets affines
faiblement contractiles du site pro-Zariski d'un schéma affine.

\begin{lemma}
\label{lemma-w-local-extremally-disconnected}
Soit $A$ un anneau. Les assertions suivantes sont équivalentes :
\begin{enumerate}
\item tout homomorphisme d'anneaux fidèlement plat $A \to B$ qui identifie les anneaux locaux
admet une rétraction,
\item tout homomorphisme d'anneaux fidèlement plat et ind-Zariski $A \to B$ admet une rétraction, et
\item $A$ satisfait aux conditions suivantes :
\begin{enumerate}
\item $\Spec(A)$ est w-local, et
\item $\pi_0(\Spec(A))$ est extrêmement discontinu.
\end{enumerate}
\end{enumerate}
\end{lemma}

\begin{proof}
L'équivalence de (1) et (2) découle immédiatement de la
proposition \ref{proposition-maps-wich-identify-local-rings}.

\medskip\noindent
Supposons (3)(a) et (3)(b). Soit $A \to B$ un homomorphisme d'anneaux fidèlement plat et ind-Zariski.
Nous utiliserons, sans plus le mentionner, le fait qu'un homomorphisme plat
$A \to B$ est fidèlement plat si et seulement si tout point fermé
de $\Spec(A)$ appartient à l'image de $\Spec(B) \to \Spec(A)$.
Nous allons montrer que $A \to B$ admet une rétraction.

\medskip\noindent
Soit $I \subset A$ un idéal tel que $V(I) \subset \Spec(A)$ soit
l'ensemble des points fermés de $\Spec(A)$.
Nous pouvons remplacer $B$ par l'anneau $C$ construit dans le
lemme \ref{lemma-w-local-algebraic-residue-field-extensions}
pour $A \to B$ et $I \subset A$.
Ainsi, nous pouvons supposer que $\Spec(B)$ est w-local et que l'ensemble des
points fermés de $\Spec(B)$ est $V(IB)$.

\medskip\noindent
Supposons que $\Spec(B)$ soit w-local et que l'ensemble des points fermés de $\Spec(B)$
soit $V(IB)$. Choisissons une section continue de l'application
continue surjective $V(IB) \to V(I)$. Cela est possible puisque
$V(I) \cong \pi_0(\Spec(A))$ est extrêmement discontinu ; voir la
proposition de Topologie
\ref{topology-proposition-projective-in-category-hausdorff-qc}.
Son image est un sous-espace fermé $T \subset \pi_0(\Spec(B)) \cong V(IB)$
qui est envoyé homéomorphiquement sur $\pi_0(A)$. En remplaçant $B$ par l'anneau quotient ind-Zariski
construit dans le lemme \ref{lemma-construct},
nous voyons que nous pouvons supposer que $\pi_0(\Spec(B)) \to \pi_0(\Spec(A))$
est bijective. À ce stade, $A \to B$ est un isomorphisme d'après le
lemme \ref{lemma-w-local-morphism-equal-points-stalks-is-iso}.

\medskip\noindent
Supposons (1), ou de manière équivalente (2). Soit $A \to A_w$ l'homomorphisme d'anneaux construit dans le
lemme \ref{lemma-make-w-local}. D'après (1), il existe une rétraction $A_w \to A$.
Ainsi, $\Spec(A)$ est homéomorphe à un sous-ensemble fermé de $\Spec(A_w)$. Par le
lemme \ref{lemma-closed-subspace-w-local}, nous voyons que (3)(a) est vérifiée.
Enfin, soit $T \to \pi_0(A)$ une application surjective, où $T$ est un espace topologique
extrêmement discontinu, quasi-compact et séparé
(Topologie, lemme \ref{topology-lemma-existence-projective-cover}).
Choisissons $A \to B$ comme dans le lemme \ref{lemma-construct-profinite},
adapté à $T \to \pi_0(\Spec(A))$. D'après (1), il existe une rétraction
$B \to A$. Ainsi, nous voyons que $T = \pi_0(\Spec(B)) \to \pi_0(\Spec(A))$
admet une section. Un argument formel de théorie des catégories, utilisant la
proposition de Topologie
\ref{topology-proposition-projective-in-category-hausdorff-qc},
implique que $\pi_0(\Spec(A))$ est extrêmement discontinu.
\end{proof}

\begin{lemma}
\label{lemma-find-Zariski-w-contractible}
Soit $A$ un anneau. Il existe un homomorphisme d'anneaux fidèlement plat et ind-Zariski
$A \to B$ tel que $B$ satisfasse aux conditions équivalentes
du lemme \ref{lemma-w-local-extremally-disconnected}.
\end{lemma}

\begin{proof}
Nous appliquons d'abord le lemme \ref{lemma-make-w-local} pour voir que nous pouvons supposer que
$\Spec(A)$ est w-local.
Choisissons un espace extrêmement discontinu $T$ et une application continue surjective
$T \to \pi_0(\Spec(A))$, voir le
lemme de Topologie \ref{topology-lemma-existence-projective-cover}.
Remarquons que $T$ est profini. Appliquons le lemme \ref{lemma-construct-profinite}
pour trouver un homomorphisme d'anneaux ind-Zariski $A \to B$ tel que
$\pi_0(\Spec(B)) \to \pi_0(\Spec(A))$ réalise $T \to \pi_0(\Spec(A))$
et tel que
$$
\xymatrix{
\Spec(B) \ar[r] \ar[d] & \pi_0(\Spec(B)) \ar[d] \\
\Spec(A) \ar[r] & \pi_0(\Spec(A))
}
$$
est cartésien dans la catégorie des espaces topologiques. Remarquons que $\Spec(B)$
est w-local, que $\Spec(B) \to \Spec(A)$ est w-locale et que
l'ensemble des points fermés de $\Spec(B)$ est l'image réciproque de
l'ensemble des points fermés de $\Spec(A)$, voir le lemme \ref{lemma-silly}.
Ainsi, la condition (3) du
lemme \ref{lemma-w-local-extremally-disconnected}
est satisfaite pour $B$.
\end{proof}

\begin{remark}
\label{remark-slightly-stronger}
Dans les lemmes \ref{lemma-construct}, \ref{lemma-construct-profinite},
dans la proposition \ref{proposition-maps-wich-identify-local-rings} et dans le
lemme \ref{lemma-find-Zariski-w-contractible}, nous construisons un homomorphisme d'anneaux ind-Zariski
possédant certaines propriétés. Dans l'article \cite{BS}, les auteurs utilisent la notion
d'ind-(localisation de Zariski), c'est-à-dire de limite inductive filtrante de produits finis
de localisations principales. Dans chacun des résultats énumérés ci-dessus, on peut remplacer ind-Zariski
par ind-(localisation de Zariski).
Nous n'en avons toutefois pas besoin, et la notion d'homomorphisme ind-Zariski
d'anneaux définie ici possède des propriétés formelles légèrement meilleures. De plus,
la notion d'homomorphisme d'anneaux ind-Zariski est l'analogue naturel de la
notion d'homomorphisme d'anneaux ind-\'etale définie à la section suivante.
\end{remark}




\section{Algèbre ind-\'etale}
\label{section-ind-etale}

\noindent
Commençons par une définition.

\begin{definition}
\label{definition-ind-etale}
Un homomorphisme d'anneaux $A \to B$ est dit {\it ind-\'etale} si $B$ peut s'écrire
comme une limite inductive filtrante d'algèbres \'etales sur $A$.
\end{definition}

\noindent
La catégorie des algèbres ind-\'etales est stable par un certain nombre d'opérations
naturelles.

\begin{lemma}
\label{lemma-base-change-ind-etale}
Soient $A \to B$ et $A \to A'$ des homomorphismes d'anneaux. Soit $B' = B \otimes_A A'$
l'algèbre obtenue à partir de $B$ par changement de base.
Si $A \to B$ est ind-\'etale, alors $A' \to B'$ est ind-\'etale.
\end{lemma}

\begin{proof}
C'est le lemme d'Algèbre \ref{algebra-lemma-base-change-colimit-etale}.
\end{proof}

\begin{lemma}
\label{lemma-composition-ind-etale}
Soient $A \to B$ et $B \to C$ des homomorphismes d'anneaux. Si $A \to B$ et $B \to C$
sont ind-\'etales, alors $A \to C$ est ind-\'etale.
\end{lemma}

\begin{proof}
C'est le lemme d'Algèbre \ref{algebra-lemma-composition-colimit-etale}.
\end{proof}

\begin{lemma}
\label{lemma-ind-ind-etale}
Une limite inductive filtrante d'algèbres ind-\'etales sur $A$ est ind-\'etale sur $A$.
\end{lemma}

\begin{proof}
C'est le lemme d'Algèbre \ref{algebra-lemma-colimit-colimit-etale}.
\end{proof}

\begin{lemma}
\label{lemma-ind-etale-permanence}
Soit $A$ un anneau. Soit $B \to C$ un homomorphisme entre des $A$-algèbres ind-\'etales
sur $A$. Alors $C$ est une $B$-algèbre ind-\'etale.
\end{lemma}

\begin{proof}
C'est le lemme d'Algèbre \ref{algebra-lemma-colimits-of-etale}.
\end{proof}

\begin{lemma}
\label{lemma-ind-etale-implies}
Soit $A \to B$ ind-\'etale. Alors $A \to B$ est faiblement \'etale
(Compléments d'Algèbre, définition \ref{more-algebra-definition-weakly-etale}).
\end{lemma}

\begin{proof}
Cela résulte du lemme de Compléments d'Algèbre
\ref{more-algebra-lemma-when-weakly-etale}.
\end{proof}

\begin{lemma}
\label{lemma-lift-ind-etale}
Soient $A$ un anneau et $I \subset A$ un idéal. Le foncteur de changement de base
$$
\text{algèbres ind-\'etales sur }A\text{}
\longrightarrow
\text{algèbres ind-\'etales sur }A/I\text{},\quad
C \longmapsto C/IC
$$
admet un adjoint à droite pleinement fidèle $v$. En particulier, étant donnée
une $A/I$-algèbre ind-\'etale $\overline{C}$, il existe
une $A$-algèbre ind-\'etale $C = v(\overline{C})$ telle que
$\overline{C} = C/IC$.
\end{lemma}

\begin{proof}
Soit $\overline{C}$ une $A/I$-algèbre ind-\'etale.
Considérons la catégorie $\mathcal{C}$ des factorisations
$A \to B \to \overline{C}$ où $A \to B$ est \'etale.
(Nous laissons de côté certaines questions ensemblistes dans cette démonstration.)
Nous allons montrer que cette catégorie est filtrante et que
$C = \colim_\mathcal{C} B$ est une $A$-algèbre ind-\'etale
telle que $\overline{C} = C/IC$.

\medskip\noindent
Montrons d'abord que $\mathcal{C}$ est filtrante
(Catégories, définition \ref{categories-definition-directed}).
La catégorie est non vide, puisque $A \to A \to \overline{C}$ en est un objet.
Supposons que $A \to B \to \overline{C}$ et $A \to B' \to \overline{C}$
soient deux objets de $\mathcal{C}$. Alors $A \to B \otimes_A B' \to \overline{C}$
en est un autre (utiliser le lemme d'Algèbre \ref{algebra-lemma-etale}).
Supposons que $f, g : B \to B'$ soient deux morphismes entre les
objets $A \to B \to \overline{C}$ et $A \to B' \to \overline{C}$
de $\mathcal{C}$. Un coégalisateur est alors
$A \to B' \otimes_{f, B, g} B' \to \overline{C}$.
C'est un objet de $\mathcal{C}$ d'après les
lemmes d'Algèbre \ref{algebra-lemma-etale} et
\ref{algebra-lemma-map-between-etale}.
Ainsi, la catégorie $\mathcal{C}$ est filtrante.

\medskip\noindent
Écrivons $\overline{C} = \colim \overline{B_i}$ comme limite inductive filtrante, avec
$\overline{B_i}$ \'etale sur $A/I$. Pour tout $i$,
il existe un homomorphisme $A \to B_i$ qui est \'etale, tel que $\overline{B_i} = B_i/IB_i$, voir le
lemme d'Algèbre \ref{algebra-lemma-lift-etale}.
Ainsi, $C \to \overline{C}$ est surjectif.
Comme $C/IC \to \overline{C}$ est ind-\'etale
(lemme \ref{lemma-ind-etale-permanence}),
il est plat. Par conséquent, $\overline{C}$ est une localisation de
$C/IC$ par une partie multiplicative $S \subset C/IC$
(Algèbre, lemme \ref{algebra-lemma-pure}).
Soit $f \in C$ dont l'image appartient à $S \subset C/IC$.
Choisissons $A \to B \to \overline{C}$ dans $\mathcal{C}$ et $g \in B$
dont l'image dans la limite inductive est $f$. Alors $A \to B_g \to \overline{C}$
est encore un objet de $\mathcal{C}$. Ainsi, $f$ est un élément inversible
de $C$. Il s'ensuit que $C/IC = \overline{C}$.

\medskip\noindent
Montrons ensuite que, pour toute algèbre ind-\'etale $D$ sur $A$, on a
$$
\Mor_A(D, C) = \Mor_{A/I}(D/ID, \overline{C})
$$
En effet, soit $D/ID \to \overline{C}$ un homomorphisme de $A/I$-algèbres.
Écrivons $D = \colim_{i \in I} D_i$ comme limite inductive indexée par un ensemble préordonné filtrant $I$,
avec $D_i$ \'etale sur $A$. Par le choix de $\mathcal{C}$,
nous obtenons un foncteur $I \to \mathcal{C}$, et donc un homomorphisme
$D \to C$ compatible avec les homomorphismes vers $\overline{C}$. D'où l'assertion.

\medskip\noindent
Il s'ensuit que le foncteur $v$ défini par la règle
$$
\overline{C} 
\longmapsto
v(\overline{C}) = \colim_{A \to B \to \overline{C}} B
$$
est un adjoint à droite du foncteur de changement de base $u$, comme l'affirme le lemme.
Le foncteur $v$ est pleinement fidèle car
$u \circ v = \text{id}$ par construction, voir le
lemme de Catégories \ref{categories-lemma-adjoint-fully-faithful}.
\end{proof}









\section{Construction d'algèbres ind-\'etales}
\label{section-construction-ind-etale}

\noindent
Soit $A$ un anneau. Rappelons que tout homomorphisme d'anneaux \'etale $A \to B$ est isomorphe
à un homomorphisme lisse standard de dimension relative $0$. Un tel homomorphisme
est de la forme
$$
A \longrightarrow A[x_1, \ldots, x_n]/(f_1, \ldots, f_n)
$$
où le déterminant de la matrice $n \times n$ dont les coefficients sont les
$\partial f_i/\partial x_j$ est inversible dans l'anneau quotient. Voir le
lemme d'Algèbre \ref{algebra-lemma-etale-standard-smooth}.

\medskip\noindent
Soit $S(A)$ l'ensemble de toutes les algèbres {\it fidèlement plates}\footnote{En présence
de platitude, par exemple pour des homomorphismes d'anneaux lisses ou \'etales,
cela signifie simplement que l'application induite sur les spectres est surjective. Voir le
lemme d'Algèbre \ref{algebra-lemma-ff-rings}.}
lisses standard sur $A$ et de dimension relative $0$.
Soit $I(A)$ l'ensemble, partiellement ordonné par inclusion, des sous-ensembles finis
$E$ de $S(A)$. Remarquons que $I(A)$ est un ensemble partiellement ordonné
filtrant. Pour $E = \{A \to B_1, \ldots, A \to B_n\}$, posons
$$
B_E = B_1 \otimes_A \ldots \otimes_A B_n
$$
Remarquons que $B_E$ est une $A$-algèbre \'etale fidèlement plate.
Pour $E \subset E'$, il existe un homomorphisme de transition canonique $B_E \to B_{E'}$
de $A$-algèbres \'etales. En effet, supposons $E = \{A \to B_1, \ldots, A \to B_n\}$
et $E' = \{A \to B_1, \ldots, A \to B_{n + m}\}$; alors
$B_E \to B_{E'}$ envoie $b_1 \otimes \ldots \otimes b_n$ sur
l'élément $b_1 \otimes \ldots \otimes b_n \otimes 1 \otimes \ldots \otimes 1$
de $B_{E'}$. Cette construction définit un système d'algèbres fidèlement plates
\'etales sur $A$, indexé par $I(A)$, et nous posons
$$
T(A) = \colim_{E \in I(A)} B_E
$$
Remarquons que $T(A)$ est une $A$-algèbre ind-\'etale fidèlement plate
(Algèbre, lemme \ref{algebra-lemma-colimit-faithfully-flat}). Par construction,
pour toute $A$-algèbre \'etale fidèlement plate $B$, il existe un homomorphisme, non unique,
de $A$-algèbres $B \to T(A)$. En effet, choisissons $(A \to B_0) \in S(A)$
et un isomorphisme $B \cong B_0$. Alors la coprojection canonique
$$
B \to B_0 \to 
T(A) = \colim_{E \in I(A)} B_E
$$
est le morphisme recherché.

\begin{lemma}
\label{lemma-first-construction}
Pour tout anneau $A$, il existe une $A$-algèbre ind-\'etale et fidèlement plate $C$
telle que tout homomorphisme d'anneaux \'etale et fidèlement plat $C \to B$ admette une rétraction.
\end{lemma}

\begin{proof}
Posons $T^1(A) = T(A)$ et $T^{n + 1}(A) = T(T^n(A))$. Soit
$$
C = \colim T^n(A)
$$
Cette algèbre est fidèlement plate sur chaque $T^n(A)$, et en particulier
sur $A$, voir le
lemme d'Algèbre \ref{algebra-lemma-colimit-faithfully-flat}.
De plus, $C$ est ind-\'etale sur $A$ d'après le lemme \ref{lemma-ind-ind-etale}.
Si $C \to B$ est \'etale, il existe alors un entier $n$ et un homomorphisme d'anneaux \'etale
$T^n(A) \to B'$ tels que $B = C \otimes_{T^n(A)} B'$, voir le
lemme d'Algèbre \ref{algebra-lemma-etale}.
Si $C \to B$ est fidèlement plat, alors $\Spec(B) \to \Spec(C) \to \Spec(T^n(A))$
est surjectif; par conséquent, $\Spec(B') \to \Spec(T^n(A))$ est surjectif.
Autrement dit, $T^n(A) \to B'$ est fidèlement plat.
Par construction, il existe un homomorphisme de $T^n(A)$-algèbres
$B' \to T^{n + 1}(A)$. Il induit un homomorphisme de $C$-algèbres $B \to C$,
ce qui achève la démonstration.
\end{proof}

\begin{remark}
\label{remark-size-T}
Soit $A$ un anneau. Soit $\kappa$ un cardinal infini supérieur ou
égal au cardinal de $A$. Alors le cardinal de $T(A)$
est au plus $\kappa$. En effet, chaque $B_E$ est de cardinal au plus
$\kappa$, et l'ensemble d'indices $I(A)$ est lui aussi de cardinal au plus $\kappa$.
Le résultat s'ensuit puisque $\kappa \otimes \kappa = \kappa$, voir la
section d'Ensembles \ref{sets-section-cardinals}. Il s'ensuit également que
l'anneau construit dans la démonstration du lemme \ref{lemma-first-construction}
est de cardinal au plus $\kappa$.
\end{remark}

\begin{remark}
\label{remark-first-construction-functorial}
La construction $A \mapsto T(A)$ est fonctorielle au sens suivant :
si $A \to A'$ est un homomorphisme d'anneaux, nous pouvons construire un diagramme commutatif
$$
\xymatrix{
A \ar[r] \ar[d] & T(A) \ar[d] \\
A' \ar[r] & T(A')
}
$$
En effet, étant donné $(A \to A[x_1, \ldots, x_n]/(f_1, \ldots, f_n))$ dans
$S(A)$, nous pouvons utiliser l'homomorphisme d'anneaux $\varphi : A \to A'$ pour obtenir l'élément correspondant
$(A' \to A'[x_1, \ldots, x_n]/(f^\varphi_1, \ldots, f^\varphi_n))$
de $S(A')$, où $f^\varphi$ désigne le polynôme obtenu en appliquant
$\varphi$ aux coefficients du polynôme $f$.
De plus, on a un diagramme commutatif
$$
\xymatrix{
A \ar[r] \ar[d] & A[x_1, \ldots, x_n]/(f_1, \ldots, f_n) \ar[d] \\
A' \ar[r] & A'[x_1, \ldots, x_n]/(f^\varphi_1, \ldots, f^\varphi_n)
}
$$
dans la catégorie des anneaux. Pour $E \subset S(A)$ fini, posons
$E' = \varphi(E)$ et définissons $B_E \to B_{E'}$ de la manière évidente.
En passant à la limite inductive, on obtient l'homomorphisme recherché $T(A) \to T(A')$, voir le
lemme de Catégories \ref{categories-lemma-functorial-colimit}.
\end{remark}

\begin{lemma}
\label{lemma-have-sections-quotient}
Soit $A$ un anneau tel que tout homomorphisme d'anneaux \'etale et fidèlement plat
$A \to B$ admette une rétraction. Il en va alors de même pour tout anneau quotient
$A/I$.
\end{lemma}

\begin{proof}
Soit $A/I \to \overline{B}$ \'etale et fidèlement plat. D'après le lemme d'Algèbre
\ref{algebra-lemma-lift-etale}, on peut écrire $\overline{B} = B/IB$ pour
un homomorphisme d'anneaux \'etale $A \to B$. L'image $U$ de $\Spec(B) \to \Spec(A)$
est ouverte et contient $V(I)$. Le complémentaire $Z = \Spec(A) \setminus U$
est donc quasi-compact et disjoint de $V(I)$. Par conséquent,
$Z \subset D(f_1) \cup \ldots \cup D(f_r)$ pour un certain $r \geq 0$
et des $f_i \in I$. Alors $A \to B' = B \times \prod A_{f_i}$
est \'etale et fidèlement plat, et $\overline{B} = B'/IB'$.
La rétraction $B' \to A$ de $A \to B'$ induit donc
une rétraction de $A/I \to \overline{B}$.
\end{proof}

\begin{lemma}
\label{lemma-have-sections-strictly-henselian}
Soit $A$ un anneau tel que tout homomorphisme d'anneaux \'etale et fidèlement plat
$A \to B$ admette une rétraction. Alors tout anneau local de $A$ en un idéal
maximal est strictement hensélien.
\end{lemma}

\begin{proof}
Soit $\mathfrak m$ un idéal maximal de $A$. Soit $A \to B$ un homomorphisme d'anneaux
\'etale et soit $\mathfrak q \subset B$ un idéal premier
au-dessus de $\mathfrak m$. D'après la description de l'hensélisé strict
$A_\mathfrak m^{sh}$ donnée dans le
lemme d'Algèbre \ref{algebra-lemma-strict-henselization-different},
il suffit de montrer que $A_\mathfrak m = B_\mathfrak q$.
Remarquons qu'il n'existe qu'un nombre fini d'idéaux premiers
$\mathfrak q = \mathfrak q_1, \mathfrak q_2, \ldots, \mathfrak q_n$
au-dessus de $\mathfrak m$ et qu'il n'y a pas de spécialisations
entre eux, car un homomorphisme d'anneaux \'etale est quasi-fini, voir le
lemme d'Algèbre \ref{algebra-lemma-etale-quasi-finite}.
Ainsi, $\mathfrak q_i$ est un idéal maximal et nous pouvons trouver
$g \in \mathfrak q_2 \cap \ldots \cap \mathfrak q_n$, $g \not \in \mathfrak q$
(Algèbre, lemme \ref{algebra-lemma-silly}).
Après avoir remplacé $B$ par $B_g$, nous voyons que $\mathfrak q$
est le seul idéal premier de $B$ au-dessus de $\mathfrak m$.
L'image $U \subset \Spec(A)$ de $\Spec(B) \to \Spec(A)$ est
ouverte (Algèbre, proposition \ref{algebra-proposition-fppf-open}).
Le complémentaire $\Spec(A) \setminus U$ est donc fermé,
et nous pouvons trouver $f \in A$, $f \not \in \mathfrak p$, tel que
$\Spec(A) = U \cup D(f)$. L'homomorphisme d'anneaux $A \to B \times A_f$
est \'etale et fidèlement plat; il admet donc une rétraction
$\sigma : B \times A_f \to A$ en vertu de l'hypothèse faite sur $A$.
Remarquons que $\sigma$ est \'etale, donc plat comme homomorphisme entre algèbres \'etales
sur $A$ (Algèbre, lemme \ref{algebra-lemma-map-between-etale}).
Puisque $\mathfrak q$ est le seul idéal premier de $B \times A_f$ au-dessus
de $A$, nous trouvons que $A_\mathfrak p \to B_\mathfrak q$ admet
une rétraction qui est elle aussi plate. Ainsi,
$A_\mathfrak p \to B_\mathfrak q \to A_\mathfrak p$
sont des homomorphismes locaux plats d'anneaux locaux dont la composée est l'identité. Comme
un homomorphisme local plat d'anneaux locaux est injectif, nous en concluons que ces
homomorphismes sont les isomorphismes recherchés.
\end{proof}

\begin{lemma}
\label{lemma-have-sections-localize}
Soit $A$ un anneau tel que tout homomorphisme d'anneaux \'etale et fidèlement plat
$A \to B$ admette une rétraction. Soit $Z \subset \Spec(A)$ un sous-schéma fermé.
Soit $A \to A_Z^\sim$ construit comme dans le lemme \ref{lemma-localization}.
Alors tout homomorphisme d'anneaux \'etale et fidèlement plat $A_Z^\sim \to C$ admet
une rétraction.
\end{lemma}

\begin{proof}
Il existe un homomorphisme d'anneaux \'etale $A \to B'$ tel que
$C = B' \otimes_A A_Z^\sim$ comme $A_Z^\sim$-algèbres.
L'image $U' \subset \Spec(A)$ de $\Spec(B') \to \Spec(A)$
est ouverte et contient $V(I)$; il existe donc $f \in I$ tel
que $\Spec(A) = U' \cup D(f)$. Alors $A \to B' \times A_f$
est \'etale et fidèlement plat. Par hypothèse, il existe une rétraction
$B' \times A_f \to A$. En localisant, on obtient la rétraction voulue
$C \to A_Z^\sim$.
\end{proof}

\begin{lemma}
\label{lemma-get-w-local-algebraic-residue-field-extensions}
Soit $A \to B$ un homomorphisme d'anneaux induisant des extensions algébriques des corps résiduels.
Il existe un diagramme commutatif
$$
\xymatrix{
B \ar[r] & D \\
A \ar[r] \ar[u] & C \ar[u]
}
$$
qui possède les propriétés suivantes:
\begin{enumerate}
\item $A \to C$ est fidèlement plat et ind-\'etale,
\item $B \to D$ est fidèlement plat et ind-\'etale,
\item $\Spec(C)$ est w-local,
\item $\Spec(D)$ est w-local,
\item $\Spec(D) \to \Spec(C)$ est w-locale,
\item l'ensemble des points fermés de $\Spec(D)$ est l'image réciproque
de l'ensemble des points fermés de $\Spec(C)$,
\item l'ensemble des points fermés de $\Spec(C)$ a pour image tout $\Spec(A)$,
\item l'ensemble des points fermés de $\Spec(D)$ a pour image tout $\Spec(B)$,
\item pour tout idéal maximal $\mathfrak m \subset C$, l'anneau local
$C_\mathfrak m$ est strictement hensélien.
\end{enumerate}
\end{lemma}

\begin{proof}
Il existe un homomorphisme d'anneaux fidèlement plat et ind-Zariski $A \to A'$ tel que
$\Spec(A')$ soit w-local et que l'ensemble des points fermés de
$\Spec(A')$ a pour image tout $\Spec(A)$; voir le lemme \ref{lemma-make-w-local}.
Soit $I \subset A'$ l'idéal tel que $V(I)$ soit l'ensemble
des points fermés de $\Spec(A')$.
Choisissons $A' \to C'$ comme dans le lemme \ref{lemma-first-construction}.
Notons que les anneaux locaux $C'_{\mathfrak m'}$, pour les idéaux maximaux
$\mathfrak m' \subset C'$, sont strictement henséliens d'après le
lemme \ref{lemma-have-sections-strictly-henselian}.
Appliquons le lemme \ref{lemma-w-local-algebraic-residue-field-extensions}
à $A' \to C'$ et $I \subset A'$ pour obtenir $C' \to C$ avec $C'/IC' \cong C/IC$.
Puisque $A' \to C'$ est fidèlement plat, $\Spec(C'/IC')$
se projette sur l'ensemble des points fermés de $A'$, et donc en particulier
sur $\Spec(A)$. De plus, comme $V(IC) \subset \Spec(C)$
est l'ensemble des points fermés de $C$ et que $C' \to C$ est ind-Zariski
(et identifie les anneaux locaux), on obtient les propriétés (1), (3), (7) et (9).

\medskip\noindent
Notons $J \subset C$ l'idéal tel que $V(J)$ soit l'ensemble des points fermés
de $\Spec(C)$. Posons $D' = B \otimes_A C$. L'homomorphisme d'anneaux
$C \to D'$ induit des extensions algébriques des corps résiduels. Remarquons que,
puisque $V(J) \to \Spec(A)$ est surjective, l'application $T = V(JD') \to \Spec(B)$
est elle aussi surjective. Appliquons le
lemme \ref{lemma-w-local-algebraic-residue-field-extensions}
à $C \to D'$ et $J \subset C$ pour obtenir 
$D' \to D$ avec $D'/JD' \cong D/JD$.
Toutes les autres propriétés de l'énoncé résultent
immédiatement du
lemme \ref{lemma-w-local-algebraic-residue-field-extensions}.
\end{proof}








\section{Faiblement \'etale et pro-\'etale}
\label{section-weakly-etale}

\noindent
Rappelons qu'un homomorphisme d'anneaux $A \to B$ est {\it faiblement \'etale}
si $A \to B$ est plat et si $B \otimes_A B \to B$ est plat. Nous avons
établi certaines propriétés de ces homomorphismes dans
Compléments d'Algèbre, section \ref{more-algebra-section-weakly-etale}.
En particulier, si $A \to B$ est un homomorphisme local et si $A$ est un
anneau local strictement hensélien, alors $A = B$; voir
Compléments d'Algèbre, théorème \ref{more-algebra-theorem-olivier}.
En utilisant ce théorème et le travail accompli ci-dessus, on obtient
le théorème de structure suivant pour les homomorphismes d'anneaux faiblement \'etales.

\begin{proposition}
\label{proposition-weakly-etale}
Soit $A \to B$ un homomorphisme d'anneaux faiblement \'etale.
Il existe alors un homomorphisme d'anneaux fidèlement plat et ind-\'etale
$B \to B'$ tel que $A \to B'$ soit ind-\'etale.
\end{proposition}

\begin{proof}
L'homomorphisme d'anneaux $A \to B$ induit des extensions algébriques (séparables) des
corps résiduels; voir Compléments d'Algèbre, lemme
\ref{more-algebra-lemma-weakly-etale-residue-field-extensions}.
On peut donc appliquer le
lemme \ref{lemma-get-w-local-algebraic-residue-field-extensions}
et choisir un diagramme
$$
\xymatrix{
B \ar[r] & D \\
A \ar[r] \ar[u] & C \ar[u]
}
$$
ayant les propriétés énumérées dans le lemme. Notons que $C \to D$
est faiblement \'etale d'après
Compléments d'Algèbre, lemme \ref{more-algebra-lemma-weakly-etale-permanence}.
Choisissons un idéal maximal $\mathfrak m \subset D$. Par construction,
il est au-dessus d'un idéal maximal $\mathfrak m' \subset C$.
D'après Compléments d'Algèbre, théorème \ref{more-algebra-theorem-olivier},
l'homomorphisme d'anneaux $C_{\mathfrak m'} \to D_\mathfrak m$ est un isomorphisme.
Comme tout point de $\Spec(C)$ se spécialise en un point fermé, on en conclut que
$C \to D$ identifie les anneaux locaux.
La proposition \ref{proposition-maps-wich-identify-local-rings}
s'applique donc à l'homomorphisme d'anneaux $C \to D$. Choisissons $D \to D'$ fidèlement plat
et ind-Zariski tel que $C \to D'$ soit ind-Zariski. Alors
$B \to D'$ résout le problème posé dans la proposition.
\end{proof}





\section{La topologie V et la topologie pro-h}
\label{section-V-versus-pro-h}

\noindent
La topologie V a été introduite dans
Topologies, section \ref{topologies-section-V}.
La topologie h a été introduite dans
Compléments sur la platitude, section \ref{flat-section-h}.
Une sorte de topologie intermédiaire, la topologie ph,
a été introduite dans Topologies, section \ref{topologies-section-ph}.

\medskip\noindent
Étant donnée une topologie $\tau$ sur une catégorie convenable $\mathcal{C}$
de schémas, on peut introduire une ``topologie pro-$\tau$''
sur $\mathcal{C}$ comme suit. Rappelons que, pour tout objet $X$ de $\mathcal{C}$,
on note $h_X$ le préfaisceau représentable associé à $X$.
Disons provisoirement qu'un morphisme $X \to Y$ de $\mathcal{C}$
est $\tau$-couvrant\footnote{Il ne faut pas confondre ceci avec
la notion de recouvrement. Par exemple, si $\tau = \etale$,
tout morphisme $X \to Y$ qui admet une section est $\tau$-couvrant. Mais notre
définition des recouvrements \'etales $\{V_i \to Y\}_{i \in I}$
impose que chaque $V_i \to Y$ soit \'etale.}
si la $\tau$-faisceautisation du morphisme $h_X \to h_Y$
est surjective. On peut alors définir la topologie pro-$\tau$
comme la topologie la moins fine telle que
\begin{enumerate}
\item la topologie pro-$\tau$ soit plus fine que la topologie $\tau$, et
\item $X \to Y$ soit pro-$\tau$-couvrant si
$Y$ est affine et si $X = \lim X_\lambda$ est une limite projective
filtrante de schémas affines $X_\lambda$ au-dessus de $Y$ telle que
$h_{X_\lambda} \to h_Y$ soit $\tau$-couvrant pour tout $\lambda$.
\end{enumerate}
Nous employons cette formulation minutieuse parce que nous ne voulons pas
fixer un choix de recouvrements pro-$\tau$: selon $\tau$,
différents choix de familles de recouvrements conviennent.
Par exemple, nous verrons à la section \ref{section-proetale} que
pour définir la topologie pro-\'etale, il est commode de considérer des familles
de morphismes faiblement \'etales possédant une certaine propriété de finitude.
Plus généralement, la construction proposée dans ce paragraphe
vise surtout à motiver les résultats de cette section, et nous ne définirons jamais
implicitement une topologie pro-$\tau$ par cette méthode.

\medskip\noindent
Le lemme suivant affirme que la
topologie pro-V est égale à la topologie V.

\begin{lemma}
\label{lemma-pro-V-V}
Soit $Y$ un schéma affine. Soit $X = \lim X_i$ une limite projective filtrante
de schémas affines au-dessus de $Y$. Les conditions suivantes sont équivalentes:
\begin{enumerate}
\item $\{X \to Y\}$ est un recouvrement V standard
(Topologies, définition \ref{topologies-definition-standard-V-covering}), et
\item $\{X_i \to Y\}$ est un recouvrement V standard pour tout $i$.
\end{enumerate}
\end{lemma}

\begin{proof}
Le singleton $\{X \to Y\}$ est un recouvrement V standard si et seulement si,
pour tout morphisme $g : \Spec(V) \to Y$, il existe une extension
d'anneaux de valuation $V \subset W$ et un diagramme commutatif
$$
\xymatrix{
\Spec(W) \ar[r] \ar[d] & X \ar[d] \\
\Spec(V) \ar[r]^g & Y
}
$$
L'implication (1) $\Rightarrow$ (2) résulte donc immédiatement de la définition.
Réciproquement, supposons (2) et considérons $g : \Spec(V) \to Y$ comme ci-dessus.
Écrivons $\Spec(V) \times_Y X_i = \Spec(A_i)$.
Puisque $\{X_i \to Y\}$ est un recouvrement V standard, on peut choisir
un anneau de valuation $W_i$ et un homomorphisme d'anneaux $A_i \to W_i$ tels que
la composée $V \to A_i \to W_i$ soit une extension d'anneaux de
valuation. En particulier, le quotient $A'_i$ de $A_i$ par sa
$V$-torsion est une $V$-algèbre fidèlement plate. La platitude résulte de
Compléments d'Algèbre, lemme
\ref{more-algebra-lemma-valuation-ring-torsion-free-flat} et
la surjectivité sur les spectres du fait que $A_i \to W_i$ se factorise par $A'_i$.
Ainsi,
$$
A = \colim A'_i
$$
est une $V$-algèbre fidèlement plate
(Algèbre, lemme \ref{algebra-lemma-colimit-faithfully-flat}).
Puisque $\{\Spec(A) \to \Spec(V)\}$ est un recouvrement fpqc standard, c'est
un recouvrement V standard
(Topologies, lemme \ref{topologies-lemma-standard-fpqc-standard-V});
on peut donc choisir $\Spec(W) \to \Spec(A)$ tel que
$V \to W$ soit une extension d'anneaux de valuation. En composant
avec le morphisme $\Spec(A) \to X = \Spec(\colim A_i)$,
on achève la démonstration.
\end{proof}

\noindent
Le lemme suivant affirme que la topologie pro-h est égale à la
topologie pro-ph, elle-même égale à la topologie V.

\begin{lemma}
\label{lemma-pro-h-V}
Soit $X \to Y$ un morphisme de schémas affines. Les conditions suivantes sont équivalentes:
\begin{enumerate}
\item $\{X \to Y\}$ est un recouvrement V standard
(Topologies, définition \ref{topologies-definition-standard-V-covering}),
\item $X = \lim X_i$ est une limite projective filtrante de schémas affines sur $Y$
telle que $\{X_i \to Y\}$ soit un recouvrement ph pour tout $i$, et
\item $X = \lim X_i$ est une limite projective filtrante de schémas affines sur $Y$
telle que $\{X_i \to Y\}$ soit un recouvrement h pour tout $i$.
\end{enumerate}
\end{lemma}

\begin{proof}
Démonstration de (2) $\Rightarrow$ (1). Rappelons qu'un recouvrement V donné par une
seule flèche entre schémas affines est un recouvrement V standard; voir
Topologies, définition \ref{topologies-definition-V-covering} et
lemme \ref{topologies-lemma-refine-standard-V}.
Rappelons que tout recouvrement ph est un recouvrement V; voir
Topologies, lemme
\ref{topologies-lemma-zariski-etale-smooth-syntomic-fppf-fpqc-ph-V}.
Ainsi, si $X = \lim X_i$ comme en (2), alors $\{X_i \to Y\}$
est un recouvrement V standard pour tout $i$. Le
lemme \ref{lemma-pro-V-V} montre donc que (1) est vraie.

\medskip\noindent
Démonstration de (3) $\Rightarrow$ (2). C'est clair, car un recouvrement h
est toujours un recouvrement ph; voir
Compléments sur la platitude, définition \ref{flat-definition-h-covering}.

\medskip\noindent
Démonstration de (1) $\Rightarrow$ (3). C'est l'implication intéressante, mais
le fond de la démonstration est contenu dans Compléments sur la platitude, lemme
\ref{flat-lemma-equivalence-h-v-locally-finite-presentation}.
Écrivons $X = \Spec(A)$ et $Y = \Spec(R)$. On peut écrire
$A = \colim A_i$, où $A_i$ est de présentation finie sur $R$; voir
Algèbre, lemme \ref{algebra-lemma-ring-colimit-fp}.
Posons $X_i = \Spec(A_i)$. Alors $\{X_i \to Y\}$ est un recouvrement V standard
pour tout $i$, d'après (1) et
Topologies, lemme \ref{topologies-lemma-refine-standard-V}.
Ainsi $\{X_i \to Y\}$ est un recouvrement h d'après
Compléments sur la platitude, définition \ref{flat-definition-h-covering}.
Ceci achève la démonstration.
\end{proof}

\noindent
Le lemme suivant affirme, en gros, qu'un faisceau h
qui commute aux limites satisfait la condition de faisceau pour les recouvrements V.
Comparer aussi avec la remarque \ref{remark-h-limit-preserving}.

\begin{lemma}
\label{lemma-h-limit-preserving}
Soit $S$ un schéma. Soit $F$ un foncteur contravariant défini
sur la catégorie de tous les schémas au-dessus de $S$. Si
\begin{enumerate}
\item $F$ satisfait la condition de faisceau pour la topologie h, et
\item $F$ commute aux limites
(Limites, remarque \ref{limits-remark-limit-preserving}),
\end{enumerate}
alors $F$ satisfait la condition de faisceau pour la topologie V.
\end{lemma}

\begin{proof}
Nous allons le prouver en vérifiant (1) et (2') de
Topologies, lemme \ref{topologies-lemma-sheaf-property-V}.
La condition de faisceau pour les recouvrements de Zariski résulte
du fait que $F$ satisfait la condition de faisceau pour tout recouvrement h.
Enfin, supposons que $X \to Y$ soit un morphisme de schémas affines
au-dessus de $S$ tel que $\{X \to Y\}$ soit un recouvrement V.
D'après le lemme \ref{lemma-pro-h-V}, on peut écrire $X = \lim X_i$
comme une limite projective filtrante de schémas affines
au-dessus de $Y$, telle que $\{X_i \to Y\}$ soit un recouvrement h pour tout $i$.
On obtient
\begin{align*}
&
\text{Égalisateur}(
\xymatrix{
F(X) \ar@<1ex>[r] \ar@<-1ex>[r] &
F(X \times_Y X)
}
)
\\
& =
\text{Égalisateur}(
\xymatrix{
\colim F(X_i) \ar@<1ex>[r] \ar@<-1ex>[r] &
\colim F(X_i \times_Y X_i)
}
)
\\
& =
\colim
\text{Égalisateur}(
\xymatrix{
F(X_i) \ar@<1ex>[r] \ar@<-1ex>[r] &
F(X_i \times_Y X_i)
}
)
\\
& =
\colim F(Y) = F(Y)
\end{align*}
ce qui est l'assertion voulue.
La première égalité vient de ce que $F$ commute aux limites, $X = \lim X_i$ et
$X \times_Y X = \lim X_i \times_Y X_i$.
La deuxième vient de l'exactitude des limites inductives filtrantes.
La troisième vient de ce que $F$ satisfait la condition de faisceau
pour les recouvrements h.
\end{proof}

\begin{remark}
\label{remark-h-limit-preserving}
Soit $S$ un schéma contenu dans un grand site $\Sch_h$. Soit $F$ un faisceau
d'ensembles sur $(\Sch/S)_h$ tel que $F(T) = \colim F(T_i)$ chaque fois que
$T = \lim T_i$ est une limite projective filtrante de schémas affines dans $(\Sch/S)_h$.
Dans cette situation, $F$ se prolonge de manière unique en un foncteur contravariant $F'$
sur la catégorie de tous les schémas au-dessus de $S$, tel que (a) $F'$ satisfasse la
condition de faisceau pour la topologie h et (b) $F'$ commute aux limites.
Voir Compléments sur la platitude, lemme \ref{flat-lemma-extend-sheaf-h}.
Dans cette situation, le lemme \ref{lemma-h-limit-preserving}
affirme que $F'$ satisfait
la condition de faisceau pour la topologie V.
\end{remark}







\section{Construction de recouvrements w-contractiles}
\label{section-w-contractible}

\noindent
Dans cette section, nous construisons des recouvrements w-contractiles de schémas affines.

\begin{definition}
\label{definition-w-contractible}
Soit $A$ un anneau. On dit que $A$ est {\it w-contractile} si tout
homomorphisme d'anneaux fidèlement plat et faiblement \'etale $A \to B$ admet une rétraction.
\end{definition}

\noindent
Remarquons que, d'après la proposition \ref{proposition-weakly-etale},
une définition équivalente consisterait à demander que tout homomorphisme d'anneaux
fidèlement plat et ind-\'etale $A \to B$ admet une rétraction.
Voici une observation essentielle qui nous permettra de construire
des anneaux w-contractiles.

\begin{lemma}
\label{lemma-w-local-strictly-henselian-extremally-disconnected}
Soit $A$ un anneau. Les conditions suivantes sont équivalentes:
\begin{enumerate}
\item $A$ est w-contractile,
\item tout homomorphisme d'anneaux fidèlement plat et ind-\'etale $A \to B$ admet
une rétraction, et
\item $A$ satisfait les conditions suivantes:
\begin{enumerate}
\item $\Spec(A)$ est w-local,
\item $\pi_0(\Spec(A))$ est extrêmement discontinu, et
\item pour tout idéal maximal $\mathfrak m \subset A$, l'anneau
local $A_\mathfrak m$ est strictement hensélien.
\end{enumerate}
\end{enumerate}
\end{lemma}

\begin{proof}
L'équivalence de (1) et (2) résulte immédiatement de la
proposition \ref{proposition-weakly-etale}.

\medskip\noindent
Supposons (3)(a), (3)(b) et (3)(c). Soit $A \to B$ un homomorphisme d'anneaux fidèlement plat
et ind-\'etale. Nous utiliserons sans autre mention le fait qu'un homomorphisme plat
$A \to B$ est fidèlement plat si et seulement si tout point fermé
de $\Spec(A)$ appartient à l'image de $\Spec(B) \to \Spec(A)$.
Montrons que $A \to B$ admet une rétraction.

\medskip\noindent
Soit $I \subset A$ un idéal tel que $V(I) \subset \Spec(A)$ soit
l'ensemble des points fermés de $\Spec(A)$. 
On peut remplacer $B$ par l'anneau $C$ construit au
lemme \ref{lemma-w-local-algebraic-residue-field-extensions}
pour $A \to B$ et $I \subset A$.
On peut ainsi supposer que $\Spec(B)$ est w-local et que l'ensemble des
points fermés de $\Spec(B)$ est égal à $V(IB)$. Dans ce cas, $A \to B$
identifie les anneaux locaux par la condition (3)(c), car il suffit de le vérifier
aux idéaux maximaux de $B$ situés au-dessus d'idéaux maximaux de $A$.
Ainsi $A \to B$ admet une rétraction d'après le
lemme \ref{lemma-w-local-extremally-disconnected}.

\medskip\noindent
Supposons (1) ou, ce qui revient au même, (2). On a (3)(c) d'après le
lemme \ref{lemma-have-sections-strictly-henselian}.
Les propriétés (3)(a) et (3)(b) résultent du
lemme \ref{lemma-w-local-extremally-disconnected}.
\end{proof}

\begin{proposition}
\label{proposition-find-w-contractible}
Pour tout anneau $A$, il existe un homomorphisme d'anneaux fidèlement plat et ind-\'etale
$A \to D$ tel que $D$ soit w-contractile.
\end{proposition}

\begin{proof}
En appliquant le lemme \ref{lemma-get-w-local-algebraic-residue-field-extensions}
à $\text{id}_A : A \to A$, on obtient un homomorphisme d'anneaux fidèlement plat et ind-\'etale
$A \to C$ tel que $C$ soit w-local et que tout anneau local en un
idéal maximal de $C$ soit strictement hensélien.
Choisissons un espace extrêmement discontinu $T$ et une application continue
surjective $T \to \pi_0(\Spec(C))$; voir
Topologie, lemme \ref{topology-lemma-existence-projective-cover}.
Notons que $T$ est profini. Appliquons le lemme \ref{lemma-construct-profinite}
pour obtenir un homomorphisme d'anneaux ind-Zariski $C \to D$ tel que
$\pi_0(\Spec(D)) \to \pi_0(\Spec(C))$ réalise $T \to \pi_0(\Spec(C))$
et que
$$
\xymatrix{
\Spec(D) \ar[r] \ar[d] & \pi_0(\Spec(D)) \ar[d] \\
\Spec(C) \ar[r] & \pi_0(\Spec(C))
}
$$
soit cartésien dans la catégorie des espaces topologiques. Notons que $\Spec(D)$
est w-local, que $\Spec(D) \to \Spec(C)$ est w-locale et que
l'ensemble des points fermés de $\Spec(D)$ est l'image réciproque de
l'ensemble des points fermés de $\Spec(C)$; voir le lemme \ref{lemma-silly}.
Il reste donc vrai que les anneaux locaux de $D$ en ses idéaux
maximaux sont strictement henséliens (puisqu'ils sont isomorphes aux
anneaux locaux de $C$ aux idéaux maximaux correspondants).
Il résulte du
lemme \ref{lemma-w-local-strictly-henselian-extremally-disconnected}
que $D$ est w-contractile.
\end{proof}

\begin{remark}
\label{remark-size-w-contractible}
Soit $A$ un anneau. Soit $\kappa$ un cardinal infini supérieur ou
égal au cardinal de $A$. Alors le cardinal de l'anneau
$D$ construit dans la proposition \ref{proposition-find-w-contractible}
est au plus
$$
\kappa^{2^{2^{2^\kappa}}}.
$$
En effet, l'homomorphisme d'anneaux $A \to D$
s'obtient comme le composé
$$
A \to A_w = A' \to C' \to C \to D.
$$
Les trois premières étapes de la construction sont effectuées
dans le premier paragraphe de la démonstration du
lemme \ref{lemma-get-w-local-algebraic-residue-field-extensions}.
Pour la première étape, on a $|A_w| \leq \kappa$ d'après la
remarque \ref{remark-size-w}.
On a $|C'| \leq \kappa$ d'après la
remarque \ref{remark-size-T}.
Puis $|C| \leq \kappa$, car $C$ est une localisation de $(C')_w$
(il est construit à partir de $C'$ par une application du
lemme \ref{lemma-localize-along-closed-profinite}
dans la démonstration du lemme \ref{lemma-w-local-algebraic-residue-field-extensions}).
Ainsi $C$ possède au plus $2^\kappa$ idéaux maximaux.
Enfin, l'homomorphisme d'anneaux $C \to D$ identifie les anneaux locaux, et le
cardinal de l'ensemble des idéaux maximaux de $D$ est au plus
$2^{2^{2^\kappa}}$ d'après
Topologie, remarque \ref{topology-remark-size-projective-cover}.
Puisque $D \subset \prod_{\mathfrak m \subset D} D_\mathfrak m$, on voit
que le cardinal de $D$ est au plus celui affiché ci-dessus.
\end{remark}

\begin{lemma}
\label{lemma-finite-finitely-presented-over-extremally-disconnected}
Soit $A \to B$ un homomorphisme d'anneaux quasi-fini et de présentation finie.
Si les corps résiduels de $A$ sont séparablement clos
et si $\Spec(A)$ est séparé et extrêmement discontinu, alors $\Spec(B)$ est
extrêmement discontinu.
\end{lemma}

\begin{proof}
Posons $X = \Spec(A)$ et $Y = \Spec(B)$. Choisissons une partition finie
$X = \coprod X_i$ et des morphismes $X'_i \to X_i$ comme dans
Cohomologie \'etale, lemme
\ref{etale-cohomology-lemma-decompose-quasi-finite-morphism}.
L'application d'espaces topologiques $\coprod X_i \to X$ (dont la source
est la somme disjointe dans la catégorie des espaces topologiques)
admet une section d'après la proposition de Topologie
\ref{topology-proposition-projective-in-category-hausdorff-qc}.
Ainsi, $X$ est, comme espace topologique, la somme disjointe des
strates $X_i$. On peut donc remplacer $X$ par les $X_i$ et supposer qu'il
existe un morphisme fini localement libre et surjectif $X' \to X$ tel que
$(X' \times_X Y)_{red}$ soit isomorphe à une somme disjointe finie de
copies de $X'_{red}$. Représentons la situation par le diagramme
$$
\xymatrix{
\coprod_{i = 1, \ldots, r} X' \ar[r] \ar[d] & Y \ar[d] \\
X' \ar[r] & X
}
$$
L'hypothèse sur les corps résiduels de $A$ implique que
ce diagramme est cartésien au niveau des ensembles de points
sous-jacents (détails omis).
Puisque $X$ est extrêmement discontinu et que $X'$ est séparé
(lemme \ref{lemma-profinite-goes-up}), l'application continue
$X' \to X$ admet une section continue $\sigma$. Alors
$\coprod_{i = 1, \ldots, r} \sigma(X) \to Y$ est une application continue
bijective. D'après
Topologie, lemme \ref{topology-lemma-bijective-map},
c'est un homéomorphisme, ce qui achève la démonstration.
\end{proof}

\begin{lemma}
\label{lemma-finite-finitely-presented-over-w-contractible}
Soit $A \to B$ un homomorphisme d'anneaux fini et de présentation finie.
Si $A$ est w-contractile, alors $B$ l'est aussi.
\end{lemma}

\begin{proof}
Nous utiliserons le critère du
lemme \ref{lemma-w-local-strictly-henselian-extremally-disconnected}.
Posons $X = \Spec(A)$ et $Y = \Spec(B)$, et notons $f : Y \to X$
le morphisme induit.
Comme $f : Y \to X$ est un morphisme fini, l'ensemble des points fermés
$Y_0$ de $Y$ est l'image réciproque de l'ensemble des points fermés
$X_0$ de $X$. Soit $y \in Y$ d'image $x \in X$. Alors
$x$ se spécialise en un unique point fermé $x_0 \in X$.
Écrivons $f^{-1}(\{x_0\}) = \{y_1, \ldots, y_n\}$, où les $y_i$
sont fermés dans $Y$. Puisque $R = \mathcal{O}_{X, x_0}$ est strictement hensélien
et que $f$ est fini, $Y \times_{f, X} \Spec(R)$ est
égal à $\coprod_{i = 1, \ldots, n} \Spec(R_i)$, où
chaque $R_i$ est un anneau local fini sur $R$
dont l'idéal maximal correspond à $y_i$; voir
Algèbre, lemme \ref{algebra-lemma-characterize-henselian}, partie (10).
Alors $y$ appartient à exactement l'un de ces $\Spec(R_i)$,
et l'on voit que $y$ se spécialise en exactement l'un des $y_i$.
Autrement dit, tout point de $Y$ se spécialise en un unique
point de $Y_0$. Ainsi, $Y$ est w-local.
Pour tout $y \in Y_0$ d'image $x \in X_0$,
l'anneau $\mathcal{O}_{Y, y}$ est strictement hensélien d'après le
lemme d'Algèbre \ref{algebra-lemma-finite-over-henselian},
appliqué à $\mathcal{O}_{X, x} \to B \otimes_A \mathcal{O}_{X, x}$.
Il reste à montrer que $Y_0$ est extrêmement discontinu.
Pour cela, considérons $X_0 \times_X Y \to X_0$,
où $X_0 \subset X$ est muni de la structure de sous-schéma réduite induite.
Notons que l'espace topologique sous-jacent à
$X_0 \times_X Y$ coïncide avec $Y_0$. Le résultat voulu découle alors du
lemme \ref{lemma-finite-finitely-presented-over-extremally-disconnected}.
\end{proof}

\begin{lemma}
\label{lemma-localization-w-contractible}
Soit $A$ un anneau. Soit $Z \subset \Spec(A)$ un sous-ensemble fermé
de la forme $Z = V(f_1, \ldots, f_r)$. Posons $B = A_Z^\sim$; voir le
lemme \ref{lemma-localization}. Si $A$ est w-contractile, alors $B$ l'est aussi.
\end{lemma}

\begin{proof}
Soit $A_Z^\sim \to B$ un homomorphisme d'anneaux faiblement \'etale et fidèlement plat.
Considérons l'homomorphisme d'anneaux
$$
A \longrightarrow A_{f_1} \times \ldots \times A_{f_r} \times B
$$
qui est fidèlement plat et faiblement \'etale. Si $A$ est w-contractile,
cet homomorphisme admet une rétraction $\sigma$. Considérons le morphisme
$$
\Spec(A_Z^\sim) \to \Spec(A) \xrightarrow{\Spec(\sigma)}
\coprod \Spec(A_{f_i}) \amalg \Spec(B)
$$
Tout point de $Z \subset \Spec(A_Z^\sim)$ a son image dans la composante
$\Spec(B)$. Puisque tout point de $\Spec(A_Z^\sim)$ se spécialise en un
point de $Z$, on obtient un morphisme $\Spec(A_Z^\sim) \to \Spec(B)$
comme on le voulait.
\end{proof}






\section{Le site pro-\'etale}
\label{section-proetale}

\noindent
Dans cette section, nous nous bornons à définir et à construire
les différents sites pro-\'etales ainsi que les morphismes entre eux. L'existence
d'objets faiblement contractiles sera établie dans la
section \ref{section-weakly-contractible}.

\medskip\noindent
La topologie pro-\'etale ressemble quelque peu à
la topologie fpqc (voir Topologies, section \ref{topologies-section-fpqc}),
en ce que le topos des faisceaux sur le petit site pro-\'etale d'un schéma
dépend du choix de la catégorie sous-jacente de schémas. On ne peut donc pas
parler du {\it topos} pro-\'etale d'un schéma. Toutefois, les
groupes de cohomologie d'un faisceau ne changent pas si l'on agrandit
la catégorie sous-jacente de schémas; voir la section \ref{section-change-universe}.

\medskip\noindent
Nous définirons les recouvrements pro-\'etales au moyen de morphismes de schémas faiblement \'etales;
voir Compléments sur les morphismes, section \ref{more-morphisms-section-weakly-etale}.
La raison est que, d'une part, il est quelque peu malaisé de définir
la notion de morphisme pro-\'etale de schémas et que, d'autre part, la
proposition \ref{proposition-weakly-etale}
garantit que l'on obtient les mêmes faisceaux\footnote{Plus précisément,
la topologie pro-\'etale obtenue avec notre choix de recouvrements
est la même que celle issue de la procédure générale expliquée
à la section \ref{section-V-versus-pro-h} en partant de $\tau = \etale$.}
avec la définition suivante.

\begin{definition}
\label{definition-fpqc-covering}
Soit $T$ un schéma. Un {\it recouvrement pro-\'etale de $T$} est une famille
de morphismes de schémas $\{f_i : T_i \to T\}_{i \in I}$
telle que chaque $f_i$ soit faiblement \'etale et que, pour tout ouvert affine
$U \subset T$, il existe $n \geq 0$, une application
$a : \{1, \ldots, n\} \to I$ et des ouverts affines
$V_j \subset T_{a(j)}$, $j = 1, \ldots, n$,
tels que $\bigcup_{j = 1}^n f_{a(j)}(V_j) = U$.
\end{definition}

\noindent
Bien entendu, cette condition implique que $T = \bigcup f_i(T_i)$.
Voici un lemme qui permettra de reconnaître les recouvrements pro-\'etales.
Il permettra aussi de ramener de nombreux lemmes sur les recouvrements pro-\'etales
aux résultats correspondants sur les recouvrements fpqc.

\begin{lemma}
\label{lemma-recognize-proetale-covering}
Soit $T$ un schéma. Soit $\{f_i : T_i \to T\}_{i \in I}$ une famille de
morphismes de schémas de but $T$. Les conditions suivantes sont équivalentes :
\begin{enumerate}
\item $\{f_i : T_i \to T\}_{i \in I}$ est un recouvrement pro-\'etale,
\item chaque $f_i$ est faiblement \'etale et $\{f_i : T_i \to T\}_{i \in I}$
est un recouvrement fpqc,
\item chaque $f_i$ est faiblement \'etale et, pour tout ouvert affine $U \subset T$,
il existe des ouverts quasi-compacts $U_i \subset T_i$, presque tous vides,
tels que $U = \bigcup f_i(U_i)$,
\item chaque $f_i$ est faiblement \'etale et il existe un recouvrement ouvert affine
$T = \bigcup_{\alpha \in A} U_\alpha$ tel que, pour chaque $\alpha \in A$,
il existe $i_{\alpha, 1}, \ldots, i_{\alpha, n(\alpha)} \in I$
et des ouverts quasi-compacts $U_{\alpha, j} \subset T_{i_{\alpha, j}}$ tels que
$U_\alpha =
\bigcup_{j = 1, \ldots, n(\alpha)} f_{i_{\alpha, j}}(U_{\alpha, j})$.
\end{enumerate}
Si $T$ est quasi-séparé, ces conditions sont encore équivalentes à la suivante :
\begin{enumerate}
\item[(5)] chaque $f_i$ est faiblement \'etale et, pour tout $t \in T$, il existe
$i_1, \ldots, i_n \in I$ et des ouverts quasi-compacts $U_j \subset T_{i_j}$
tels que $\bigcup_{j = 1, \ldots, n} f_{i_j}(U_j)$ soit un
voisinage de $t$ dans $T$ (non nécessairement ouvert).
\end{enumerate}
\end{lemma}

\begin{proof}
L'équivalence de (1) et (2) découle immédiatement des définitions.
Le lemme résulte donc de
Topologies, lemme \ref{topologies-lemma-recognize-fpqc-covering}.
\end{proof}

\begin{lemma}
\label{lemma-etale-proetale}
Tout recouvrement \'etale et tout recouvrement de Zariski sont des recouvrements pro-\'etales.
\end{lemma}

\begin{proof}
Cela résulte du résultat correspondant pour les recouvrements fpqc
(Topologies, lemme
\ref{topologies-lemma-zariski-etale-smooth-syntomic-fppf-fpqc}),
du lemme \ref{lemma-recognize-proetale-covering} et
du fait qu'un morphisme \'etale est faiblement \'etale; voir
Compléments sur les morphismes, lemme \ref{more-morphisms-lemma-when-weakly-etale}.
\end{proof}

\begin{lemma}
\label{lemma-proetale}
Soit $T$ un schéma.
\begin{enumerate}
\item Si $T' \to T$ est un isomorphisme, alors $\{T' \to T\}$
est un recouvrement pro-\'etale de $T$.
\item Si $\{T_i \to T\}_{i\in I}$ est un recouvrement pro-\'etale et si, pour chaque
$i$, on a un recouvrement pro-\'etale $\{T_{ij} \to T_i\}_{j\in J_i}$, alors
$\{T_{ij} \to T\}_{i \in I, j\in J_i}$ est un recouvrement pro-\'etale.
\item Si $\{T_i \to T\}_{i\in I}$ est un recouvrement pro-\'etale
et si $T' \to T$ est un morphisme de schémas, alors
$\{T' \times_T T_i \to T'\}_{i\in I}$ est un recouvrement pro-\'etale.
\end{enumerate}
\end{lemma}

\begin{proof}
Cela résulte du fait que les composés et les changements de base
de morphismes faiblement \'etales sont faiblement \'etales
(Compléments sur les morphismes, lemmes
\ref{more-morphisms-lemma-composition-weakly-etale} et
\ref{more-morphisms-lemma-base-change-weakly-etale}),
du lemme \ref{lemma-recognize-proetale-covering} et
des résultats correspondants pour les recouvrements fpqc; voir
Topologies, lemme \ref{topologies-lemma-fpqc}.
\end{proof}

\begin{lemma}
\label{lemma-proetale-affine}
Soit $T$ un schéma affine. Soit $\{T_i \to T\}_{i \in I}$ un recouvrement pro-\'etale
de $T$. Il existe alors un recouvrement pro-\'etale
$\{U_j \to T\}_{j = 1, \ldots, n}$ qui raffine
$\{T_i \to T\}_{i \in I}$ et dans lequel chaque $U_j$ est un schéma
affine. De plus, on peut choisir chaque $U_j$ comme ouvert affine
de l'un des $T_i$.
\end{lemma}

\begin{proof}
Cela résulte directement de la définition.
\end{proof}

\noindent
Nous définissons donc comme suit les recouvrements standards correspondants dans le cas affine.

\begin{definition}
\label{definition-standard-proetale}
Soit $T$ un schéma affine. Un {\it recouvrement pro-\'etale standard}
de $T$ est une famille $\{f_i : T_i \to T\}_{i = 1, \ldots, n}$
où chaque $T_j$ est affine, chaque $f_i$ est faiblement \'etale et
$T = \bigcup f_i(T_i)$.
\end{definition}

\noindent
Nous suivons le schéma général donné dans
Topologies, section \ref{topologies-section-procedure},
pour construire le grand site pro-\'etale avec lequel nous travaillerons.
Toutefois, comme nous avons besoin d'anneaux un peu plus grands pour tenir compte de la taille
de certaines constructions, nous modifions légèrement celles-ci.

\begin{definition}
\label{definition-big-proetale-site}
Un {\it grand site pro-\'etale} est un site $\Sch_\proetale$ quelconque
au sens de Sites, définition \ref{sites-definition-site}, construit comme suit :
\begin{enumerate}
\item Choisir un ensemble quelconque de schémas $S_0$ et un ensemble quelconque de recouvrements pro-\'etales
$\text{Cov}_0$ parmi ces schémas.
\item Remplacer la fonction $Bound$ de
Ensembles, équation (\ref{sets-equation-bound}), par
$$
Bound(\kappa) = \max\{\kappa^{2^{2^{2^\kappa}}}, \kappa^{\aleph_0}, \kappa^+\}.
$$
\item Prendre pour catégorie sous-jacente une catégorie $\Sch_\alpha$ quelconque
construite comme dans Ensembles, lemme \ref{sets-lemma-construct-category},
à partir de l'ensemble $S_0$ et de la fonction $Bound$.
\item Choisir un ensemble quelconque de recouvrements comme dans
Ensembles, lemme \ref{sets-lemma-coverings-site}, à partir de la
catégorie $\Sch_\alpha$, de la classe des recouvrements pro-\'etales
et de l'ensemble $\text{Cov}_0$ choisi ci-dessus.
\end{enumerate}
\end{definition}

\noindent
Voir les remarques qui suivent Topologies,
définition \ref{topologies-definition-big-zariski-site},
pour la motivation et les explications relatives à la définition des grands sites.

\medskip\noindent
Il résultera du lemme \ref{lemma-proetale-induced} que la
topologie d'un grand site pro-\'etale $\Sch_\proetale$ est, en un certain sens,
induite par la topologie pro-\'etale sur la catégorie de tous les schémas.

\begin{definition}
\label{definition-big-small-proetale}
Soit $S$ un schéma. Soit $\Sch_\proetale$ un grand site pro-\'etale
contenant $S$.
\begin{enumerate}
\item Le {\it grand site pro-\'etale de $S$}, noté
$(\Sch/S)_\proetale$, est le site $\Sch_\proetale/S$
introduit dans Sites, section \ref{sites-section-localize}.
\item Le {\it petit site pro-\'etale de $S$}, que nous notons
$S_\proetale$, est la sous-catégorie pleine de $(\Sch/S)_\proetale$
dont les objets sont les $U/S$ tels que $U \to S$ soit faiblement \'etale.
Un recouvrement de $S_\proetale$ est tout recouvrement $\{U_i \to U\}$ de
$(\Sch/S)_\proetale$ avec $U \in \Ob(S_\proetale)$.
\item Le {\it grand site pro-\'etale affine de $S$}, noté
$(\textit{Aff}/S)_\proetale$, est la sous-catégorie pleine de
$(\Sch/S)_\proetale$ dont les objets sont les $U/S$ affines.
Un recouvrement de $(\textit{Aff}/S)_\proetale$ est tout recouvrement
$\{U_i \to U\}$ de $(\Sch/S)_\proetale$ qui est un
recouvrement pro-\'etale standard.
\end{enumerate}
\end{definition}

\noindent
Il n'est pas tout à fait évident que le petit site pro-\'etale et
le grand site pro-\'etale affine soient des sites. Vérifions-le maintenant.

\begin{lemma}
\label{lemma-verify-site-proetale}
Soit $S$ un schéma. Soit $\Sch_\proetale$ un grand site pro-\'etale
contenant $S$. Alors $S_\proetale$ et $(\textit{Aff}/S)_\proetale$ sont tous deux des sites.
\end{lemma}

\begin{proof}
Montrons que $S_\proetale$ est un site. C'est une catégorie munie d'un
ensemble donné de familles de morphismes à cible fixée. Il faut donc
vérifier les propriétés (1), (2) et (3) de
Sites, définition \ref{sites-definition-site}.
Puisque $(\Sch/S)_\proetale$ est un site, il suffit de montrer
que, pour tout recouvrement $\{U_i \to U\}$ de $(\Sch/S)_\proetale$
avec $U \in \Ob(S_\proetale)$, on a aussi $U_i \in \Ob(S_\proetale)$.
Cela découle des définitions,
car la composée de morphismes faiblement \'etales est faiblement \'etale.

\medskip\noindent
Pour montrer que $(\textit{Aff}/S)_\proetale$ est un site, en raisonnant comme ci-dessus,
il suffit de montrer que la famille des recouvrements pro-\'etales standards
de schémas affines satisfait les propriétés (1), (2) et (3) de
Sites, définition \ref{sites-definition-site}.
Cela découle du lemme \ref{lemma-recognize-proetale-covering}
et du résultat correspondant pour les recouvrements fpqc standards
(Topologies, lemme \ref{topologies-lemma-fpqc-affine-axioms}).
\end{proof}

\begin{lemma}
\label{lemma-fibre-products-proetale}
Soit $S$ un schéma. Soit $\Sch_\proetale$ un grand site pro-\'etale
contenant $S$. Soit $\Sch$ la catégorie de tous les schémas.
\begin{enumerate}
\item Les catégories $\Sch_\proetale$, $(\Sch/S)_\proetale$,
$S_\proetale$ et $(\textit{Aff}/S)_\proetale$ possèdent des produits fibrés
qui coïncident avec les produits fibrés dans $\Sch$.
\item Les catégories $\Sch_\proetale$, $(\Sch/S)_\proetale$ et
$S_\proetale$ possèdent des égalisateurs qui coïncident avec les égalisateurs dans $\Sch$.
\item Les catégories $(\Sch/S)_\proetale$ et $S_\proetale$ possèdent toutes deux
un objet final, à savoir $S/S$.
\item La catégorie $\Sch_\proetale$ possède un objet final qui coïncide
avec l'objet final de $\Sch$, à savoir $\Spec(\mathbf{Z})$.
\end{enumerate}
\end{lemma}

\begin{proof}
La catégorie $\Sch_\proetale$ contient $\Spec(\mathbf{Z})$ et
est stable par produits et produits fibrés par construction; voir
Ensembles, lemme \ref{sets-lemma-what-is-in-it}.
Supposons donnés des morphismes $U \to S$, $V \to U$, $W \to U$
de schémas avec $U, V, W \in \Ob(\Sch_\proetale)$.
Le produit fibré $V \times_U W$ dans $\Sch_\proetale$
est un produit fibré dans $\Sch$ et
est le produit fibré de $V/S$ et $W/S$ au-dessus de $U/S$ dans
la catégorie de tous les schémas au-dessus de $S$; c'est donc aussi un
produit fibré dans $(\Sch/S)_\proetale$.
Cela démontre le résultat pour $(\Sch/S)_\proetale$.
Si $U \to S$, $V \to U$ et $W \to U$ sont faiblement \'etales, alors
$V \times_U W \to S$ l'est aussi (voir
Compléments sur les morphismes, section \ref{more-morphisms-section-weakly-etale});
on obtient donc les produits fibrés dans $S_\proetale$.
Si $U, V, W$ sont affines, alors $V \times_U W$ l'est aussi; on obtient donc
les produits fibrés dans $(\textit{Aff}/S)_\proetale$.

\medskip\noindent
Soient $a, b : U \to V$ deux morphismes de $\Sch_\proetale$.
Dans ce cas, l'égalisateur de $a$ et $b$ (dans la catégorie des schémas) est
$$
V
\times_{\Delta_{V/\Spec(\mathbf{Z})}, V \times_{\Spec(\mathbf{Z})} V, (a, b)}
(U \times_{\Spec(\mathbf{Z})} U)
$$
qui est un objet de $\Sch_\proetale$ d'après ce que nous venons de voir.
Ainsi, $\Sch_\proetale$ possède des égalisateurs. Si $a$ et $b$ sont des morphismes au-dessus de $S$,
alors l'égalisateur (dans la catégorie des schémas) est aussi donné par
$$
V \times_{\Delta_{V/S}, V \times_S V, (a, b)} (U \times_S U)
$$
et l'on voit donc que $(\Sch/S)_\proetale$ possède des égalisateurs. De plus, si
$U$ et $V$ sont faiblement \'etales sur $S$, alors l'égalisateur ci-dessus
l'est aussi, puisqu'il est un produit fibré de schémas faiblement \'etales sur $S$.
Ainsi, $S_\proetale$ possède des égalisateurs. Les assertions sur les objets finaux
sont claires.
\end{proof}

\noindent
Vérifions ensuite que le grand site pro-\'etale affine définit le même
topos que le grand site pro-\'etale.

\begin{lemma}
\label{lemma-affine-big-site-proetale}
Soit $S$ un schéma. Soit $\Sch_\proetale$ un grand site pro-\'etale
contenant $S$.
Le foncteur $(\textit{Aff}/S)_\proetale \to (\Sch/S)_\proetale$
est un foncteur spécial cocontinu. Il induit donc une équivalence
de topos de $\Sh((\textit{Aff}/S)_\proetale)$ vers
$\Sh((\Sch/S)_\proetale)$.
\end{lemma}

\begin{proof}
La notion de foncteur spécial cocontinu est introduite dans
Sites, définition \ref{sites-definition-special-cocontinuous-functor}.
Il faut donc vérifier les hypothèses (1) -- (5) de
Sites, lemme \ref{sites-lemma-equivalence}.
Notons le foncteur d'inclusion
$u : (\textit{Aff}/S)_\proetale \to (\Sch/S)_\proetale$.
La cocontinuité signifie simplement que tout recouvrement pro-\'etale de
$T/S$, avec $T$ affine, peut être raffiné par un recouvrement pro-\'etale
standard de $T$. C'est le contenu du
lemme \ref{lemma-proetale-affine}.
Ainsi, (1) est vérifiée. Le foncteur $u$ est continu simplement parce qu'un recouvrement
pro-\'etale standard est un recouvrement pro-\'etale. Ainsi, (2) est vérifiée.
Les parties (3) et (4) découlent immédiatement du fait que $u$ est
pleinement fidèle. Enfin, la condition (5) résulte du
fait que tout schéma admet un recouvrement ouvert affine.
\end{proof}

\begin{lemma}
\label{lemma-put-in-T}
Soit $\Sch_\proetale$ un grand site pro-\'etale.
Soit $f : T \to S$ un morphisme de $\Sch_\proetale$.
Le foncteur $T_\proetale \to (\Sch/S)_\proetale$
est cocontinu et induit un morphisme de topos
$$
i_f :
\Sh(T_\proetale)
\longrightarrow
\Sh((\Sch/S)_\proetale)
$$
Pour tout faisceau $\mathcal{G}$ sur $(\Sch/S)_\proetale$,
on a la formule $(i_f^{-1}\mathcal{G})(U/T) = \mathcal{G}(U/S)$.
Le foncteur $i_f^{-1}$ admet aussi un adjoint à gauche $i_{f, !}$ qui commute
aux produits fibrés et aux égalisateurs.
\end{lemma}

\begin{proof}
Notons $u : T_\proetale \to (\Sch/S)_\proetale$ ce foncteur.
Autrement dit, étant donné un morphisme faiblement \'etale $j : U \to T$ correspondant
à un objet de $T_\proetale$, on pose $u(U \to T) = (f \circ j : U \to S)$.
Ce foncteur commute aux produits fibrés ; voir le
lemme \ref{lemma-fibre-products-proetale}.
De plus, $T_\proetale$ possède des égalisateurs et $u$ commute avec eux
d'après le lemme \ref{lemma-fibre-products-proetale}.
Il est manifestement cocontinu.
Il est aussi continu, car $u$ transforme les recouvrements en recouvrements et
commute aux produits fibrés. Le lemme résulte donc de
Sites, lemmes \ref{sites-lemma-when-shriek}
et \ref{sites-lemma-preserve-equalizers}.
\end{proof}

\begin{lemma}
\label{lemma-at-the-bottom}
Soit $S$ un schéma. Soit $\Sch_\proetale$ un grand site pro-\'etale
contenant $S$.
Le foncteur d'inclusion $S_\proetale \to (\Sch/S)_\proetale$
satisfait aux hypothèses de Sites, lemme \ref{sites-lemma-bigger-site},
et induit donc un morphisme de sites
$$
\pi_S : (\Sch/S)_\proetale \longrightarrow S_\proetale
$$
et un morphisme de topos
$$
i_S : \Sh(S_\proetale) \longrightarrow \Sh((\Sch/S)_\proetale)
$$
tels que $\pi_S \circ i_S = \text{id}$. De plus, $i_S = i_{\text{id}_S}$,
où $i_{\text{id}_S}$ est défini au lemme \ref{lemma-put-in-T}. En particulier, le
foncteur $i_S^{-1} = \pi_{S, *}$ est décrit par la règle
$i_S^{-1}(\mathcal{G})(U/S) = \mathcal{G}(U/S)$.
\end{lemma}

\begin{proof}
Dans ce cas, le foncteur $u : S_\proetale \to (\Sch/S)_\proetale$,
outre les propriétés établies dans la démonstration du
lemme \ref{lemma-put-in-T} ci-dessus, est pleinement fidèle
et envoie l'objet final sur l'objet final.
Le lemme découle de Sites, lemme \ref{sites-lemma-bigger-site}.
\end{proof}

\begin{definition}
\label{definition-restriction-small-proetale}
Dans la situation du
lemme \ref{lemma-at-the-bottom},
le foncteur $i_S^{-1} = \pi_{S, *}$ est souvent
appelé la {\it restriction au petit site pro-\'etale} ; pour un faisceau
$\mathcal{F}$ sur le grand site pro-\'etale, on note
$\mathcal{F}|_{S_\proetale}$ cette restriction.
\end{definition}

\noindent
Avec cette notation, on a, pour un faisceau $\mathcal{F}$ sur le
grand site et un faisceau $\mathcal{G}$ sur le petit site,
\begin{align*}
\Mor_{\Sh(S_\proetale)}(\mathcal{F}|_{S_\proetale}, \mathcal{G})
& =
\Mor_{\Sh((\Sch/S)_\proetale)}(\mathcal{F},
i_{S, *}\mathcal{G}) \\
\Mor_{\Sh(S_\proetale)}(\mathcal{G}, \mathcal{F}|_{S_\proetale})
& =
\Mor_{\Sh((\Sch/S)_\proetale)}(\pi_S^{-1}\mathcal{G}, \mathcal{F})
\end{align*}
De plus, $(i_{S, *}\mathcal{G})|_{S_\proetale} = \mathcal{G}$
et $(\pi_S^{-1}\mathcal{G})|_{S_\proetale} = \mathcal{G}$.

\begin{lemma}
\label{lemma-morphism-big}
Soit $\Sch_\proetale$ un grand site pro-\'etale.
Soit $f : T \to S$ un morphisme de $\Sch_\proetale$.
Le foncteur
$$
u : (\Sch/T)_\proetale \longrightarrow (\Sch/S)_\proetale, \quad
V/T \longmapsto V/S
$$
est cocontinu et admet un adjoint à droite continu
$$
v : (\Sch/S)_\proetale \longrightarrow (\Sch/T)_\proetale, \quad
(U \to S) \longmapsto (U \times_S T \to T).
$$
Ils induisent le même morphisme de topos
$$
f_{big} :
\Sh((\Sch/T)_\proetale)
\longrightarrow
\Sh((\Sch/S)_\proetale)
$$
On a $f_{big}^{-1}(\mathcal{G})(U/T) = \mathcal{G}(U/S)$.
On a $f_{big, *}(\mathcal{F})(U/S) = \mathcal{F}(U \times_S T/T)$.
De plus, $f_{big}^{-1}$ admet un adjoint à gauche $f_{big!}$ qui commute aux
produits fibrés et aux égalisateurs.
\end{lemma}

\begin{proof}
Le foncteur $u$ est cocontinu, continu et commute aux produits fibrés
et aux égalisateurs (détails omis ; comparer avec la démonstration du
lemme \ref{lemma-put-in-T}). Par conséquent,
Sites, lemmes \ref{sites-lemma-when-shriek} et
\ref{sites-lemma-preserve-equalizers}
s'appliquent, et l'on en déduit la formule
pour $f_{big}^{-1}$ ainsi que l'existence de $f_{big!}$. De plus,
le foncteur $v$ est adjoint à droite, car, pour $U/T$ et $V/S$,
on a $\Mor_S(u(U), V) = \Mor_T(U, V \times_S T)$,
ce qui est l'assertion voulue. On peut donc appliquer
Sites, lemmes \ref{sites-lemma-have-functor-other-way} et
\ref{sites-lemma-have-functor-other-way-morphism}
pour obtenir la formule pour $f_{big, *}$.
\end{proof}

\begin{lemma}
\label{lemma-morphism-big-small}
Soit $\Sch_\proetale$ un grand site pro-\'etale.
Soit $f : T \to S$ un morphisme de $\Sch_\proetale$.
\begin{enumerate}
\item On a $i_f = f_{big} \circ i_T$, où $i_f$ est défini au
lemme \ref{lemma-put-in-T} et $i_T$ au
lemme \ref{lemma-at-the-bottom}.
\item Le foncteur $S_\proetale \to T_\proetale$,
$(U \to S) \mapsto (U \times_S T \to T)$ est continu et induit
un morphisme de topos
$$
f_{small} : \Sh(T_\proetale) \longrightarrow \Sh(S_\proetale).
$$
On a $f_{small, *}(\mathcal{F})(U/S) = \mathcal{F}(U \times_S T/T)$.
\item On a un diagramme commutatif de morphismes de sites
$$
\xymatrix{
T_\proetale \ar[d]_{f_{small}} &
(\Sch/T)_\proetale \ar[d]^{f_{big}} \ar[l]^{\pi_T}\\
S_\proetale &
(\Sch/S)_\proetale \ar[l]_{\pi_S}
}
$$
de sorte que $f_{small} \circ \pi_T = \pi_S \circ f_{big}$ comme morphismes de topos.
\item On a $f_{small} = \pi_S \circ f_{big} \circ i_T = \pi_S \circ i_f$.
\end{enumerate}
\end{lemma}

\begin{proof}
L'égalité $i_f = f_{big} \circ i_T$ résulte de
l'égalité $i_f^{-1} = i_T^{-1} \circ f_{big}^{-1}$, qui ressort
clairement des descriptions de ces foncteurs ci-dessus.
Cela démontre (1).

\medskip\noindent
Le foncteur $u : S_\proetale \to T_\proetale$,
$u(U \to S) = (U \times_S T \to T)$,
transforme les recouvrements en recouvrements et commute aux produits fibrés ;
voir les lemmes \ref{lemma-proetale} et \ref{lemma-fibre-products-proetale}.
De plus, $S_\proetale$ et $T_\proetale$ ont tous deux un objet final,
à savoir $S/S$ et $T/T$, et $u(S/S) = T/T$. D'après
Sites, proposition \ref{sites-proposition-get-morphism},
le foncteur $u$ correspond donc à un morphisme de sites
$T_\proetale \to S_\proetale$. Celui-ci donne à son tour le
morphisme de topos ; voir
Sites, lemme \ref{sites-lemma-morphism-sites-topoi}. La description
de l'image directe résulte clairement de ces références.

\medskip\noindent
La partie (3) résulte du fait que $\pi_S$ et $\pi_T$ sont donnés par les
foncteurs d'inclusion, et $f_{small}$ et $f_{big}$ par les
foncteurs de changement de base $U \mapsto U \times_S T$.

\medskip\noindent
L'assertion (4) résulte de (3) en précomposant par $i_T$.
\end{proof}

\noindent
Dans la situation du lemme, avec la terminologie de la
définition \ref{definition-restriction-small-proetale},
on a, pour $\mathcal{F}$ un faisceau sur le grand site pro-\'etale de $T$,
\begin{equation}
\label{equation-compare-big-small}
(f_{big, *}\mathcal{F})|_{S_\proetale} =
f_{small, *}(\mathcal{F}|_{T_\proetale}),
\end{equation}
Cette égalité résulte clairement de la commutativité du diagramme de
sites du lemme, puisque la restriction au petit site pro-\'etale de
$T$, resp.\ $S$, est donnée par $\pi_{T, *}$, resp.\ $\pi_{S, *}$. Une formule
analogue faisant intervenir les images inverses et les restrictions est fausse.

\begin{lemma}
\label{lemma-composition-proetale}
Étant donnés des schémas $X$, $Y$, $Y$ dans $\Sch_\proetale$
et des morphismes $f : X \to Y$, $g : Y \to Z$, on a
$g_{big} \circ f_{big} = (g \circ f)_{big}$ et
$g_{small} \circ f_{small} = (g \circ f)_{small}$.
\end{lemma}

\begin{proof}
Cela découle de la description explicite des foncteurs image directe
et image inverse sur les grands sites donnée dans le
lemme \ref{lemma-morphism-big}. Pour les foncteurs
sur les petits sites, cela découle de la description des
foncteurs image directe du lemme \ref{lemma-morphism-big-small}.
\end{proof}

\begin{lemma}
\label{lemma-morphism-big-small-cartesian-diagram}
Soit $\Sch_\proetale$ un grand site pro-\'etale. Considérons un diagramme cartésien
$$
\xymatrix{
T' \ar[r]_{g'} \ar[d]_{f'} & T \ar[d]^f \\
S' \ar[r]^g & S
}
$$
dans $\Sch_\proetale$. Alors
$i_g^{-1} \circ f_{big, *} = f'_{small, *} \circ (i_{g'})^{-1}$
et $g_{big}^{-1} \circ f_{big, *} = f'_{big, *} \circ (g'_{big})^{-1}$.
\end{lemma}

\begin{proof}
Puisque le diagramme est cartésien, on a, pour $U'/S'$,
$U' \times_{S'} T' = U' \times_S T$. Ainsi,
$i_g^{-1} \circ f_{big, *}$ et $f'_{small, *} \circ (i_{g'})^{-1}$
associent tous deux à un faisceau $\mathcal{F}$ sur $(\Sch/T)_\proetale$ le faisceau
$U' \mapsto \mathcal{F}(U' \times_{S'} T')$ sur $S'_\proetale$
(d'après les lemmes \ref{lemma-put-in-T} et \ref{lemma-morphism-big}).
La seconde égalité se démontre de la même manière ou peut se
déduire du résultat très général suivant :
Sites, lemme \ref{sites-lemma-localize-morphism}.
\end{proof}

\noindent
On peut considérer un faisceau sur le grand site pro-\'etale de $S$ comme une famille
de faisceaux sur les petits sites pro-\'etales des schémas au-dessus de $S$.

\begin{lemma}
\label{lemma-characterize-sheaf-big}
Soit $S$ un schéma contenu dans un grand site pro-\'etale $\Sch_\proetale$.
Un faisceau $\mathcal{F}$ sur le grand site pro-\'etale $(\Sch/S)_\proetale$
est déterminé par les données suivantes :
\begin{enumerate}
\item pour tout $T/S \in \Ob((\Sch/S)_\proetale)$, un faisceau
$\mathcal{F}_T$ sur $T_\proetale$,
\item pour tout $f : T' \to T$ dans
$(\Sch/S)_\proetale$, un morphisme
$c_f : f_{small}^{-1}\mathcal{F}_T \to \mathcal{F}_{T'}$.
\end{enumerate}
Ces données sont soumises aux conditions suivantes :
\begin{enumerate}
\item[(a)] pour tous $f : T' \to T$ et $g : T'' \to T'$ dans
$(\Sch/S)_\proetale$, le composé
$c_g \circ g_{small}^{-1}c_f$ est égal à $c_{f \circ g}$, et
\item[(b)] si $f : T' \to T$ dans $(\Sch/S)_\proetale$
est faiblement \'etale, alors $c_f$ est un isomorphisme.
\end{enumerate}
\end{lemma}

\begin{proof}
La démonstration est identique à celle du
Topologies, lemme \ref{topologies-lemma-characterize-sheaf-big-etale}.
\end{proof}

\begin{lemma}
\label{lemma-alternative}
Soit $S$ un schéma. Notons $S_{affine, \proetale}$ la sous-catégorie pleine
de $S_\proetale$ formée des objets affines. Un recouvrement de
$S_{affine, \proetale}$ sera un recouvrement pro-\'etale standard ; voir la
définition \ref{definition-standard-proetale}.
Alors la restriction
$$
\mathcal{F} \longmapsto \mathcal{F}|_{S_{affine, \etale}}
$$
définit une équivalence de topos
$\Sh(S_\proetale) \cong \Sh(S_{affine, \proetale})$.
\end{lemma}

\begin{proof}
On peut le montrer directement à partir des définitions, et c'est un bon exercice.
Mais cela découle aussi immédiatement de
Sites, lemme \ref{sites-lemma-equivalence},
en vérifiant que le foncteur d'inclusion
$S_{affine, \proetale} \to S_\proetale$
est un foncteur spécial cocontinu (voir
Sites, définition \ref{sites-definition-special-cocontinuous-functor}).
\end{proof}

\begin{lemma}
\label{lemma-affine-alternative}
Soit $S$ un schéma affine. Notons $S_{app}$ la sous-catégorie pleine
de $S_\proetale$ formée des objets affines $U$ tels que
$\mathcal{O}(S) \to \mathcal{O}(U)$ soit ind-\'etale. Un recouvrement de
$S_{app}$ sera un recouvrement pro-\'etale standard ; voir la
définition \ref{definition-standard-proetale}.
Alors la restriction
$$
\mathcal{F} \longmapsto \mathcal{F}|_{S_{app}}
$$
définit une équivalence de topos $\Sh(S_\proetale) \cong \Sh(S_{app})$.
\end{lemma}

\begin{proof}
D'après le lemme \ref{lemma-alternative}, on peut remplacer $S_\proetale$ par
$S_{affine, \proetale}$.
Le lemme découle de Sites, lemme \ref{sites-lemma-equivalence},
en vérifiant que le foncteur d'inclusion $S_{app} \to S_{affine, \proetale}$
est un foncteur spécial cocontinu ; voir
Sites, définition \ref{sites-definition-special-cocontinuous-functor}.
Les conditions énoncées dans Sites, lemme \ref{sites-lemma-equivalence},
découlent immédiatement de la définition et des faits suivants :
(a) tout objet $U$ de $S_{affine, \proetale}$ admet un recouvrement
$\{V \to U\}$ avec $V$ ind-\'etale sur $X$
(proposition \ref{proposition-weakly-etale})
et (b) le foncteur $u$ est pleinement fidèle.
\end{proof}

\begin{lemma}
\label{lemma-proetale-subcanonical}
Soit $S$ un schéma. La topologie de chacun des sites pro-\'etales
$\Sch_\proetale$, $S_\proetale$, $(\Sch/S)_\proetale$,
$S_{affine, \proetale}$ et $(\textit{Aff}/S)_\proetale$ est sous-canonique.
\end{lemma}

\begin{proof}
Il suffit de combiner le lemme \ref{lemma-recognize-proetale-covering} avec
Descente, lemme \ref{descent-lemma-fpqc-universal-effective-epimorphisms}.
\end{proof}





\section{Objets faiblement contractiles}
\label{section-weakly-contractible}

\noindent
Dans cette section, nous démontrons le fait essentiel que nos sites pro-\'etales
contiennent de nombreux objets faiblement contractiles. En fait, la démonstration du
lemme \ref{lemma-get-many-weakly-contractible} explique la forme de la fonction
$Bound$ dans la
définition \ref{definition-big-proetale-site} (bien que, pour les lecteurs
qui ne se préoccupent pas des questions ensemblistes, cette information
soit sans contenu).

\noindent
Commençons par exprimer la notion d'anneau w-contractile en termes de
recouvrements pro-\'etales.

\begin{lemma}
\label{lemma-w-contractible-proetale-cover}
Soit $T = \Spec(A)$ un schéma affine. Les conditions suivantes sont équivalentes :
\begin{enumerate}
\item $A$ est w-contractile ;
\item tout recouvrement pro-\'etale de $T$ peut être raffiné par
un recouvrement de Zariski de la forme $T = \coprod_{i = 1, \ldots, n} U_i$.
\end{enumerate}
\end{lemma}

\begin{proof}
Supposons $A$ w-contractile. D'après le lemme \ref{lemma-proetale-affine},
il suffit de montrer que tout recouvrement pro-\'etale standard
$\{f_i : T_i \to T\}_{i = 1, \ldots, n}$ peut être raffiné par un recouvrement de Zariski de $T$.
Le morphisme $\coprod T_i \to T$ est un morphisme faiblement \'etale et surjectif
de schémas affines. D'après la définition \ref{definition-w-contractible},
il existe donc un morphisme $\sigma : T \to \coprod T_i$ au-dessus de $T$.
Alors le recouvrement de Zariski $T = \coprod \sigma^{-1}(T_i)$
raffine $\{f_i : T_i \to T\}$.

\medskip\noindent
Réciproquement, supposons (2). Si $A \to B$ est fidèlement plat et faiblement \'etale,
alors $\{\Spec(B) \to T\}$ est un recouvrement pro-\'etale.
Il existe donc un recouvrement de Zariski $T = \coprod U_i$
et des morphismes $U_i \to \Spec(B)$ au-dessus de $T$. Puisque $T = \coprod U_i$,
on obtient $T \to \Spec(B)$, c'est-à-dire un homomorphisme de $A$-algèbres $B \to A$.
Cela signifie que $A$ est w-contractile.
\end{proof}

\begin{lemma}
\label{lemma-w-contractible-is-weakly-contractible}
Soit $\Sch_\proetale$ un grand site pro-\'etale comme dans la
définition \ref{definition-big-proetale-site}.
Soit $T = \Spec(A)$ un objet affine de $\Sch_\proetale$.
Les conditions suivantes sont équivalentes :
\begin{enumerate}
\item $A$ est w-contractile,
\item $T$ est un objet faiblement contractile
(Sites, définition \ref{sites-definition-w-contractible})
de $\Sch_\proetale$, et
\item tout recouvrement pro-\'etale de $T$ peut être raffiné par
un recouvrement de Zariski de la forme $T = \coprod_{i = 1, \ldots, n} U_i$.
\end{enumerate}
\end{lemma}

\begin{proof}
Nous avons établi l'équivalence de (1) et (3) dans le
lemme \ref{lemma-w-contractible-proetale-cover}.

\medskip\noindent
Supposons (3), et soit $\mathcal{F} \to \mathcal{G}$ une surjection de faisceaux sur
$\Sch_\proetale$. Soit $s \in \mathcal{G}(T)$. Pour démontrer (2), montrons que
$s$ appartient à l'image de $\mathcal{F}(T) \to \mathcal{G}(T)$. On peut trouver un
recouvrement $\{T_i \to T\}$ de $\Sch_\proetale$ tel que $s$ se relève
en une section de $\mathcal{F}$ sur $T_i$
(Sites, définition \ref{sites-definition-sheaves-injective-surjective}).
D'après (3), on peut supposer donné un recouvrement fini
$T = \coprod_{j = 1, \ldots, m} U_j$ par des sous-ensembles ouverts et fermés
et des $t_j \in \mathcal{F}(U_j)$ dont l'image est $s|_{U_j}$.
Puisque les recouvrements de Zariski sont des recouvrements de $\Sch_\proetale$
(lemme \ref{lemma-etale-proetale}), on en conclut que
$\mathcal{F}(T) = \prod \mathcal{F}(U_j)$.
Ainsi, $t = (t_1, \ldots, t_m) \in \mathcal{F}(T)$
est une section dont l'image est $s$.

\medskip\noindent
Supposons (2). Soit $A \to D$ comme dans la
proposition \ref{proposition-find-w-contractible}.
Alors $\{V \to T\}$ est un recouvrement de $\Sch_\proetale$.
(Notons que $V = \Spec(D)$ est un objet de $\Sch_\proetale$
d'après la remarque \ref{remark-size-w-contractible},
combinée avec notre choix de la fonction
$Bound$ dans la définition \ref{definition-big-proetale-site}
et le calcul de la taille des schémas affines dans
Ensembles, lemme \ref{sets-lemma-bound-size}.)
Puisque la topologie de $\Sch_\proetale$ est sous-canonique
(lemme \ref{lemma-proetale-subcanonical}),
$h_V \to h_T$ est un morphisme surjectif de faisceaux
(Sites, lemme \ref{sites-lemma-covering-surjective-after-sheafification}).
Puisque $T$ est supposé faiblement contractile, il existe un élément
$f \in h_V(T) = \Mor(T, V)$ dont l'image dans $h_T(T)$ est $\text{id}_T$.
Ainsi, $A \to D$ admet une rétraction $\sigma : D \to A$.
Si maintenant $A \to B$ est fidèlement plat et faiblement \'etale, alors
$D \to D \otimes_A B$ possède les mêmes propriétés ; il existe donc
une rétraction $D \otimes_A B \to D$ qui, composée avec
$\sigma$, fournit une rétraction $B \to D \otimes_A B \to D \to A$
de $A \to B$. Ainsi, $A$ est w-contractile et (1) est vérifiée.
\end{proof}

\begin{lemma}
\label{lemma-get-many-weakly-contractible}
Soit $\Sch_\proetale$ un grand site pro-\'etale comme dans la
définition \ref{definition-big-proetale-site}.
Pour tout objet $T$ de $\Sch_\proetale$, il existe
un recouvrement $\{T_i \to T\}$ dans $\Sch_\proetale$
tel que chaque $T_i$ soit affine et soit le spectre d'un anneau
w-contractile. En particulier, $T_i$ est faiblement contractile dans $\Sch_\proetale$.
\end{lemma}

\begin{proof}
Pour les lecteurs qui ne se préoccupent pas des questions ensemblistes,
ce lemme est une conséquence immédiate du
lemme \ref{lemma-w-contractible-is-weakly-contractible} et de la
proposition \ref{proposition-find-w-contractible}.
Voici les détails.
Choisissons un recouvrement ouvert affine $T = \bigcup U_i$. Écrivons $U_i = \Spec(A_i)$.
Choisissons des homomorphismes d'anneaux fidèlement plats et ind-\'etales $A_i \to D_i$
tels que $D_i$ soit w-contractile comme dans la
proposition \ref{proposition-find-w-contractible}.
La famille de morphismes $\{\Spec(D_i) \to T\}$ est un
recouvrement pro-\'etale.
Si l'on peut montrer que $\Spec(D_i)$ est isomorphe à un objet, disons $T_i$,
de $\Sch_\proetale$, alors $\{T_i \to T\}$ sera combinatoirement
équivalente à un recouvrement de $\Sch_\proetale$ par la construction
de $\Sch_\proetale$ dans la définition \ref{definition-big-proetale-site}
et, plus précisément, par l'application de
Ensembles, lemme \ref{sets-lemma-coverings-site}, à la dernière étape.
Pour démontrer que $\Spec(D_i)$ est isomorphe à un objet de
$\Sch_\proetale$, il suffit de montrer que
$|D_i| \leq Bound(\text{size}(T))$ par la construction
de $\Sch_\proetale$ dans la définition \ref{definition-big-proetale-site}
et, plus précisément, par l'application de
Ensembles, lemme \ref{sets-lemma-construct-category}, à l'étape (3).
Puisque $|A_i| \leq \text{size}(U_i) \leq \text{size}(T)$
d'après Ensembles, lemmes \ref{sets-lemma-bound-affine} et
\ref{sets-lemma-bound-finite-type}, on obtient
$|D_i| \leq \kappa^{2^{2^{2^\kappa}}}$, où $\kappa = \text{size}(T)$,
d'après la remarque \ref{remark-size-w-contractible}.
Le résultat découle donc de notre choix de la fonction $Bound$ dans la
définition \ref{definition-big-proetale-site}.
\end{proof}

\begin{lemma}
\label{lemma-proetale-enough-w-contractible}
Soit $S$ un schéma. Les sites pro-\'etales
$S_\proetale$, $(\Sch/S)_\proetale$, $S_{affine, \proetale}$,
$(\textit{Aff}/S)_\proetale$, ainsi que, lorsque le schéma $S$ est affine, $S_{app}$,
possèdent suffisamment d'objets quasi-compacts, affines et faiblement contractiles ;
voir Sites, définition \ref{sites-definition-w-contractible}.
\end{lemma}

\begin{proof}
Cela découle immédiatement du lemme \ref{lemma-get-many-weakly-contractible}.
\end{proof}

\begin{lemma}
\label{lemma-weakly-contractible-cover}
Soit $S$ un schéma. Les sites pro-\'etales
$\Sch_\proetale$, $S_\proetale$ et $(\Sch/S)_\proetale$
possèdent la propriété suivante : pour tout objet
$U$, il existe un recouvrement $\{V \to U\}$ où $V$ est un
objet faiblement contractile. Si $U$ est quasi-compact, on peut
choisir $V$ affine et faiblement contractile.
\end{lemma}

\begin{proof}
Supposons que $V = \coprod_{j \in J} V_j$ soit un objet de $(\Sch/S)_\proetale$
qui soit la somme disjointe d'objets faiblement contractiles $V_j$.
Puisqu'une décomposition en somme disjointe est un recouvrement pro-\'etale,
on a $\mathcal{F}(V) = \prod_{j \in J} \mathcal{F}(V_j)$ pour tout
faisceau pro-\'etale $\mathcal{F}$. Soit $\mathcal{F} \to \mathcal{G}$
un morphisme surjectif de faisceaux d'ensembles. Puisque $V_j$ est faiblement contractile,
le morphisme $\mathcal{F}(V_j) \to \mathcal{G}(V_j)$ est surjectif ; voir
Sites, définition \ref{sites-definition-w-contractible}.
Ainsi, $\mathcal{F}(V) \to \mathcal{G}(V)$ est surjectif, car c'est un produit
de morphismes surjectifs d'ensembles ; on en conclut que $V$ est faiblement contractile.

\medskip\noindent
Choisissons un recouvrement $\{U_i \to U\}_{i \in I}$ où les $U_i$ sont affines et
faiblement contractiles, comme dans le lemme \ref{lemma-get-many-weakly-contractible}.
Prenons $V = \coprod_{i \in I} U_i$ (il y a ici une question ensembliste
que nous traiterons ci-dessous). Alors $\{V \to U\}$ est le recouvrement
pro-\'etale voulu par un objet faiblement contractile
(pour vérifier qu'il s'agit d'un recouvrement, on applique le
lemme \ref{lemma-recognize-proetale-covering}).
Si $U$ est quasi-compact, le
lemme \ref{lemma-recognize-proetale-covering}
montre immédiatement que l'on peut choisir un sous-ensemble fini $I' \subset I$ tel que
$\{U_i \to U\}_{i \in I'}$ soit encore un recouvrement ;
alors $\{\coprod_{i \in I'} U_i \to U\}$ est le recouvrement voulu
par un objet affine et faiblement contractile.

\medskip\noindent
Dans ce paragraphe, que nous conseillons vivement au lecteur de ne pas lire, nous traitons
les problèmes ensemblistes. Pour savoir que la somme disjointe
appartient à notre univers partiel, il faut borner le cardinal de
l'ensemble d'indices $I$. Il ressort immédiatement de la construction du
recouvrement $\{U_i \to U\}_{i \in I}$ dans la démonstration du
lemme \ref{lemma-get-many-weakly-contractible}
que $|I| \leq \text{size}(U)$, où la taille d'un schéma est celle
définie dans Ensembles, section \ref{sets-section-categories-schemes}.
De plus, pour tout $i$, on a
$\text{size}(U_i) \leq Bound(\text{size}(U))$;
cela résulte de la borne du cardinal de
$\Gamma(U_i, \mathcal{O}_{U_i})$ donnée dans
la démonstration du lemme \ref{lemma-get-many-weakly-contractible}
et d'Ensembles, lemme \ref{sets-lemma-bound-affine}.
Ainsi, $\text{size}(\coprod_{i \in I} U_i)) \leq Bound(\text{size}(U))$
d'après Ensembles, lemme \ref{sets-lemma-bound-size}.
D'après la construction du grand site pro-\'etale donnée dans
Ensembles, lemme \ref{sets-lemma-construct-category},
on voit donc que $\coprod_{i \in I} U_i$ est isomorphe à un objet
de notre site, ce qui achève la démonstration.
\end{proof}






\section{Hyperrecouvrements faiblement contractiles}
\label{section-hypercovering}

\noindent
Les résultats de la section \ref{section-weakly-contractible}
entraînent l'existence d'hyperrecouvrements constitués d'objets
faiblement contractiles.

\begin{lemma}
\label{lemma-w-contractible-hypercovering}
Soit $X$ un schéma.
\begin{enumerate}
\item Pour tout objet $U$ de $X_\proetale$, il existe un hyperrecouvrement
$K$ de $U$ dans $X_\proetale$ tel que chaque terme $K_n$ soit constitué d'un
seul objet faiblement contractile de $X_\proetale$ qui recouvre $U$.
\item Pour tout objet quasi-compact et quasi-séparé $U$ de $X_\proetale$,
il existe un hyperrecouvrement $K$ de $U$ dans $X_\proetale$ tel que chaque
terme $K_n$ soit constitué d'un seul objet affine et faiblement contractile de
$X_\proetale$ qui recouvre $U$.
\end{enumerate}
\end{lemma}

\begin{proof}
Soit $\mathcal{B} \subset \Ob(X_\proetale)$ l'ensemble des objets faiblement contractiles
de $X_\proetale$. Tout objet $T$ de $X_\proetale$ admet un
recouvrement $\{T_i \to T\}_{i \in I}$ où $I$ est fini et $T_i \in \mathcal{B}$, d'après le
lemme \ref{lemma-weakly-contractible-cover}.
D'après Hyperrecouvrements, lemme \ref{hypercovering-lemma-w-contractible},
on obtient un hyperrecouvrement $K$ de $U$ tel que $K_n = \{U_{n, i}\}_{i \in I_n}$,
où $I_n$ est fini et $U_{n, i}$ faiblement contractile.
On peut alors remplacer $K$ par l'hyperrecouvrement de $U$ donné par $\{U_n\}$
en degré $n$, où $U_n = \coprod_{i \in I_n} U_{n, i}$.
Cette opération est permise d'après Hyperrecouvrements,
remarque \ref{hypercovering-remark-take-unions-hypercovering-X}.

\medskip\noindent
Soit $X_{qcqs, \proetale} \subset X_\proetale$ la sous-catégorie pleine
formée des objets quasi-compacts et quasi-séparés.
Un recouvrement de $X_{qcqs, \proetale}$ sera un recouvrement pro-\'etale fini.
Alors $X_{qcqs, \proetale}$ est un site, admet des produits fibrés et
le foncteur d'inclusion $X_{qcqs, \proetale} \to X_\proetale$ est continu
et commute aux produits fibrés.
En particulier, si $K$ est un hyperrecouvrement d'un objet $U$ dans
$X_{qcqs, \proetale}$, alors $K$ est un hyperrecouvrement de $U$ dans $X_\proetale$
d'après Hyperrecouvrements, lemme
\ref{hypercovering-lemma-hypercovering-continuous-functor}.
Soit $\mathcal{B} \subset \Ob(X_{qcqs, \proetale})$ l'ensemble des objets
affines et faiblement contractiles. D'après le
lemme \ref{lemma-get-many-weakly-contractible}
et le fait qu'une somme disjointe finie de schémas affines est affine,
pour tout objet $U$ de $X_{qcqs, \proetale}$, il existe un recouvrement
$\{V \to U\}$ de $X_{qcqs, \proetale}$ tel que $V \in \mathcal{B}$.
D'après Hyperrecouvrements, lemme \ref{hypercovering-lemma-w-contractible},
on obtient un hyperrecouvrement $K$ de $U$ tel que $K_n = \{U_{n, i}\}_{i \in I_n}$,
où $I_n$ est fini et $U_{n, i}$ affine et faiblement contractile.
On peut alors remplacer $K$ par l'hyperrecouvrement de $U$
donné par $\{U_n\}$ en degré $n$, où
$U_n = \coprod_{i \in I_n} U_{n, i}$. Cette opération est permise d'après
Hyperrecouvrements, remarque \ref{hypercovering-remark-take-unions-hypercovering-X}.
\end{proof}

\noindent
Dans le lemme suivant, nous utilisons le complexe de {\v C}ech $s(\mathcal{F}(K))$
associé à un hyperrecouvrement $K$ dans un site. Voir
Hyperrecouvrements, section \ref{hypercovering-section-hyper-cech}.
Si $K$ est un hyperrecouvrement de $U$ et $K_n = \{U_n \to U\}$, alors
le complexe de {\v C}ech est de la forme suivante :
$$
s(\mathcal{F}(K)) =
\left(
\mathcal{F}(U_0) \to \mathcal{F}(U_1) \to \mathcal{F}(U_2) \to \ldots
\right)
$$
où $s(\mathcal{F}(U_n))$ est placé en degré cohomologique $n$.

\begin{lemma}
\label{lemma-compute-cohomology}
Soit $X$ un schéma. Soit $E \in D^+(X_\proetale)$ représenté par
un complexe borné inférieurement $\mathcal{E}^\bullet$ de faisceaux abéliens.
Soit $K$ un hyperrecouvrement de $U \in \Ob(X_\proetale)$ avec
$K_n = \{U_n \to U\}$, où $U_n$ est un objet faiblement contractile de
$X_\proetale$. Alors
$$
R\Gamma(U, E) = \text{Tot}(s(\mathcal{E}^\bullet(K)))
$$
dans $D(\textit{Ab})$.
\end{lemma}

\begin{proof}
Si $\mathcal{E}$ est un faisceau
abélien sur $X_\proetale$, alors la suite spectrale donnée par
Hyperrecouvrements, lemme \ref{hypercovering-lemma-cech-spectral-sequence},
implique que
$$
R\Gamma(X_\proetale, \mathcal{E}) = s(\mathcal{E}(K))
$$
car les groupes de cohomologie supérieure de tout faisceau sur $U_n$ s'annulent; voir
Cohomologie sur les sites, lemme \ref{sites-cohomology-lemma-w-contractible}.

\medskip\noindent
Si $\mathcal{E}^\bullet$ est borné inférieurement, on peut choisir une résolution
injective $\mathcal{E}^\bullet \to \mathcal{I}^\bullet$ et considérer
le morphisme de complexes
$$
\text{Tot}(s(\mathcal{E}^\bullet(K)))
\longrightarrow
\text{Tot}(s(\mathcal{I}^\bullet(K)))
$$
Pour tout $n$, le morphisme $\mathcal{E}^\bullet(U_n) \to \mathcal{I}^\bullet(U_n)$
est un quasi-isomorphisme, car le foncteur des sections sur $U_n$ est exact.
Le morphisme affiché est donc un quasi-isomorphisme d'après l'une des suites
spectrales fournies par
Homologie, lemme \ref{homology-lemma-first-quadrant-ss}.
En utilisant le résultat du premier paragraphe, on voit que pour tout $p$,
le complexe $s(\mathcal{I}^p(K))$ est acyclique en degrés $n > 0$ et
calcule $\mathcal{I}^p(U)$ en degré $0$. L'autre suite
spectrale fournie par
Homologie, lemme \ref{homology-lemma-first-quadrant-ss},
montre donc que $\text{Tot}(s(\mathcal{I}^\bullet(K)))$ calcule
$R\Gamma(U, E) = \mathcal{I}^\bullet(U)$.
\end{proof}

\begin{lemma}
\label{lemma-quasi-compact-quasi-separated-commutes-direct-sums}
Soit $X$ un schéma quasi-compact et quasi-séparé.
Le foncteur $R\Gamma(X, -) : D^+(X_\proetale) \to D(\textit{Ab})$
commute aux sommes directes et aux colimites homotopiques.
\end{lemma}

\begin{proof}
L'assertion signifie ce qui suit. Supposons donnée une famille d'objets
$E_i$ de $D^+(X_\proetale)$ telle que $\bigoplus E_i$ soit un objet
de $D^+(X_\proetale)$. Alors
$R\Gamma(X, \bigoplus E_i) = \bigoplus R\Gamma(X, E_i)$.
Pour le voir, choisissons un hyperrecouvrement $K$ de $X$ tel que $K_n = \{U_n \to X\}$,
où $U_n$ est un schéma affine et faiblement contractile; voir le
lemme \ref{lemma-w-contractible-hypercovering}.
Soit $N$ un entier tel que $H^p(E_i) = 0$ pour $p < N$.
Choisissons un complexe de faisceaux abéliens $\mathcal{E}_i^\bullet$
représentant $E_i$ tel que $\mathcal{E}_i^p = 0$ pour $p < N$.
La somme directe terme à terme $\bigoplus \mathcal{E}_i^\bullet$ représente
$\bigoplus E_i$ dans $D(X_\proetale)$; voir
Injectifs, lemme \ref{injectives-lemma-derived-products}.
D'après le lemme \ref{lemma-compute-cohomology}, on a
$$
R\Gamma(X, \bigoplus E_i) =
\text{Tot}(s((\bigoplus \mathcal{E}^\bullet_i)(K)))
$$
et
$$
R\Gamma(X, E_i) = \text{Tot}(s(\mathcal{E}^\bullet_i(K)))
$$
Puisque chaque $U_n$ est quasi-compact, on voit que
$$
\text{Tot}(s((\bigoplus \mathcal{E}^\bullet_i)(K))) =
\bigoplus \text{Tot}(s(\mathcal{E}^\bullet_i(K)))
$$
d'après Modules sur les sites, lemme
\ref{sites-modules-lemma-sections-over-quasi-compact}.
L'assertion relative aux colimites homotopiques est une conséquence formelle du fait
que $R\Gamma$ est un foncteur exact de catégories triangulées et du
fait (que nous venons de démontrer) qu'il commute aux sommes directes.
\end{proof}

\begin{remark}
\label{remark-extend-to-all}
Soit $X$ un schéma. Puisque $X_\proetale$ a suffisamment d'objets faiblement contractiles,
pour tout $K$ dans $D(X_\proetale)$, on a $K = R\lim \tau_{\geq -n}K$
d'après
Cohomologie sur les sites, proposition
\ref{sites-cohomology-proposition-enough-weakly-contractibles}.
Puisque $R\Gamma$ commute à $R\lim$ d'après
Injectifs, lemme \ref{injectives-lemma-RF-commutes-with-Rlim},
on voit que
$$
R\Gamma(X, K) = R\lim R\Gamma(X, \tau_{\geq -n}K)
$$
dans $D(\textit{Ab})$. Cela nous permettra parfois d'étendre à tous les complexes
des résultats établis pour les complexes bornés inférieurement.
\end{remark}







\section{Engendrement compact}
\label{section-compact-generation}

\noindent
Dans cette section, nous démontrons que diverses catégories dérivées associées
à nos sites pro-\'etales sont compactement engendrées au sens de
Catégories dérivées, définition \ref{derived-definition-compactly-generated}.

\begin{lemma}
\label{lemma-enough-compact-proetale}
Soit $S$ un schéma. Soit $\Lambda$ un anneau.
\begin{enumerate}
\item $D(S_\proetale)$ est compactement engendrée,
\item $D(S_\proetale, \Lambda)$ est compactement engendrée,
\item $D(S_\proetale, \mathcal{A})$ est compactement engendrée
pour tout faisceau d'anneaux $\mathcal{A}$ sur $S_\proetale$,
\item $D((\Sch/S)_\proetale)$ est compactement engendrée,
\item $D((\Sch/S)_\proetale, \Lambda)$ est compactement engendrée,
\item $D((\Sch/S)_\proetale, \mathcal{A})$ est compactement engendrée
pour tout faisceau d'anneaux $\mathcal{A}$ sur $(\Sch/S)_\proetale$.
\end{enumerate}
\end{lemma}

\begin{proof}
Démonstration de (3). Soit $U$ un objet de $S_\proetale$, affine et
faiblement contractile. Alors $j_{U!}\mathcal{A}_U$ est un objet compact
de la catégorie dérivée $D(S_\proetale, \mathcal{A})$; voir
Cohomologie sur les sites, lemme
\ref{sites-cohomology-lemma-quasi-compact-weakly-contractible-compact}.
Choisissons un ensemble $I$ et, pour chaque $i \in I$, un objet
$U_i$ de $S_\proetale$, affine et faiblement contractile, de telle sorte que tout objet affine et faiblement contractile
de $S_\proetale$ soit isomorphe à l'un des $U_i$. Cela est possible
car $\Ob(S_\proetale)$ est un ensemble. Pour achever la démonstration de (3), il suffit
de montrer que $\bigoplus j_{U_i, !}\mathcal{A}_{U_i}$ est un générateur
de $D(S_\proetale, \mathcal{A})$; voir
Catégories dérivées, définition \ref{derived-definition-generators}.
Pour le voir, soit $K$ un objet non nul
de $D(S_\proetale, \mathcal{A})$. Il existe alors un objet $T$
de notre site $S_\proetale$ et un élément non nul $\xi$ de $H^n(K)(T)$.
Autrement dit, $\xi$ est une section non nulle du $n$-ième faisceau de cohomologie
de $K$.
On peut supposer que $K$ est représenté par un complexe $\mathcal{K}^\bullet$
de faisceaux de $\mathcal{A}$-modules et que $\xi$ est la classe d'une
section $s \in \mathcal{K}^n(T)$ telle que $\text{d}(s) = 0$.
En effet, $\xi$ est localement représenté par la classe d'une section
(on obtient donc le résultat après avoir remplacé $T$ par un membre d'un recouvrement de $T$).
Choisissons ensuite un recouvrement $\{T_j \to T\}_{j \in J}$
comme dans le lemme \ref{lemma-get-many-weakly-contractible}.
Puisque $H^n(K)$ est un faisceau, on voit que, pour un certain $j$, la restriction
$\xi|_{T_j}$ reste non nulle. Ainsi, $s|_{T_j}$ définit un morphisme non nul
$j_{T_j, !}\mathcal{A}_{T_j} \to K$ dans $D(S_\proetale, \mathcal{A})$.
Comme $T_j \cong U_i$ pour un certain $i \in I$, on conclut.

\medskip\noindent
Le même argument s'applique mot pour mot au grand site pro-\'etale de $S$.
\end{proof}







\section{Comparaison des topologies}
\label{section-compare-topologies}

\noindent
Cette section est l'analogue de
Cohomologie \'etale, section \ref{etale-cohomology-section-compare-topologies}.

\begin{lemma}
\label{lemma-presheaf-value-weakly-contractible}
Soit $X$ un schéma. Soit $\mathcal{F}$ un préfaisceau d'ensembles sur $X_\proetale$
qui transforme les sommes disjointes finies en produits. Alors
$\mathcal{F}^\#(W) = \mathcal{F}(W)$ si $W$ est un objet affine et faiblement contractile
de $X_\proetale$.
\end{lemma}

\begin{proof}
Rappelons que $\mathcal{F}^\#$ est égal à $(\mathcal{F}^+)^+$ ; voir
Sites, théorème \ref{sites-theorem-plus}, où $\mathcal{F}^+$
est le préfaisceau qui associe à tout objet $U$ de $X_\proetale$
$\colim H^0(\mathcal{U}, \mathcal{F})$, où la colimite est prise
sur tous les recouvrements pro-\'etales $\mathcal{U}$ de $U$.
Il suffit donc de démontrer que (a) $\mathcal{F}^+$ transforme les sommes
disjointes finies en produits et que (b) sa valeur en $W$ est $\mathcal{F}(W)$.
Si $U = U_1 \amalg U_2$, alors, étant donné un recouvrement pro-\'etale
$\mathcal{U} = \{f_j : V_j \to U\}$ de $U$, on obtient des recouvrements pro-\'etales
$\mathcal{U}_i = \{f_j^{-1}(U_i) \to U_i\}$ et l'on a clairement
$$
H^0(\mathcal{U}, \mathcal{F}) =
H^0(\mathcal{U}_1, \mathcal{F}) \times
H^0(\mathcal{U}_2, \mathcal{F})
$$
car $\mathcal{F}$ transforme les sommes disjointes finies en produits (cela comprend
la condition selon laquelle $\mathcal{F}$ associe au schéma vide le singleton).
Cela démontre (a).
Enfin, tout recouvrement pro-\'etale de $W$ peut être raffiné par une décomposition
en somme disjointe finie $W = W_1 \amalg \ldots W_n$ d'après le
lemme \ref{lemma-w-contractible-is-weakly-contractible}.
Ainsi, $\mathcal{F}^+(W) = \mathcal{F}(W)$ précisément parce que la valeur de
$\mathcal{F}$ sur $W$ est le produit des valeurs de $\mathcal{F}$
sur les $W_j$. Cela démontre (b).
\end{proof}

\begin{lemma}
\label{lemma-small-pullback-weakly-contractible}
Soit $f : X \to Y$ un morphisme de schémas. Soit $\mathcal{F}$ un faisceau
d'ensembles sur $Y_\proetale$. Si $W$ est un objet affine et faiblement contractile
de $X_\proetale$, alors
$$
f_{small}^{-1}\mathcal{F}(W) = \colim_{W \to V} \mathcal{F}(V)
$$
où la colimite est prise sur les morphismes $W \to V$ au-dessus de $Y$
tels que $V \in Y_\proetale$.
\end{lemma}

\begin{proof}
Rappelons que $f_{small}^{-1}\mathcal{F}$ est le faisceau associé au
préfaisceau
$$
u_p\mathcal{F} : U \mapsto \colim_{U \to V} \mathcal{F}(V)
$$
sur $X_\etale$ ; voir Sites, sections \ref{sites-section-morphism-sites} et
\ref{sites-section-continuous-functors} ; nous avons omis de la notation le fait
que la colimite est prise sur la catégorie opposée de la catégorie
$\{U \to V, V \in Y_\proetale\}$. D'après le
lemme \ref{lemma-presheaf-value-weakly-contractible},
il suffit de démontrer que $u_p\mathcal{F}$ transforme les sommes disjointes finies
en produits.
Supposons que $U = U_1 \amalg U_2$ soit une somme disjointe de sous-schémas
ouverts et fermés. Il existe un foncteur
$$
\{U_1 \to V_1\} \times \{U_2 \to V_2\} \longrightarrow
\{U \to V\},\quad
(U_1 \to V_1, U_2 \to V_2) \longmapsto (U \to V_1 \amalg V_2)
$$
qui est initial (Catégories, définition \ref{categories-definition-initial}).
Le foncteur correspondant sur les catégories opposées est donc cofinal ;
d'après Catégories, lemme \ref{categories-lemma-cofinal}, on voit
que $u_p\mathcal{F}$ évalué en $U$ est la colimite des valeurs
$\mathcal{F}(V_1 \amalg V_2)$ indexées par la catégorie produit.
Puisque $\mathcal{F}$ est un faisceau, il transforme les sommes disjointes en produits et
on conclut qu'il en est de même pour $u_p\mathcal{F}$.
\end{proof}

\begin{lemma}
\label{lemma-describe-pullback}
Soit $S$ un schéma. Considérons le morphisme
$$
\pi_S : (\Sch/S)_\proetale \longrightarrow S_\proetale
$$
du lemme \ref{lemma-at-the-bottom}. Soit $\mathcal{F}$ un faisceau sur
$S_\proetale$. Alors $\pi_S^{-1}\mathcal{F}$ est donné par la règle
$$
(\pi_S^{-1}\mathcal{F})(T) = \Gamma(T_\proetale, f_{small}^{-1}\mathcal{F})
$$
où $f : T \to S$. De plus, $\pi_S^{-1}\mathcal{F}$ satisfait la
condition de faisceau relativement aux recouvrements fpqc.
\end{lemma}

\begin{proof}
Observons que l'on a un morphisme
$i_f : \Sh(T_\proetale) \to \Sh(\Sch/S)_\proetale)$
tel que $\pi_S \circ i_f = f_{small}$ en tant que morphismes
$T_\proetale \to S_\proetale$ ; voir le lemme \ref{lemma-put-in-T}.
Puisque l'image inverse est transitive, on voit que
$i_f^{-1} \pi_S^{-1}\mathcal{F} = f_{small}^{-1}\mathcal{F}$, comme souhaité.

\medskip\noindent
Soit $\{g_i : T_i \to T\}_{i \in I}$ un recouvrement fpqc. L'assertion finale
signifie ceci. Soient $\mathcal{G}$ un faisceau sur $T_\proetale$ et
des sections $s_i \in \Gamma(T_i, g_{i, small}^{-1}\mathcal{G})$ dont les images inverses
sur $T_i \times_T T_j$ coïncident, il existe une unique section $s$ de $\mathcal{G}$
sur $T$ dont l'image inverse sur $T_i$ coïncide avec $s_i$. Nous démontrerons cette
assertion lorsque $T$ est affine et que le recouvrement est donné par un unique
morphisme plat surjectif $T' \to T$ entre schémas affines et nous omettons la réduction
du cas général à ce cas.

\medskip\noindent
Soit $g : T' \to T$ un morphisme plat surjectif entre schémas affines et soit
$s' \in g_{small}^{-1}\mathcal{G}(T')$ une section telle que
$\text{pr}_0^*s' = \text{pr}_1^*s'$ sur $T' \times_T T'$.
Choisissons un morphisme surjectif faiblement \'etale $W \to T'$ où
$W$ est affine et faiblement contractile ; voir le
lemme \ref{lemma-weakly-contractible-cover}. D'après le
lemme \ref{lemma-small-pullback-weakly-contractible},
la restriction $s'|_W$ est un élément de $\colim_{W \to U} \mathcal{G}(U)$.
Choisissons $\phi : W \to U_0$ et $s_0 \in \mathcal{G}(U_0)$ correspondant à $s'$.
Choisissons un morphisme surjectif faiblement \'etale $V \to W \times_T W$ où
$V$ est affine et faiblement contractile.
Notons $a, b : V \to W$ les morphismes induits.
Puisque $a^*(s'|_W) = b^*(s'|_W)$ et que la catégorie
$\{V \to U, U \in T_\proetale\}$ est cofiltrante
(ce qui est clair ; voir néanmoins
Sites, lemme \ref{sites-lemma-directed-morphism}, en cas de doute),
on voit que les deux morphismes $\phi \circ a , \phi \circ b : V \to U_0$
doivent être égaux. D'après les résultats de
Descente, section \ref{descent-section-fpqc-universal-effective-epimorphisms}
(notamment
Descente, lemme \ref{descent-lemma-fpqc-universal-effective-epimorphisms}),
il existe donc un unique morphisme $T \to U_0$ tel que $\phi$
soit le composé de ce morphisme avec le morphisme structural $W \to T$
(nous omettons un petit détail). On peut alors prendre pour $s$ l'image inverse
de $s_0$ par ce morphisme. Nous omettons de vérifier que
$s$ a pour image inverse $s'$ sur $T'$.
\end{proof}









\section{Comparaison des grands et des petits topos}
\label{section-compare}

\noindent
Cette section est l'analogue de Cohomologie \'etale, section
\ref{etale-cohomology-section-compare}. Dans la suite, nous
noterons souvent $\mathcal{F} \mapsto \mathcal{F}|_{S_\proetale}$
le foncteur image inverse $i_S^{-1}$ correspondant au
morphisme de topos $i_S : \Sh(S_\proetale) \to \Sh((\Sch/S)_\proetale)$
du lemme \ref{lemma-at-the-bottom}.

\begin{lemma}
\label{lemma-compare-injectives}
Soit $S$ un schéma. Soit $T$ un objet de $(\Sch/S)_\proetale$.
\begin{enumerate}
\item Si $\mathcal{I}$ est injectif dans $\textit{Ab}((\Sch/S)_\proetale)$, alors
\begin{enumerate}
\item $i_f^{-1}\mathcal{I}$ est injectif dans $\textit{Ab}(T_\proetale)$,
\item $\mathcal{I}|_{S_\proetale}$ est injectif dans $\textit{Ab}(S_\proetale)$,
\end{enumerate}
\item Si $\mathcal{I}^\bullet$ est un complexe K-injectif
dans $\textit{Ab}((\Sch/S)_\proetale)$, alors
\begin{enumerate}
\item $i_f^{-1}\mathcal{I}^\bullet$ est un complexe K-injectif dans
$\textit{Ab}(T_\proetale)$,
\item $\mathcal{I}^\bullet|_{S_\proetale}$ est un complexe K-injectif dans
$\textit{Ab}(S_\proetale)$,
\end{enumerate}
\end{enumerate}
\end{lemma}

\begin{proof}
Démonstration de (1)(a) et (2)(a) : $i_f^{-1}$ est adjoint à droite
au foncteur exact $i_{f, !}$. En effet, rappelons que $i_f$ correspond
à un foncteur cocontinu $u : T_\proetale \to (\Sch/S)_\proetale$
qui est continu et commute aux produits fibrés et aux égalisateurs ; voir le
lemme \ref{lemma-put-in-T} et sa démonstration. On obtient donc $i_{f, !}$ d'après
Modules sur les sites, lemme \ref{sites-modules-lemma-g-shriek-adjoint}.
Il est montré dans Modules sur les sites, lemme
\ref{sites-modules-lemma-exactness-lower-shriek}, que ce foncteur est exact.
On conclut alors que (1)(a) et (2)(a) sont vérifiées d'après
Homologie, lemme \ref{homology-lemma-adjoint-preserve-injectives}, et
Catégories dérivées, lemme \ref{derived-lemma-adjoint-preserve-K-injectives}.

\medskip\noindent
Les parties (1)(b) et (2)(b) sont des cas particuliers de (1)(a) et (2)(a),
puisque $i_S = i_{\text{id}_S}$.
\end{proof}

\begin{lemma}
\label{lemma-compare-higher-direct-image}
Soit $f : T \to S$ un morphisme de schémas.
Pour $K$ dans $D((\Sch/T)_\proetale)$, on a
$$
(Rf_{big, *}K)|_{S_\proetale} = Rf_{small, *}(K|_{T_\proetale})
$$
dans $D(S_\proetale)$. Plus généralement, soit $S' \in \Ob((\Sch/S)_\proetale)$
de morphisme structural $g : S' \to S$. Considérons le produit fibré
$$
\xymatrix{
T' \ar[r]_{g'} \ar[d]_{f'} & T \ar[d]^f \\
S' \ar[r]^g & S
}
$$
Alors, pour $K$ dans $D((\Sch/T)_\proetale)$, on a
$$
i_g^{-1}(Rf_{big, *}K) = Rf'_{small, *}(i_{g'}^{-1}K)
$$
dans $D(S'_\proetale)$ et
$$
g_{big}^{-1}(Rf_{big, *}K) = Rf'_{big, *}((g'_{big})^{-1}K)
$$
dans $D((\Sch/S')_\proetale)$.
\end{lemma}

\begin{proof}
La première égalité découle du lemme \ref{lemma-compare-injectives}
et de (\ref{equation-compare-big-small}),
en choisissant un complexe K-injectif de faisceaux abéliens représentant $K$.
La deuxième égalité découle du lemme \ref{lemma-compare-injectives}
et du lemme
\ref{lemma-morphism-big-small-cartesian-diagram},
en choisissant un complexe K-injectif de faisceaux abéliens représentant $K$.
La troisième égalité découle de même de
Cohomologie sur les sites, lemmes \ref{sites-cohomology-lemma-cohomology-of-open} et
\ref{sites-cohomology-lemma-restrict-K-injective-to-open},
ainsi que du lemme \ref{lemma-morphism-big-small-cartesian-diagram},
en choisissant un complexe K-injectif de faisceaux abéliens représentant $K$.
\end{proof}

\noindent
Soit $S$ un schéma et soit $\mathcal{H}$ un faisceau abélien sur
$(\Sch/S)_\proetale$. Rappelons que $H^n_\proetale(U, \mathcal{H})$ désigne
la cohomologie de $\mathcal{H}$ sur un objet $U$ de $(\Sch/S)_\proetale$.

\begin{lemma}
\label{lemma-compare-cohomology-big-small}
Soit $f : T \to S$ un morphisme de schémas. Pour $K$ dans $D(S_\proetale)$,
on a
$$
H^n_\proetale(S, \pi_S^{-1}K) = H^n(S_\proetale, K)
$$
et
$$
H^n_\proetale(T, \pi_S^{-1}K) = H^n(T_\proetale, f_{small}^{-1}K).
$$
Pour $M$ dans $D((\Sch/S)_\proetale)$, on a
$$
H^n_\proetale(T, M) = H^n(T_\proetale, i_f^{-1}M).
$$
\end{lemma}

\begin{proof}
Pour démontrer la dernière égalité, représentons $M$
par un complexe K-injectif de faisceaux abéliens,
puis appliquons le lemme \ref{lemma-compare-injectives}
et explicitons les définitions. La deuxième égalité en découle,
car $i_f^{-1} \circ \pi_S^{-1} = f_{small}^{-1}$. La première
égalité est un cas particulier
de la deuxième.
\end{proof}

\begin{lemma}
\label{lemma-cohomological-descent-proetale}
Soit $S$ un schéma. Pour $K \in D(S_\proetale)$, le morphisme
$$
K \longrightarrow R\pi_{S, *}\pi_S^{-1}K
$$
est un isomorphisme.
\end{lemma}

\begin{proof}
Cela résulte du fait que $\pi_S^{-1}$ et $\pi_{S, *} = i_S^{-1}$ sont tous deux
des foncteurs exacts et que le foncteur composé $\pi_{S, *} \circ \pi_S^{-1}$
est le foncteur identité.
\end{proof}






















\section{Points du site pro-\'etale}
\label{section-points}

\noindent
Nous appliquons d'abord le critère de Deligne pour montrer qu'il y a suffisamment de points.

\begin{lemma}
\label{lemma-points-proetale}
Soit $S$ un schéma. Les sites pro-\'etales $\Sch_\proetale$,
$S_\proetale$, $(\Sch/S)_\proetale$, $S_{affine, \proetale}$ et
$(\textit{Aff}/S)_\proetale$ ont suffisamment de points.
\end{lemma}

\begin{proof}
Le grand topos pro-\'etale de $S$ est équivalent au topos défini par
$(\textit{Aff}/S)_\proetale$ ; voir le
lemme \ref{lemma-affine-big-site-proetale}.
Le topos des faisceaux sur $S_\proetale$ est équivalent au topos
associé à $S_{affine, \proetale}$ ; voir le
lemme \ref{lemma-alternative}.
Le résultat pour les sites $(\textit{Aff}/S)_\proetale$ et
$S_{affine, \proetale}$ découle immédiatement du résultat de Deligne
(Sites, lemme \ref{sites-lemma-criterion-points}).
Le cas de $\Sch_\proetale$ résulte du fait que ce site est égal à
$(\Sch/\Spec(\mathbf{Z}))_\proetale$.
\end{proof}

\noindent
Soit $S$ un schéma. Soit $\overline{s} : \Spec(k) \to S$ un point
géométrique. On appelle {\it voisinage pro-\'etale} de $\overline{s}$ un
diagramme commutatif
$$
\xymatrix{
\Spec(k) \ar[r]_-{\overline{u}} \ar[rd]_{\overline{s}} & U \ar[d] \\
& S
}
$$
avec $U \to S$ faiblement \'etale.

\begin{lemma}
\label{lemma-cofinal-etale}
Soit $S$ un schéma et soit $\overline{s} : \Spec(k) \to S$ un
point géométrique. La catégorie des voisinages pro-\'etales de
$\overline{s}$ est cofiltrante.
\end{lemma}

\begin{proof}
La démonstration est identique à celle de
Cohomologie \'etale, lemme \ref{etale-cohomology-lemma-cofinal-etale},
mais en utilisant les résultats correspondants sur les morphismes faiblement \'etales
démontrés dans
Compléments sur les morphismes, lemmes
\ref{more-morphisms-lemma-composition-weakly-etale},
\ref{more-morphisms-lemma-base-change-weakly-etale} et
\ref{more-morphisms-lemma-weakly-etale-permanence}.
\end{proof}

\begin{lemma}
\label{lemma-geometric-lift-to-cover}
Soit $S$ un schéma. Soit $\overline{s}$ un point géométrique de $S$.
Soit $\mathcal{U} = \{\varphi_i : S_i \to S\}_{i\in I}$ un
recouvrement pro-\'etale. Il existe alors $i \in I$ et un point
géométrique $\overline{s}_i$ de $S_i$ dont l'image est $\overline{s}$.
\end{lemma}

\begin{proof}
Cela découle immédiatement du fait que $\coprod \varphi_i$ est surjectif
et que les extensions de corps résiduels induites par des morphismes faiblement \'etales
sont algébriques séparables (voir, par exemple,
Compléments sur les morphismes, lemme
\ref{more-morphisms-lemma-weakly-etale-strictly-henselian-local-rings}).
\end{proof}

\noindent
Soit $S$ un schéma et soit $\overline{s}$ un point géométrique de $S$.
Pour $\mathcal{F}$ dans $\Sh(S_\proetale)$, définissons la
{\it fibre de $\mathcal{F}$ en $\overline{s}$} par la formule
$$
\mathcal{F}_{\overline{s}} = \colim_{(U, \overline{u})} \mathcal{F}(U)
$$
où la colimite est prise sur tous les voisinages pro-\'etales $(U, \overline{u})$
de $\overline{s}$ tels que $U \in \Ob(S_\proetale)$. Il résulte des
deux lemmes précédents que le foncteur
$$
S_\proetale \textit{Ens},\quad
U \longmapsto
\{\overline{u}\text{ point géométrique de }U\text{ dont l'image est }\overline{s}\}
$$
définit un point du site $S_\proetale$ ; voir
Sites, définition \ref{sites-definition-point} et
le lemme \ref{sites-lemma-neighbourhoods-cofiltered}.
Le foncteur $\mathcal{F} \mapsto \mathcal{F}_{\overline{s}}$
définit donc un point du topos $\Sh(S_\proetale)$ ; voir
Sites, définition \ref{sites-definition-point-topos} et
le lemme \ref{sites-lemma-point-site-topos}. En particulier, ce
foncteur est exact et commute aux colimites quelconques.
En fait, ce foncteur admet une autre description.

\begin{lemma}
\label{lemma-classical-point}
Dans la situation ci-dessus, le schéma $\Spec(\mathcal{O}_{S, \overline{s}}^{sh})$
est un objet de $X_\proetale$ et il existe un isomorphisme canonique
$$
\mathcal{F}(\Spec(\mathcal{O}_{S, \overline{s}}^{sh})) =
\mathcal{F}_{\overline{s}}
$$
fonctoriel en $\mathcal{F}$.
\end{lemma}

\begin{proof}
La première assertion résulte clairement de la construction de l'hensélisé strict
comme limite inductive filtrante d'algèbres \'etales sur $S$, ou de la caractérisation
des morphismes faiblement \'etales donnée dans
Compléments sur les morphismes, lemme
\ref{more-morphisms-lemma-weakly-etale-strictly-henselian-local-rings}.
La deuxième assertion découle du théorème d'Olivier
(Compléments d'Algèbre, théorème \ref{more-algebra-theorem-olivier}),
selon lequel le schéma $\Spec(\mathcal{O}_{S, \overline{s}}^{sh})$
est un objet initial de la catégorie des voisinages pro-\'etales
de $\overline{s}$.
\end{proof}

\noindent
Contrairement au cas du topos \'etale de $S$, il n'est pas vrai
que tout point de $\Sh(S_\proetale)$ soit de cette forme, ni que
la famille des points associés aux points géométriques soit
conservative. En effet, supposons que $S = \Spec(k)$, où $k$ est un
corps algébriquement clos. Soit $A$ un groupe abélien non nul.
Considérons le faisceau $\mathcal{F}$ sur $S_\proetale$ défini par
$$
\mathcal{F}(U) =
\frac{\{\text{applications }U \to A\}}{\{\text{applications localement constantes}\}}
$$
pour $U$ affine et, en général, par passage au faisceau associé ; voir
l'exemple \ref{example-functions-mod-locally-constant-functions}.
Alors $\mathcal{F}(U) = 0$ si $U = S = \Spec(k)$, mais, en général, $\mathcal{F}$
n'est pas nul. En effet, $S_\proetale$ contient des objets affines
ayant une infinité de points. Par exemple, soit $E = \lim E_n$
une limite projective d'ensembles finis dont les morphismes de transition sont surjectifs ;
on peut prendre $E = \mathbf{Z}_p = \lim \mathbf{Z}/p^n\mathbf{Z}$. Le schéma
$U = \Spec(\colim \text{Map}(E_n, k))$ est un objet de $S_\proetale$
car $\colim \text{Map}(E_n, k)$ est faiblement \'etale (et même ind-Zariski)
sur $k$. Ainsi, $\mathcal{F}(U)$ est non nul, car il existe des applications $E \to A$
qui ne sont pas localement constantes.
Ainsi, $\mathcal{F}$ est un faisceau abélien non nul dont la fibre en l'unique
point géométrique de $S$ est nulle. Puisque nous savons que $S_\proetale$
a suffisamment de points, on conclut qu'il doit exister un point du site pro-\'etale
qui ne provient pas de la construction expliquée ci-dessus.

\medskip\noindent
Pour remplacer les arguments qui utilisent les points, on emploie des objets affines
faiblement contractiles.
D'une part, il y a suffisamment d'objets affines faiblement contractiles d'après le
lemme \ref{lemma-proetale-enough-w-contractible}.
D'autre part, si $W \in \Ob(S_\proetale)$ est affine et faiblement contractile,
alors le foncteur
$$
\Sh(S_\proetale) \longrightarrow \textit{Ens},\quad
\mathcal{F} \longmapsto \mathcal{F}(W)
$$
est un foncteur exact $\Sh(S_\proetale) \to \textit{Ens}$ qui commute
à toutes les limites. Le foncteur
$$
\textit{Ab}(S_\proetale) \longrightarrow \textit{Ab},\quad
\mathcal{F} \longmapsto \mathcal{F}(W)
$$
est exact et commute aux sommes directes (car $W$ est quasi-compact ; voir
Sites, lemme \ref{sites-lemma-directed-colimits-sections}) ; il
commute donc à toutes les limites et à toutes les colimites. De plus, on peut vérifier l'exactitude d'un
complexe de faisceaux abéliens par évaluation sur ces objets affines
faiblement contractiles de $S_\proetale$ ; voir
Cohomologie sur les sites, proposition
\ref{sites-cohomology-proposition-enough-weakly-contractibles}.

\medskip\noindent
Remarquons enfin que le foncteur $\mathcal{F} \mapsto \mathcal{F}(W)$,
pour $W$ affine et faiblement contractile, n'est en général pas le foncteur fibre associé à un
point de $S_\proetale$, car il ne préserve pas les coproduits de faisceaux
d'ensembles si $W$ est non connexe. En fait, $W$ est non connexe dès que $W$
possède plus de $1$ point fermé, c'est-à-dire lorsque $W$ n'est pas le spectre d'un
anneau local strictement hensélien (qui est le cas particulier examiné ci-dessus).





\section{Comparaison avec le site \'etale}
\label{section-comparison}

\noindent
Soit $X$ un schéma. Avec des choix convenables de sites\footnote{Choisissons
un grand site pro-\'etale $\Sch_\proetale$ contenant $X$ comme dans la
définition \ref{definition-big-proetale-site}. Puis munissons $\Sch_\etale$
de la même catégorie sous-jacente que $\Sch_\proetale$, en prenant
pour recouvrements exactement les recouvrements pro-\'etales qui sont aussi des
recouvrements \'etales. Avec ces choix, soient $X_\etale$ et $X_\proetale$
les sous-catégories définies dans
la définition \ref{definition-big-small-proetale} et
Topologies, définition \ref{topologies-definition-big-small-etale}.
Comparer avec
Topologies, remarque \ref{topologies-remark-choice-sites}.},
le foncteur $u : X_\etale \to X_\proetale$ qui associe
$U/X$ à $U/X$ définit un morphisme de sites
$$
\epsilon : X_\proetale \longrightarrow X_\etale
$$
Cela résulte de Sites, proposition \ref{sites-proposition-get-morphism}.

\begin{lemma}
\label{lemma-describe-pullback-from-etale}
Avec les notations ci-dessus, soit $\mathcal{F}$ un faisceau sur $X_\etale$.
La règle
$$
X_\proetale \longrightarrow \textit{Ens},\quad
(f : Y \to X) \longmapsto \Gamma(Y_\etale, f_\etale^{-1}\mathcal{F})
$$
définit un faisceau égal à $\epsilon^{-1}\mathcal{F}$.
Ici, $f_\etale : Y_\etale \to X_\etale$ est le morphisme
entre les petits sites \'etales construit dans
Cohomologie \'etale, section \ref{etale-cohomology-section-functoriality}.
\end{lemma}

\begin{proof}
D'après le lemme \ref{lemma-recognize-proetale-covering}, tout recouvrement pro-\'etale
est un recouvrement fpqc. La formule définit donc un faisceau sur $X_\proetale$
d'après Cohomologie \'etale, lemme \ref{etale-cohomology-lemma-describe-pullback}.
Soit $a : \Sh(X_\etale) \to \Sh(X_\proetale)$ le foncteur
qui associe à $\mathcal{F}$ le faisceau donné par la formule du lemme.
Pour montrer que $a = \epsilon^{-1}$, il suffit de montrer que $a$ est un
adjoint à gauche de $\epsilon_*$.

\medskip\noindent
Soit $\mathcal{G}$ un objet de $\Sh(X_\proetale)$.
Rappelons que $\epsilon_*\mathcal{G}$ n'est autre que la restriction de
$\mathcal{G}$ à la sous-catégorie pleine $X_\etale$.
Soit $f : Y \to X$ un objet de $X_\proetale$.
Nous considérons $Y_\etale$ comme une sous-catégorie de $X_\proetale$.
Les morphismes de restriction du faisceau $\mathcal{G}$ définissent un morphisme
$$
\epsilon_*\mathcal{G} = \mathcal{G}|_{X_\etale}
\longrightarrow
f_{\etale, *}(\mathcal{G}|_{Y_\etale})
$$
En effet, pour $U$ dans $X_\etale$, la valeur de
$f_{\etale, *}(\mathcal{G}|_{Y_\etale})$ en $U$
est $\mathcal{G}(Y \times_X U)$ et il existe un morphisme de restriction
$\mathcal{G}(U) \to \mathcal{G}(Y \times_X U)$.
Par adjonction, cela détermine un morphisme
$$
f_\etale^{-1}(\epsilon_*\mathcal{G}) \to \mathcal{G}|_{Y_\etale}
$$
En réunissant ces morphismes pour tous les $f : Y \to X$ dans $X_\proetale$,
on obtient un morphisme canonique $a(\epsilon_*\mathcal{G}) \to \mathcal{G}$.

\medskip\noindent
Soit $\mathcal{F}$ un objet de $\Sh(X_\etale)$. On a manifestement
$\mathcal{F} = \epsilon_*a(\mathcal{F})$.

\medskip\noindent 
Nous affirmons que les morphismes $\mathcal{F} \to \epsilon_*a(\mathcal{F})$ et
$a(\epsilon_*\mathcal{G}) \to \mathcal{G}$
sont l'unité et la co-unité de l'adjonction (voir
Catégories, section \ref{categories-section-adjoint}).
Pour le voir, il suffit de montrer que les morphismes correspondants
$$
\Mor_{\Sh(X_\proetale)}(a(\mathcal{F}), \mathcal{G}) \to
\Mor_{\Sh(X_\etale)}(\mathcal{F}, \epsilon^{-1}\mathcal{G})
$$
et
$$
\Mor_{\Sh(X_\etale)}(\mathcal{F}, \epsilon^{-1}\mathcal{G}) \to
\Mor_{\Sh(X_\proetale)}(a(\mathcal{F}), \mathcal{G})
$$
sont inverses l'un de l'autre. Nous omettons la vérification détaillée.
\end{proof}

\begin{lemma}
\label{lemma-fully-faithful}
Soit $X$ un schéma. Pour tout faisceau $\mathcal{F}$ sur $X_\etale$,
le morphisme d'adjonction $\mathcal{F} \to \epsilon_*\epsilon^{-1}\mathcal{F}$ est un
isomorphisme, c'est-à-dire que $\epsilon^{-1}\mathcal{F}(U) = \mathcal{F}(U)$
pour $U$ dans $X_\etale$.
\end{lemma}

\begin{proof}
Cela découle immédiatement de la description de
$\epsilon^{-1}$ donnée dans le lemme \ref{lemma-describe-pullback-from-etale}.
\end{proof}

\begin{lemma}
\label{lemma-limit-pullback}
Soit $X$ un schéma. Soit $Y = \lim Y_i$ la limite d'un système projectif filtrant
d'objets quasi-compacts et quasi-séparés de $X_\proetale$
dont les morphismes de transition sont affines. Pour tout faisceau $\mathcal{F}$ sur $X_\etale$,
on a
$$
\epsilon^{-1}\mathcal{F}(Y) = \colim \epsilon^{-1}\mathcal{F}(Y_i)
$$
De plus, si chaque $Y_i$ appartient à $X_\etale$, on a
$\epsilon^{-1}\mathcal{F}(Y) = \colim \mathcal{F}(Y_i)$.
\end{lemma}

\begin{proof}
D'après la description de $\epsilon^{-1}\mathcal{F}$ dans le
lemme \ref{lemma-describe-pullback-from-etale}, la formule affichée
est un cas particulier de Cohomologie \'etale, théorème
\ref{etale-cohomology-theorem-colimit}.
(Lorsque $X$, $Y$ et les $Y_i$ sont tous affines, voir
Cohomologie \'etale, lemme
\ref{etale-cohomology-lemma-directed-colimit-cohomology}, dont l'énoncé est plus facile à lire.)
La dernière assertion découle immédiatement de ce qui précède
et du lemme \ref{lemma-fully-faithful}.
\end{proof}

\begin{lemma}
\label{lemma-affine-vanishing}
Soit $X$ un schéma affine. Pour tout faisceau abélien injectif $\mathcal{I}$ sur
$X_\etale$, on a $H^p(X_\proetale, \epsilon^{-1}\mathcal{I}) = 0$
pour $p > 0$.
\end{lemma}

\begin{proof}
Pour le démontrer, nous allons utiliser
Cohomologie sur les sites, lemme \ref{sites-cohomology-lemma-cech-vanish-collection}.
Soit $\mathcal{B} \subset \Ob(X_\proetale)$ l'ensemble des objets affines
$U$ de $X_\proetale$ tels que $\mathcal{O}(X) \to \mathcal{O}(U)$ soit
ind-\'etale. Soit $\text{Cov}$ l'ensemble des recouvrements pro-\'etales
$\{U_i \to U\}_{i = 1, \ldots, n}$ avec $U \in \mathcal{B}$ tels que
$\mathcal{O}(U) \to \mathcal{O}(U_i)$ soit ind-\'etale pour $i = 1, \ldots, n$.
Les propriétés (1) et (2) énoncées dans
Cohomologie sur les sites, lemme \ref{sites-cohomology-lemma-cech-vanish-collection},
sont vérifiées pour $\mathcal{B}$ et $\text{Cov}$ d'après les
lemmes \ref{lemma-composition-ind-etale},
\ref{lemma-base-change-ind-etale} et
\ref{lemma-proetale-affine}, ainsi que la
proposition \ref{proposition-weakly-etale}.

\medskip\noindent
Pour vérifier la condition (3), supposons que
$\mathcal{U} = \{U_i \to U\}_{i = 1, \ldots, n}$
soit un élément de $\text{Cov}$. Il faut montrer que les groupes de
cohomologie supérieure de Čech de $\epsilon^{-1}\mathcal{I}$
relativement à $\mathcal{U}$ sont nuls.
Écrivons d'abord $U_i = \lim_{a \in A_i} U_{i, a}$ sous la forme d'une limite projective filtrante,
avec $U_{i, a} \to U$ \'etale et $U_{i, a}$ affine.
Considérons $A_1 \times \ldots \times A_n$ comme un ensemble ordonné filtrant
pour l'ordre $(a_1, \ldots, a_n) \geq (a_1', \ldots, a_n')$
si et seulement si $a_i \geq a_i'$ pour $i = 1, \ldots, n$.
Observons que
$\mathcal{U}_{(a_1, \ldots, a_n)} = \{U_{i, a_i} \to U\}_{i = 1, \ldots, n}$
est un recouvrement \'etale pour tous
$a_1, \ldots, a_n \in A_1 \times \ldots \times A_n$.
Observons que
$$
U_{i_0} \times_U U_{i_1} \times_U \ldots \times_U U_{i_p}
=
\lim_{(a_1, \ldots, a_n) \in A_1 \times \ldots \times A_n}
U_{i_0, a_{i_0}} \times_U
U_{i_1, a_{i_1}} \times_U \ldots \times_U U_{i_p, a_{i_p}}
$$
pour tous $i_0, \ldots, i_p \in \{1, \ldots, n\}$, car les limites
commutent aux produits fibrés. Ainsi, d'après le lemme \ref{lemma-limit-pullback}
et l'exactitude des colimites filtrantes, on a
$$
\check{H}^p(\mathcal{U}, \epsilon^{-1}\mathcal{I}) =
\colim \check{H}^p(\mathcal{U}_{(a_1, \ldots, a_n)}, \epsilon^{-1}\mathcal{I}) 
$$
Il suffit donc de démontrer l'annulation pour les recouvrements \'etales
de $U$!

\medskip\noindent
Soit $\mathcal{U} = \{U_i \to U\}_{i = 1, \ldots, n}$ un recouvrement \'etale
avec $U_i$ affine. Écrivons $U = \lim_{b \in B} U_b$ comme limite projective filtrante,
avec $U_b$ affine et $U_b \to X$ \'etale.
D'après Limites, lemmes \ref{limits-lemma-descend-finite-presentation},
\ref{limits-lemma-limit-affine} et
\ref{limits-lemma-descend-etale},
on peut choisir $b_0 \in B$ tel que, pour $i = 1, \ldots, n$, il existe un
morphisme \'etale $U_{i, b_0} \to U_{b_0}$ entre schémas affines
tel que $U_i = U \times_{U_{b_0}} U_{i, b_0}$.
Posons $U_{i, b} = U_b \times_{U_{b_0}} U_{i, b_0}$ pour $b \geq b_0$.
Pour $b$ assez grand, la famille
$\mathcal{U}_b = \{U_{i, b} \to U_b\}_{i = 1, \ldots, n}$
est un recouvrement \'etale ; voir
Limites, lemme \ref{limits-lemma-descend-surjective}.
Comme précédemment, on obtient
$$
\check{H}^p(\mathcal{U}, \epsilon^{-1}\mathcal{I}) =
\colim \check{H}^p(\mathcal{U}_b, \epsilon^{-1}\mathcal{I}) =
\colim \check{H}^p(\mathcal{U}_b, \mathcal{I})
$$
où la dernière égalité résulte du lemme \ref{lemma-fully-faithful}.
Puisque chacun des complexes de {\v C}ech du membre de droite est
acyclique en degrés strictement positifs
(Cohomologie sur les sites, lemme
\ref{sites-cohomology-lemma-injective-trivial-cech}),
celui du membre de gauche l'est également. Cela démontre la condition (3) de
Cohomologie sur les sites, lemme \ref{sites-cohomology-lemma-cech-vanish-collection}.
Puisque $X \in \mathcal{B}$, le lemme en découle.
\end{proof}

\begin{lemma}
\label{lemma-relative-comparison}
Soit $X$ un schéma.
\begin{enumerate}
\item Pour un faisceau abélien $\mathcal{F}$ sur $X_\etale$,
on a $R\epsilon_*(\epsilon^{-1}\mathcal{F}) = \mathcal{F}$.
\item Pour $K \in D^+(X_\etale)$, le morphisme $K \to R\epsilon_*\epsilon^{-1}K$
est un isomorphisme.
\end{enumerate}
\end{lemma}

\begin{proof}
Soit $\mathcal{I}$ un faisceau abélien injectif sur $X_\etale$.
Rappelons que $R^q\epsilon_*(\epsilon^{-1}\mathcal{I})$ est le faisceau associé
à $U \mapsto H^q(U_\proetale, \epsilon^{-1}\mathcal{I})$ ; voir
Cohomologie sur les sites, lemme \ref{sites-cohomology-lemma-higher-direct-images}.
D'après le lemme \ref{lemma-affine-vanishing}, on voit qu'il est nul pour $q > 0$
et $U$ affine et \'etale sur $X$. Puisque tout objet de
$X_\etale$ admet un recouvrement par des objets affines, il s'ensuit que
$R^q\epsilon_*(\epsilon^{-1}\mathcal{I}) = 0$ pour $q > 0$.

\medskip\noindent
Soit $K \in D^+(X_\etale)$. Choisissons un complexe borné inférieurement $\mathcal{I}^\bullet$
de faisceaux abéliens injectifs sur $X_\etale$ représentant $K$.
Alors $\epsilon^{-1}K$ est représenté par $\epsilon^{-1}\mathcal{I}^\bullet$.
Le lemme d'acyclicité de Leray (Catégories dérivées, lemme
\ref{derived-lemma-leray-acyclicity}) montre que $R\epsilon_*\epsilon^{-1}K$
est représenté par $\epsilon_*\epsilon^{-1}\mathcal{I}^\bullet$.
D'après le lemme \ref{lemma-fully-faithful}, on conclut que
$R\epsilon_*\epsilon^{-1}\mathcal{I}^\bullet = \mathcal{I}^\bullet$,
ce qui achève la démonstration de (2). La partie (1) est un cas particulier de (2).
\end{proof}

\begin{lemma}
\label{lemma-compare-cohomology}
Soit $X$ un schéma.
\begin{enumerate}
\item Pour un faisceau abélien $\mathcal{F}$ sur $X_\etale$, on a
$$
H^i(X_\etale, \mathcal{F}) =
H^i(X_\proetale, \epsilon^{-1}\mathcal{F})
$$
pour tout $i$.
\item Pour $K \in D^+(X_\etale)$, on a
$$
R\Gamma(X_\etale, K) = R\Gamma(X_\proetale, \epsilon^{-1}K)
$$
\end{enumerate}
\end{lemma}

\begin{proof}
Conséquence immédiate du lemme \ref{lemma-relative-comparison} et de
la suite spectrale de Leray (Cohomologie sur les sites, lemme
\ref{sites-cohomology-lemma-apply-Leray}).
\end{proof}

\begin{lemma}
\label{lemma-compare-cohomology-nonabelian}
Soit $X$ un schéma. Soit $\mathcal{G}$ un faisceau de groupes (non nécessairement
commutatifs) sur $X_\etale$. On a
$$
H^1(X_\etale, \mathcal{G}) =
H^1(X_\proetale, \epsilon^{-1}\mathcal{G})
$$
où $H^1$ désigne l'ensemble des classes d'isomorphisme de
torseurs (voir
Cohomologie sur les sites, section \ref{sites-cohomology-section-h1-torsors}).
\end{lemma}

\begin{proof}
Comme le foncteur $\epsilon^{-1}$ est pleinement fidèle d'après le
lemme \ref{lemma-fully-faithful},
il est clair que l'application
$H^1(X_\etale, \mathcal{G}) \to H^1(X_\proetale, \epsilon^{-1}\mathcal{G})$
est injective. Pour montrer la surjectivité, il suffit de montrer que
tout $\epsilon^{-1}\mathcal{G}$-torseur $\mathcal{F}$ est localement trivial pour la topologie \'etale.
Pour le voir, nous pouvons supposer que $X$ est affine.
Nous nous ramenons donc à démontrer la surjectivité lorsque $X$ est affine.

\medskip\noindent
Choisissons un recouvrement $\{U \to X\}$ tel que (a) $U$ soit affine, (b)
que $\mathcal{O}(X) \to \mathcal{O}(U)$ soit ind-\'etale, et (c) que $\mathcal{F}(U)$
soit non vide. C'est possible d'après la proposition \ref{proposition-weakly-etale}
et parce que les
recouvrements pro-\'etales standards de $X$ sont cofinaux parmi tous les recouvrements pro-\'etales
de $X$ (lemme \ref{lemma-proetale-affine}).
Écrivons $U = \lim U_i$ comme une limite de schémas affines, chacun étant \'etale sur $X$.
Choisissons $s \in \mathcal{F}(U)$. Soit
$g \in \epsilon^{-1}\mathcal{G}(U \times_X U)$
l'unique section telle que $g \cdot \text{pr}_1^*s = \text{pr}_2^*s$ dans
$\mathcal{F}(U \times_X U)$. Alors $g$ satisfait la condition de cocycle
$$
\text{pr}_{12}^*g \cdot \text{pr}_{23}^*g = \text{pr}_{13}^*g
$$
dans $\epsilon^{-1}\mathcal{G}(U \times_X U \times_X U)$. D'après le
lemme \ref{lemma-limit-pullback},
on a
$$
\epsilon^{-1}\mathcal{G}(U \times_X U) =
\colim \mathcal{G}(U_i \times_X U_i)
$$
et
$$
\epsilon^{-1}\mathcal{G}(U \times_X U \times_X U) =
\colim \mathcal{G}(U_i \times_X U_i \times_X U_i)
$$
il existe donc un indice $i$ et un élément $g_i \in \mathcal{G}(U_i \times_X U_i)$
dont l'image est $g$ et qui satisfait la condition de cocycle.
Le cocycle $g_i$ définit alors un torseur sous $\mathcal{G}$ sur
$X_\etale$ dont l'image inverse est isomorphe à $\mathcal{F}$
par construction. Nous omettons certains détails, plus précisément le lien
entre les torseurs et les 1-cocycles, qu'il conviendrait d'ajouter au chapitre
Cohomologie sur les sites.
\end{proof}

\begin{lemma}
\label{lemma-compare-derived}
Soit $X$ un schéma. Soit $\Lambda$ un anneau.
\begin{enumerate}
\item L'image essentielle du foncteur pleinement fidèle
$\epsilon^{-1} : \textit{Mod}(X_\etale, \Lambda) \to
\textit{Mod}(X_\proetale, \Lambda)$
est une sous-catégorie de Serre faible $\mathcal{C}$.
\item Le foncteur $\epsilon^{-1}$ définit une équivalence de catégories
entre $D^+(X_\etale, \Lambda)$ et $D^+_\mathcal{C}(X_\proetale, \Lambda)$
dont un quasi-inverse est donné par $R\epsilon_*$.
\end{enumerate}
\end{lemma}

\begin{proof}
Pour démontrer (1), nous allons vérifier les conditions (1) -- (4) de
Homologie, lemme \ref{homology-lemma-characterize-weak-serre-subcategory}.
Comme $\epsilon^{-1}$ est pleinement fidèle (lemme \ref{lemma-fully-faithful})
et exact, tout est clair sauf la condition (4).
Cependant, si
$$
0 \to \epsilon^{-1}\mathcal{F}_1 \to \mathcal{G} \to
\epsilon^{-1}\mathcal{F}_2 \to 0
$$
est une suite exacte courte de faisceaux de $\Lambda$-modules sur $X_\proetale$,
alors on obtient
$$
0 \to \epsilon_*\epsilon^{-1}\mathcal{F}_1 \to \epsilon_*\mathcal{G} \to
\epsilon_*\epsilon^{-1}\mathcal{F}_2 \to
R^1\epsilon_*\epsilon^{-1}\mathcal{F}_1
$$
D'après le lemme \ref{lemma-relative-comparison},
cette suite se réduit à la suite exacte courte
$$
0 \to \mathcal{F}_1 \to \epsilon_*\mathcal{G} \to \mathcal{F}_2 \to 0
$$
En prenant l'image inverse, on trouve que $\mathcal{G} = \epsilon^{-1}\epsilon_*\mathcal{G}$.
Cela démontre (1).

\medskip\noindent
L'assertion (2) découle de l'assertion (1) et de
Cohomologie sur les sites, lemme \ref{sites-cohomology-lemma-equivalence-bounded}.
\end{proof}

\noindent
Soit $\Lambda$ un anneau. Dans
Modules sur les sites, section \ref{sites-modules-section-locally-constant},
nous avons défini la notion de faisceau localement constant de $\Lambda$-modules
sur un site. Si $M$ est un $\Lambda$-module, alors le faisceau $\underline{M}$ est
de présentation finie comme faisceau de $\underline{\Lambda}$-modules si et seulement
si $M$ est de présentation finie comme $\Lambda$-module ; voir
Modules sur les sites, lemme \ref{sites-modules-lemma-locally-constant-finite-type}.

\begin{lemma}
\label{lemma-compare-locally-constant}
Soit $X$ un schéma. Soit $\Lambda$ un anneau.
Le foncteur $\epsilon^{-1}$ définit une équivalence de catégories
$$
\left\{
\begin{matrix}
\text{faisceaux localement constants}\\
\text{de }\Lambda\text{-modules sur }X_\etale\\
\text{de présentation finie}
\end{matrix}
\right\}
\longleftrightarrow
\left\{
\begin{matrix}
\text{faisceaux localement constants}\\
\text{de }\Lambda\text{-modules sur }X_\proetale\\
\text{de présentation finie}
\end{matrix}
\right\}
$$
\end{lemma}

\begin{proof}
Soit $\mathcal{F}$ un faisceau localement constant de $\Lambda$-modules
sur $X_\proetale$ et de présentation finie. Choisissons un recouvrement pro-\'etale
$\{U_i \to X\}$ tel que $\mathcal{F}|_{U_i}$ soit constant, disons
$\mathcal{F}|_{U_i} \cong \underline{M_i}_{U_i}$.
Remarquons que $U_i \times_X U_j$ est vide si $M_i$ n'est pas isomorphe
à $M_j$.
Pour tout $\Lambda$-module $M$, posons $I_M = \{i \in I \mid M_i \cong M\}$.
Comme tout recouvrement pro-\'etale est un recouvrement fpqc et que
Descente, lemme \ref{descent-lemma-open-fpqc-covering},
on voit que $U_M = \bigcup_{i \in I_M} \Im(U_i \to X)$
est un ouvert de $X$. Alors $X = \coprod U_M$ est une décomposition en ouverts
disjoints de $X$. Nous pouvons remplacer $X$ par $U_M$, pour un certain $M$, et
supposer que $M_i = M$ pour tout $i$.

\medskip\noindent
Considérons le faisceau $\mathcal{I} = \mathit{Isom}(\underline{M}, \mathcal{F})$.
Ce faisceau est un torseur sous
$\mathcal{G} = \mathit{Isom}(\underline{M}, \underline{M})$.
D'après Modules sur les sites, lemme \ref{sites-modules-lemma-locally-constant},
on a $\mathcal{G} = \underline{G}$,
où $G = \mathit{Isom}_\Lambda(M, M)$.
Puisque les torseurs pour la topologie \'etale
et pour la topologie pro-\'etale coïncident d'après le
lemme \ref{lemma-compare-cohomology-nonabelian},
il s'ensuit que $\mathcal{I}$ admet localement sur $X$ des sections pour la topologie \'etale.
Ainsi, $\mathcal{F}$ est localement, pour la topologie \'etale, un faisceau constant, ce qui
est précisément ce qu'il fallait démontrer.
\end{proof}

\begin{lemma}
\label{lemma-compare-locally-constant-derived}
Soit $X$ un schéma. Soit $\Lambda$ un anneau noethérien.
Soient $D_{flc}(X_\etale, \Lambda)$, resp.\ $D_{flc}(X_\proetale, \Lambda)$,
les sous-catégories pleines respectives de
$D(X_\etale, \Lambda)$, resp.\ $D(X_\proetale, \Lambda)$,
formées des complexes dont les faisceaux de cohomologie sont localement
constants et de type fini comme faisceaux de $\Lambda$-modules. Alors
$$
\epsilon^{-1} :
D_{flc}^+(X_\etale, \Lambda)
\longrightarrow
D_{flc}^+(X_\proetale, \Lambda)
$$
est une équivalence de catégories.
\end{lemma}

\begin{proof}
Les catégories $D_{flc}(X_\etale, \Lambda)$ et $D_{flc}(X_\proetale, \Lambda)$
sont des sous-catégories triangulées strictement pleines et saturées de
$D(X_\etale, \Lambda)$ et $D(X_\proetale, \Lambda)$
d'après Modules sur les sites, lemme
\ref{sites-modules-lemma-kernel-finite-locally-constant}
et
Catégories dérivées, section \ref{derived-section-triangulated-sub}.
L'énoncé du lemme s'obtient en combinant les
lemmes \ref{lemma-compare-derived} et
\ref{lemma-compare-locally-constant}.
\end{proof}

\begin{lemma}
\label{lemma-compare-truncations}
Soit $X$ un schéma. Soit $\Lambda$ un anneau noethérien.
Soit $K$ un objet de $D(X_\proetale, \Lambda)$.
Posons $K_n = K \otimes_\Lambda^\mathbf{L} \underline{\Lambda/I^n}$.
Si $K_1$ est
\begin{enumerate}
\item dans l'image essentielle de
$\epsilon^{-1} :D(X_\etale, \Lambda/I) \to D(X_\proetale, \Lambda/I)$, et
\item d'amplitude de Tor dans $[a,\infty)$ pour un certain $a \in \mathbf{Z}$,
\end{enumerate}
alors $K_n$, considéré comme objet de $D(X_\proetale, \Lambda/I^n)$, vérifie (1) et (2).
\end{lemma}

\begin{proof}
Pour $K_n$, l'assertion (2) résulte du résultat plus général de
Cohomologie sur les sites, lemme \ref{sites-cohomology-lemma-bounded}.
Pour $K_n$, l'assertion (1) résulte, par récurrence sur $n$, des
triangles distingués
$$
K \otimes_\Lambda^\mathbf{L} \underline{I^n/I^{n + 1}} \to
K_{n + 1} \to
K_n \to
K \otimes_\Lambda^\mathbf{L} \underline{I^n/I^{n + 1}}[1]
$$
et de l'isomorphisme
$$
K \otimes_\Lambda^\mathbf{L} \underline{I^n/I^{n + 1}} =
K_1 \otimes_{\Lambda/I}^\mathbf{L} \underline{I^n/I^{n + 1}}
$$
ainsi que du fait démontré dans le lemme \ref{lemma-compare-derived},
que l'image essentielle de $\epsilon^{-1}$ est une sous-catégorie triangulée
de $D^+(X_\proetale, \Lambda/I^n)$.
\end{proof}

\begin{example}
\label{example-functions-mod-locally-constant-functions}
Soit $X$ un schéma. Soit $A$ un groupe abélien. Notons
$fun(-, A)$ le faisceau sur $X_\proetale$ qui associe à $U$
l'ensemble de toutes les applications $U \to A$ (au niveau des ensembles de points).
Considérons la suite de faisceaux
$$
0 \to \underline{A} \to fun(-, A) \to \mathcal{F} \to 0
$$
sur $X_\proetale$. Puisque le faisceau constant provient par image inverse du topos
final, on voit que $\underline{A} = \epsilon^{-1}\underline{A}$.
Cependant, si $A$ possède plus d'un élément, ni $fun(-, A)$
ni $\mathcal{F}$ ne proviennent du site \'etale de $X$ par image inverse.
Pour déterminer les valeurs de $\mathcal{F}$ dans certains cas, supposons que
tous les points de $X$ soient fermés et aient un corps résiduel séparablement clos,
et que $U$ soit affine. Alors tous les points de $U$ sont fermés et ont un corps
résiduel séparablement clos, et l'on a
$$
H^1_\proetale(U, \underline{A}) =
H^1_\etale(U, \underline{A}) = 0
$$
d'après le lemme \ref{lemma-compare-cohomology} et le
Cohomologie \'etale, lemme
\ref{etale-cohomology-lemma-affine-only-closed-points}.
On a donc, dans ce cas,
$$
\mathcal{F}(U) = fun(U, A)/\underline{A}(U)
$$
\end{example}





\section{Complétion dérivée dans le cas noethérien constant}
\label{section-derived-completion-noetherian}

\noindent
Nous poursuivons la discussion commencée dans Géométrie algébrique et formelle, section
\ref{algebraization-section-derived-completion};
nous supposons que le lecteur a lu au moins une partie de cette section.

\medskip\noindent
Soit $\mathcal{C}$ un site. Soit $\Lambda$ un anneau noethérien
et soit $I \subset \Lambda$ un idéal. Rappelons, d'après
Modules sur les sites, lemme \ref{sites-modules-lemma-completion-flat},
que
$$
\underline{\Lambda}^\wedge = \lim \underline{\Lambda/I^n}
$$
est une $\underline{\Lambda}$-algèbre plate et que le morphisme
$\underline{\Lambda} \to \underline{\Lambda}^\wedge$ identifie
les quotients modulo $I$. Ainsi,
Géométrie algébrique et formelle, lemme
\ref{algebraization-lemma-restriction-derived-complete-equivalence}
donne
$$
D_{comp}(\mathcal{C}, \Lambda) =
D_{comp}(\mathcal{C}, \underline{\Lambda}^\wedge)
$$
En particulier, les faisceaux de cohomologie $H^i(K)$ d'un objet $K$ de
$D_{comp}(\mathcal{C}, \Lambda)$ sont des faisceaux de
$\underline{\Lambda}^\wedge$-modules.
Pour simplifier les notations, nous travaillons souvent avec
$D_{comp}(\mathcal{C}, \Lambda)$.

\begin{lemma}
\label{lemma-naive-completion}
Soit $\mathcal{C}$ un site. Soit $\Lambda$ un anneau noethérien
et soit $I \subset \Lambda$ un idéal. L'adjoint à gauche
du foncteur d'inclusion
$D_{comp}(\mathcal{C}, \Lambda) \to D(\mathcal{C}, \Lambda)$
de Géométrie algébrique et formelle, proposition
\ref{algebraization-proposition-derived-completion}, associe à $K$ le complexe
$$
K^\wedge = R\lim(K \otimes_\Lambda^\mathbf{L} \underline{\Lambda/I^n})
$$
En particulier, $K$ est complet au sens dérivé si et seulement si
$K = R\lim(K \otimes_\Lambda^\mathbf{L} \underline{\Lambda/I^n})$.
\end{lemma}

\begin{proof}
Choisissons des générateurs $f_1, \ldots, f_r$ de $I$. D'après
Géométrie algébrique et formelle, lemme
\ref{algebraization-lemma-derived-completion-koszul},
on a
$$
K^\wedge = 
R\lim (K \otimes_\Lambda^\mathbf{L} \underline{K_n})
$$
où $K_n = K(\Lambda, f_1^n, \ldots, f_r^n)$.
Dans Compléments d'algèbre, lemme \ref{more-algebra-lemma-sequence-Koszul-complexes},
nous avons vu que les pro-systèmes $\{K_n\}$ et
$\{\Lambda/I^n\}$ de $D(\Lambda)$
sont isomorphes. Le lemme en résulte.
\end{proof}

\begin{lemma}
\label{lemma-pushforward-Noetherian-case}
Soit $\Lambda$ un anneau noethérien. Soit $I \subset \Lambda$ un idéal.
Soit $f : \Sh(\mathcal{D}) \to \Sh(\mathcal{C})$ un morphisme de topos.
Alors
\begin{enumerate}
\item $Rf_*$ envoie $D_{comp}(\mathcal{D}, \Lambda)$
dans $D_{comp}(\mathcal{C}, \Lambda)$,
\item le foncteur $Rf_* : D_{comp}(\mathcal{D}, \Lambda) \to
D_{comp}(\mathcal{C}, \Lambda)$ admet un adjoint à gauche
$Lf_{comp}^* : D_{comp}(\mathcal{C}, \Lambda) \to
D_{comp}(\mathcal{D}, \Lambda)$ qui s'obtient en faisant suivre $Lf^*$ de la
complétion dérivée,
\item $Rf_*$ commute à la complétion dérivée,
\item pour $K$ dans $D_{comp}(\mathcal{D}, \Lambda)$, on a
$Rf_*K = R\lim Rf_*(K \otimes^\mathbf{L}_\Lambda \underline{\Lambda/I^n})$.
\item pour $M$ dans $D_{comp}(\mathcal{C}, \Lambda)$, on a
$Lf^*_{comp}M =
R\lim Lf^*(M \otimes^\mathbf{L}_\Lambda \underline{\Lambda/I^n})$.
\end{enumerate}
\end{lemma}

\begin{proof}
Nous avons établi (1) et (2) dans
Géométrie algébrique et formelle, lemme
\ref{algebraization-lemma-pushforward-derived-complete-adjoint}.
Le point (3) résulte de
Géométrie algébrique et formelle, lemme
\ref{algebraization-lemma-pushforward-commutes-with-derived-completion}.
Pour (4), soit $K$ complet au sens dérivé. Alors
$$
Rf_*K = Rf_*( R\lim K \otimes^\mathbf{L}_\Lambda \underline{\Lambda/I^n}) =
R\lim Rf_*(K \otimes^\mathbf{L}_\Lambda \underline{\Lambda/I^n})
$$
où la première égalité résulte du lemme \ref{lemma-naive-completion},
et la seconde du fait que $Rf_*$ commute à $R\lim$
(Cohomologie sur les sites, lemme
\ref{sites-cohomology-lemma-Rf-commutes-with-Rlim}). Cela démontre (4).
Pour démontrer (5), le lemme \ref{lemma-naive-completion} donne
$$
Lf_{comp}^*M =
R\lim ( Lf^*M \otimes_\Lambda^\mathbf{L} \underline{\Lambda/I^n})
$$
Comme $Lf^*$ commute au produit tensoriel dérivé d'après
Cohomologie sur les sites, lemme \ref{sites-cohomology-lemma-pullback-tensor-product},
et que $Lf^*\underline{\Lambda/I^n} = \underline{\Lambda/I^n}$,
on obtient (5).
\end{proof}







\section{Complétion dérivée et objets faiblement contractiles}
\label{section-derived-completion-proetale}

\noindent
Nous poursuivons la discussion entreprise à la
section \ref{section-derived-completion-noetherian}.
Dans cette section, nous verrons comment l'existence
d'objets faiblement contractiles simplifie l'étude
des modules complets au sens dérivé.

\medskip\noindent
Soit $\mathcal{C}$ un site. Soit $\Lambda$ un anneau noethérien.
Soit $I \subset \Lambda$ un idéal. Bien que la théorie générale
de $D_{comp}(\mathcal{C}, \Lambda)$ soit tout à fait satisfaisante,
il est difficile de donner explicitement des exemples de complexes complets au sens dérivé.
Nous savons que
\begin{enumerate}
\item tout objet $M$ de $D(\mathcal{C}, \Lambda/I^n)$ devient par restriction un objet
complet au sens dérivé de $D(\mathcal{C}, \Lambda)$, et
\item pour tout $K \in D(\mathcal{C}, \Lambda)$, la complétion dérivée
$K^\wedge = R\lim (K \otimes_\Lambda^\mathbf{L} \underline{\Lambda/I^n})$
est complet au sens dérivé.
\end{enumerate}
Les objets du premier type sont trivialement complets et ne sont peut-être pas
intéressants. Le problème de (2) est que la complétion dérivée
est en général assez mystérieuse, même dans le cas $K = \underline{\Lambda}$.
En effet, la définition des limites homotopiques fournit
un triangle distingué
$$
R\lim(\underline{\Lambda/I^n}) \to
\prod \underline{\Lambda/I^n} \to
\prod \underline{\Lambda/I^n} \to
R\lim(\underline{\Lambda/I^n})[1]
$$
dans $D(\mathcal{C}, \Lambda)$, où les produits sont pris dans
$D(\mathcal{C}, \Lambda)$. Ces produits se calculent en prenant les produits
de résolutions injectives
(Injectifs, lemme \ref{injectives-lemma-derived-products}),
et l'on voit donc que le faisceau $H^p(\prod \underline{\Lambda/I^n})$
est le faisceau associé au préfaisceau
$$
U \longmapsto \prod H^p(U, \Lambda/I^n).
$$
Comme exemple explicite, si $X = \Spec(\mathbf{C}[t, t^{-1}])$,
$\mathcal{C} = X_\etale$, $\Lambda = \mathbf{Z}$, $I = (2)$ et
$p = 1$, on obtient le faisceau associé au préfaisceau
$$
U \mapsto \prod H^1(U_\etale, \mathbf{Z}/2^n\mathbf{Z})
$$
pour $U$ \'etale sur $X$. Observons que $H^1(X_\etale, \mathbf{Z}/m\mathbf{Z})$
est cyclique d'ordre $m$ et admet pour générateur $\alpha_m$, donné par le revêtement \'etale fini
de groupe $\mathbf{Z}/m\mathbf{Z}$ défini par l'équation $t = s^m$
(voir Cohomologie \'etale, section \ref{etale-cohomology-section-computation}).
Alors la section
$$
\alpha = (\alpha_{2^n}) \in \prod H^1(X_\etale, \mathbf{Z}/2^n\mathbf{Z})
$$
du préfaisceau ci-dessus ne devient nulle par restriction à aucun schéma \'etale
non vide au-dessus de $X$; ainsi, le faisceau associé au préfaisceau n'est pas nul.

\medskip\noindent
Cependant, ce phénomène ne se produit pas sur le site pro-\'etale.
La raison en est que nous avons suffisamment d'objets (quasi-compacts) faiblement contractiles.
Dans la proposition suivante, nous rassemblons quelques résultats sur la
complétion dérivée dans le cas d'un anneau noethérien constant, pour les sites qui possèdent suffisamment
d'objets faiblement contractiles (voir
Sites, définition \ref{sites-definition-w-contractible}).

\begin{proposition}
\label{proposition-enough-weakly-contractibles}
Soit $\mathcal{C}$ un site. Supposons que $\mathcal{C}$ possède suffisamment
d'objets faiblement contractiles.
Soit $\Lambda$ un anneau noethérien. Soit $I \subset \Lambda$ un idéal.
\begin{enumerate}
\item La catégorie des faisceaux de $\Lambda$-modules complets au sens dérivé est une
sous-catégorie de Serre faible de $\textit{Mod}(\mathcal{C}, \Lambda)$.
\item Un faisceau $\mathcal{F}$ de $\Lambda$-modules vérifie
$\mathcal{F} = \lim \mathcal{F}/I^n\mathcal{F}$ si et seulement si
$\mathcal{F}$ est complet au sens dérivé et $\bigcap I^n\mathcal{F} = 0$.
\item Le faisceau $\underline{\Lambda}^\wedge$ est complet au sens dérivé.
\item Si $\ldots \to \mathcal{F}_3 \to \mathcal{F}_2 \to \mathcal{F}_1$
est un système projectif de faisceaux de $\Lambda$-modules complets au sens dérivé,
alors $\lim \mathcal{F}_n$ est complet au sens dérivé.
\item Un objet $K \in D(\mathcal{C}, \Lambda)$ est complet au sens dérivé si
et seulement si chaque faisceau de cohomologie $H^p(K)$ est complet au sens dérivé.
\item Un objet $K \in D_{comp}(\mathcal{C}, \Lambda)$ est borné supérieurement
si et seulement si $K \otimes_\Lambda^\mathbf{L} \underline{\Lambda/I}$
est borné supérieurement.
\item Un objet $K \in D_{comp}(\mathcal{C}, \Lambda)$ est borné
si $K \otimes_\Lambda^\mathbf{L} \underline{\Lambda/I}$ est de dimension de
Tor finie.
\end{enumerate}
\end{proposition}

\begin{proof}
Soit $\mathcal{B} \subset \Ob(\mathcal{C})$ un sous-ensemble tel que tout
$U \in \mathcal{B}$ soit faiblement contractile et que tout objet de $\mathcal{C}$
admette un recouvrement par des éléments de $\mathcal{B}$.
Nous utiliserons les résultats de Cohomologie sur les sites,
lemme \ref{sites-cohomology-lemma-w-contractible} et
proposition \ref{sites-cohomology-proposition-enough-weakly-contractibles},
sans autre mention.

\medskip\noindent
Rappelons que $R\lim$ commute à $R\Gamma(U, -)$
(Injectifs, lemme \ref{injectives-lemma-RF-commutes-with-Rlim}).
Soit $f \in I$. Rappelons que $T(K, f)$ est la limite homotopique
du système
$$
\ldots \xrightarrow{f} K \xrightarrow{f} K \xrightarrow{f} K
$$
dans $D(\mathcal{C}, \Lambda)$. Ainsi,
$$
R\Gamma(U, T(K, f)) = T(R\Gamma(U, K), f).
$$
Puisque les isomorphismes entre objets de
$D(\mathcal{C}, \Lambda)$ se détectent par évaluation sur les $U \in \mathcal{B}$,
nous concluons qu'un objet $K$ de $D(\mathcal{C}, \Lambda)$
est complet au sens dérivé si et seulement si, pour tout $U \in \mathcal{B}$,
l'objet $R\Gamma(U, K)$ est complet au sens dérivé en tant qu'objet de $D(\Lambda)$.

\medskip\noindent
La remarque ci-dessus montre que les points (1) et (5) découlent des résultats correspondants
pour les modules sur les anneaux; voir
Compléments d'algèbre, lemmes \ref{more-algebra-lemma-hom-from-Af} et
\ref{more-algebra-lemma-serre-subcategory}.
De même, (2) se déduit de
Compléments d'algèbre, proposition
\ref{more-algebra-proposition-derived-complete-modules},
puisque $(I^n\mathcal{F})(U) = I^n \cdot \mathcal{F}(U)$
pour $U \in \mathcal{B}$ (par exactitude de l'évaluation en $U$).

\medskip\noindent
Démonstration de (4). La limite homotopique $R\lim \mathcal{F}_n$ appartient à
$D_{comp}(X, \Lambda)$ (voir la discussion qui suit
Géométrie algébrique et formelle, définition
\ref{algebraization-definition-derived-complete}).
D'après le point (5) que nous venons de démontrer, nous concluons que
$\lim \mathcal{F}_n = H^0(R\lim \mathcal{F}_n)$
est complet au sens dérivé.
Le point (3) est un cas particulier de (4).

\medskip\noindent
Démonstration de (6) et (7). Cela résulte du
lemme \ref{lemma-naive-completion},
du
lemme de Cohomologie sur les sites \ref{sites-cohomology-lemma-bounded}
et du calcul des limites homotopiques dans Cohomologie sur les sites,
proposition \ref{sites-cohomology-proposition-enough-weakly-contractibles}.
\end{proof}






\section{Cohomologie d'un point}
\label{section-cohomology-point}

\noindent
Soit $\Lambda$ un anneau noethérien complet pour la topologie adique définie par un idéal
$I \subset \Lambda$. Soit $k$ un corps. Dans cette section, nous ``calculons''
$$
H^i(\Spec(k)_\proetale, \underline{\Lambda}^\wedge)
$$
où $\underline{\Lambda}^\wedge = \lim_m \underline{\Lambda/I^m}$ comme précédemment.
Soit $k^{sep}$ une clôture séparable de $k$.
Alors
$$
\mathcal{U} = \{\Spec(k^{sep}) \to \Spec(k)\}
$$
est un recouvrement pro-\'etale de $\Spec(k)$. Nous utiliserons, pour ce recouvrement, la suite spectrale de {\v C}ech
aboutissant à la cohomologie.
Posons $U_0 = \Spec(k^{sep})$ et
\begin{align*}
U_n & =
\Spec(k^{sep}) \times_{\Spec(k)}
\Spec(k^{sep}) \times_{\Spec(k)} \ldots
\times_{\Spec(k)} \Spec(k^{sep}) \\
& =
\Spec(k^{sep} \otimes_k k^{sep} \otimes_k \ldots \otimes_k k^{sep})
\end{align*}
($n + 1$ facteurs). Observons que l'espace topologique sous-jacent
$|U_0|$ de $U_0$ est réduit à un point et que, pour $n \geq 1$, on a
$$
|U_n| = G \times \ldots \times G\quad (n\text{ facteurs})
$$
comme espaces profinis, où $G = \text{Gal}(k^{sep}/k)$. En effet, tout
point de $U_n$ a pour corps résiduel $k^{sep}$ et nous identifions
$(\sigma_1, \ldots, \sigma_n)$ au point correspondant à la
surjection
$$
k^{sep} \otimes_k k^{sep} \otimes_k \ldots \otimes_k k^{sep}
\longrightarrow k^{sep}, \quad
\lambda_0 \otimes \lambda_1 \otimes \ldots \lambda_n
\longmapsto \lambda_0 \sigma_1(\lambda_1) \ldots \sigma_n(\lambda_n)
$$
On calcule alors
\begin{align*}
R\Gamma((U_n)_\proetale, \underline{\Lambda}^\wedge)
& =
R\lim_m R\Gamma((U_n)_\proetale, \underline{\Lambda/I^m}) \\
& =
R\lim_m R\Gamma((U_n)_\etale, \underline{\Lambda/I^m}) \\
& =
\lim_m H^0(U_n, \underline{\Lambda/I^m}) \\
& = 
\text{Maps}_{cont}(G \times \ldots \times G, \Lambda)
\end{align*}
La première égalité résulte du fait que $R\Gamma$ commute aux limites dérivées
et que $\Lambda^\wedge$ est la limite dérivée des faisceaux
$\underline{\Lambda/I^m}$ d'après
la proposition \ref{proposition-enough-weakly-contractibles}.
La deuxième égalité résulte du lemme \ref{lemma-compare-cohomology}.
La troisième égalité résulte de Cohomologie \'etale, lemme
\ref{etale-cohomology-lemma-affine-only-closed-points}.
Pour la quatrième égalité, on utilise
Cohomologie \'etale, remarque
\ref{etale-cohomology-remark-constant-locally-constant-maps}
pour identifier les sections du faisceau constant $\underline{\Lambda/I^m}$.
Elle utilise ensuite le fait que $\Lambda$ est complet pour la topologie $I$-adique
et donc égal à $\lim_m \Lambda/I^m$ comme espace topologique, pour voir que
$\lim_m \text{Map}_{cont}(G, \Lambda/I^m) = \text{Map}_{cont}(G, \Lambda)$
et de même pour les puissances supérieures de $G$.
À ce stade, Cohomologie sur les sites, lemmes
\ref{sites-cohomology-lemma-cech-cohomology}  et
\ref{sites-cohomology-lemma-cech-spectral-sequence-application}
montrent que
$$
\Lambda \to \text{Maps}_{cont}(G, \Lambda) \to
\text{Maps}_{cont}(G \times G, \Lambda) \to \ldots
$$
calcule la cohomologie pro-\'etale. En d'autres termes, nous voyons que
$$
H^i(\Spec(k)_\proetale, \underline{\Lambda}^\wedge)
=
H^i_{cont}(G, \Lambda)
$$
où le membre de droite est la cohomologie continue de Tate; voir
Cohomologie \'etale, section
\ref{etale-cohomology-section-continuous-group-cohomology}.
Bien entendu, il devait en être ainsi.

\begin{lemma}
\label{lemma-more-general-point}
Soit $k$ un corps. Soit $G = \text{Gal}(k^{sep}/k)$ son groupe de
Galois absolu. Considérons en outre l'un des deux cas suivants :
\begin{enumerate}
\item soit $M$ un groupe abélien profini muni d'une action continue
de $G$, ou
\item soit $\Lambda$ un anneau noethérien et $I \subset \Lambda$ un idéal,
et soit $M$ un module $I$-adiquement complet sur $\Lambda$, muni d'une action continue
de $G$.
\end{enumerate}
Il existe alors un faisceau canonique $\underline{M}^\wedge$ sur
$\Spec(k)_\proetale$ associé à $M$ tel que
$$
H^i(\Spec(k), \underline{M}^\wedge) = H^i_{cont}(G, M)
$$
comme groupes abéliens ou comme $\Lambda$-modules.
\end{lemma}

\begin{proof}
Démonstration dans le cas (2). Posons $M_n = M/I^nM$. Alors $M = \lim M_n$, car $M$ est
supposé complet pour la topologie $I$-adique. Puisque l'action de $G$ est continue,
on obtient des actions continues de $G$ sur $M_n$. D'après Cohomologie \'etale, théorème
\ref{etale-cohomology-theorem-equivalence-sheaves-point},
cette action correspond à un faisceau (localement constant)
$\underline{M_n}$ de $\Lambda/I^n$-modules sur $\Spec(k)_\etale$.
Prenons son image inverse sur $\Spec(k)_\proetale$ par le morphisme de comparaison
$\epsilon$ et prenons la limite
$$
\underline{M}^\wedge = \lim \epsilon^{-1}\underline{M_n}
$$
pour obtenir le faisceau annoncé dans le lemme. Le même argument que
dans l'introduction de cette section fournit la comparaison
avec la cohomologie galoisienne continue de Tate.
\end{proof}






\section{Fonctorialité du site pro-\'etale}
\label{section-morphism}

\noindent
Soit $f : X \to Y$ un morphisme de schémas. Le foncteur
$Y_\proetale \to X_\proetale$, $V \mapsto X \times_Y V$
induit un morphisme de sites $f_\proetale : X_\proetale \to Y_\proetale$; voir
Sites, proposition \ref{sites-proposition-get-morphism}.
En fait, on obtient un diagramme commutatif de morphismes de sites
$$
\xymatrix{
X_\proetale \ar[r]_\epsilon \ar[d]_{f_\proetale} &
X_\etale \ar[d]^{f_\etale} \\
Y_\proetale \ar[r]^\epsilon & Y_\etale
}
$$
où $\epsilon$ est défini comme dans la section \ref{section-comparison}.
En particulier, on a
$\epsilon^{-1} f_\etale^{-1} = f_\proetale^{-1} \epsilon^{-1}$.
Voici le résultat correspondant pour le foncteur image directe.

\begin{lemma}
\label{lemma-morphism-comparison}
Soit $f : X \to Y$ un morphisme de schémas, supposé quasi-compact et
quasi-séparé.
\begin{enumerate}
\item Soit $\mathcal{F}$ un faisceau d'ensembles sur $X_\etale$. Alors on a
$f_{\proetale, *}\epsilon^{-1}\mathcal{F} =
\epsilon^{-1}f_{\etale, *}\mathcal{F}$.
\item Soit $\mathcal{F}$ un faisceau abélien sur $X_\etale$. Alors on a
$Rf_{\proetale, *}\epsilon^{-1}\mathcal{F} =
\epsilon^{-1}Rf_{\etale, *}\mathcal{F}$.
\end{enumerate}
\end{lemma}

\begin{proof}
Démonstration de (1). Soit $\mathcal{F}$ un faisceau d'ensembles sur $X_\etale$. Il existe
un morphisme canonique $\epsilon^{-1}f_{\etale, *}\mathcal{F} \to
f_{\proetale, *}\epsilon^{-1}\mathcal{F}$; voir
Sites, section \ref{sites-section-pullback}.
Pour montrer que c'est un isomorphisme, on peut travailler localement sur $Y$ pour la topologie de Zariski; ainsi,
on peut supposer $Y$ affine. Dans ce cas,
tout objet de $Y_\proetale$ admet un recouvrement par des objets $V = \lim V_i$
qui sont des limites de schémas affines $V_i$ \'etales sur $Y$ (d'après la
proposition \ref{proposition-weakly-etale},
par exemple). En évaluant le morphisme
$\epsilon^{-1}f_{\etale, *}\mathcal{F} \to
f_{\proetale, *}\epsilon^{-1}\mathcal{F}$
en $V$, on obtient un morphisme
$$
\colim \Gamma(X \times_Y V_i, \mathcal{F})
\longrightarrow
\Gamma(X \times_Y V, \epsilon^*\mathcal{F})
$$
Le membre de gauche est décrit dans le lemme \ref{lemma-limit-pullback}.
D'après le lemme \ref{lemma-describe-pullback-from-etale}, on a
$$
\Gamma(X \times_Y V, \epsilon^*\mathcal{F}) =
\Gamma(X \times_Y V, g_\etale^{-1}\mathcal{F})
$$
où $g : X \times_Y V \to X$ est la projection. Le résultat découle donc de
Cohomologie \'etale, lemme
\ref{etale-cohomology-lemma-directed-colimit-cohomology}.

\medskip\noindent
Démonstration de (2). En raisonnant exactement comme ci-dessus,
on voit qu'il suffit de montrer que le morphisme
$$
\colim H^i_\etale(X \times_Y V_i, \mathcal{F})
\longrightarrow
H^i_\etale(X \times_Y V, \mathcal{F})
$$
est un isomorphisme, ce qui résulte encore de Cohomologie \'etale, lemme
\ref{etale-cohomology-lemma-directed-colimit-cohomology}.
\end{proof}







\section{Morphismes finis et sites pro-\'etales}
\label{section-finite}

\noindent
Il n'est pas clair qu'un morphisme fini de schémas détermine
un foncteur image directe exact sur les faisceaux abéliens pro-\'etales.

\begin{lemma}
\label{lemma-finite}
Soit $f : Z \to X$ un morphisme fini de schémas qui est
localement de présentation finie. Alors
$f_{\proetale, *} : \textit{Ab}(Z_\proetale) \to \textit{Ab}(X_\proetale)$
est exact.
\end{lemma}

\begin{proof}
Pour le démontrer, on peut travailler localement sur $X$ pour la topologie de Zariski et supposer que $X$
est affine, écrivons $X = \Spec(A)$. Alors $Z = \Spec(B)$ pour une certaine
$A$-algèbre $B$ finie et de présentation finie. La construction donnée dans la démonstration de la
proposition \ref{proposition-find-w-contractible}
produit un homomorphisme d'anneaux fidèlement plat et ind-\'etale $A \to D$
avec $D$ w-contractile. On peut vérifier l'exactitude d'une suite de
faisceaux par évaluation sur un tel objet $U = \Spec(D)$. Alors
$f_{\proetale, *}\mathcal{F}$
évalué en $U$ est égal à $\mathcal{F}$ évalué en
$V = \Spec(D \otimes_A B)$. Puisque $D \otimes_A B$ est w-contractile
d'après le lemme \ref{lemma-finite-finitely-presented-over-w-contractible},
l'évaluation en $V$ est exacte.
\end{proof}








\section{Immersions fermées et sites pro-\'etales}
\label{section-closed-immersion}

\noindent
Il n'est pas clair (et c'est probablement faux) qu'une immersion fermée de schémas
détermine un foncteur image directe exact sur les faisceaux abéliens pro-\'etales.

\begin{lemma}
\label{lemma-closed-immersion-affines}
Soit $i : Z \to X$ une immersion fermée de schémas affines.
Notons $X_{app}$ et $Z_{app}$ les sites introduits dans le
lemme \ref{lemma-affine-alternative}.
Le foncteur de changement de base
$$
u : X_{app} \to Z_{app},\quad U \longmapsto u(U) = U \times_X Z
$$
est continu et admet un adjoint à gauche pleinement fidèle $v$.
Pour $V$ dans $Z_{app}$, le morphisme $V \to v(V)$ est une immersion fermée
qui identifie $V$ à $u(v(V)) = v(V) \times_X Z$, et tout point de
$v(V)$ se spécialise en un point de $V$.
Le foncteur $v$ est cocontinu et transforme les recouvrements en recouvrements.
\end{lemma}

\begin{proof}
L'existence de l'adjoint découle immédiatement du
lemme \ref{lemma-lift-ind-etale} et des définitions.
Il est clair que $u$ est continu d'après la définition des
recouvrements dans $X_{app}$.

\medskip\noindent
Écrivons $X = \Spec(A)$ et $Z = \Spec(A/I)$. Soit $V = \Spec(\overline{C})$
un objet de $Z_{app}$, et soit $v(V) = \Spec(C)$.
Nous avons vu dans l'énoncé du lemme \ref{lemma-lift-ind-etale}
que $V$ est égal à $v(V) \times_X Z = \Spec(C/IC)$.
Tout $g \in C$ dont l'image dans
$C/IC = \overline{C}$ est inversible est lui-même inversible dans $C$. En effet, on a les
homomorphismes de $A$-algèbres $C \to C_g \to C/IC$ et, par adjonction,
on obtient un homomorphisme de $C$-algèbres $C_g \to C$.
Ainsi, tout point de $v(V)$ se spécialise en un point de $V$.

\medskip\noindent
Supposons que $\{V_i \to V\}$ soit un recouvrement dans $Z_{app}$.
Alors $\{v(V_i) \to v(V)\}$ est une famille finie de morphismes de
$Z_{app}$ telle que tout point de $V \subset v(V)$ appartienne
à l'image de l'un des morphismes $v(V_i) \to v(V)$. Comme les
morphismes $v(V_i) \to v(V)$ sont plats (puisqu'ils sont faiblement \'etales),
on conclut que $\{v(V_i) \to v(V)\}$ est conjointement surjective.
Cela démontre que $v$ transforme les recouvrements en recouvrements.

\medskip\noindent
Soit $V$ un objet de $Z_{app}$, et soit $\{U_i \to v(V)\}$
un recouvrement dans $X_{app}$. On voit alors que
$\{u(U_i) \to u(v(V)) = V\}$ est un recouvrement dans $Z_{app}$.
Par adjonction, on obtient des morphismes $v(u(U_i)) \to U_i$.
Ainsi, la famille $\{v(u(U_i)) \to v(V)\}$ raffine le
recouvrement donné, et l'on conclut que $v$ est cocontinu.
\end{proof}

\begin{lemma}
\label{lemma-closed-immersion-affines-apply}
Soit $Z \to X$ une immersion fermée de schémas affines.
Le morphisme de topos correspondant $i = i_\proetale$
est égal au morphisme de topos
associé au foncteur pleinement fidèle et cocontinu
$v : Z_{app} \to X_{app}$ du lemme \ref{lemma-closed-immersion-affines}.
Il en résulte que
\begin{enumerate}
\item $i^{-1}\mathcal{F}$ est le faisceau associé au préfaisceau
$V \mapsto \mathcal{F}(v(V))$,
\item pour un objet faiblement contractile $V$ de $Z_{app}$, on a
$i^{-1}\mathcal{F}(V) = \mathcal{F}(v(V))$,
\item $i^{-1} : \Sh(X_\proetale) \to \Sh(Z_\proetale)$
admet un adjoint à gauche $i^{Sh}_!$,
\item $i^{-1} : \textit{Ab}(X_\proetale) \to \textit{Ab}(Z_\proetale)$
admet un adjoint à gauche $i_!$,
\item $\text{id} \to i^{-1}i^{Sh}_!$, $\text{id} \to i^{-1}i_!$, et
$i^{-1}i_* \to \text{id}$ sont des isomorphismes, et
\item $i_*$, $i^{Sh}_!$ et $i_!$ sont pleinement fidèles.
\end{enumerate}
\end{lemma}

\begin{proof}
D'après le lemme \ref{lemma-affine-alternative}, on peut décrire $i_\proetale$ au moyen
du morphisme de sites $u : X_{app} \to Z_{app}$, $V \mapsto V \times_X Z$.
La première assertion du lemme découle de
Sites, lemmes \ref{sites-lemma-have-functor-other-way} et
\ref{sites-lemma-have-functor-other-way-morphism}
(mais en inversant les rôles de $u$ et $v$).

\medskip\noindent
Démonstration de (1). D'après la description de $i$ comme morphisme de topos associé
à $v$, cela résulte de la construction; voir
Sites, lemme \ref{sites-lemma-cocontinuous-morphism-topoi}.

\medskip\noindent
Démonstration de (2). Puisque le foncteur $v$ transforme les recouvrements en recouvrements d'après le
lemme \ref{lemma-closed-immersion-affines}, on voit que le préfaisceau
$\mathcal{G} : V \mapsto \mathcal{F}(v(V))$ est un préfaisceau séparé
(Sites, définition \ref{sites-definition-separated}). Par conséquent,
le faisceau associé à $\mathcal{G}$ est $\mathcal{G}^+$; voir
Sites, théorème \ref{sites-theorem-plus}. Soit ensuite $V$ un objet
contractile de $Z_{app}$. Soit
$\mathcal{V} = \{V_i \to V\}_{i = 1, \ldots, n}$
un recouvrement quelconque dans $Z_{app}$. Posons $\mathcal{V}' = \{\coprod V_i \to V\}$.
Puisque $v$ commute aux sommes disjointes finies (en tant qu'adjoint à gauche ou par
construction) et que $\mathcal{F}$ transforme les sommes
disjointes finies en produits, on voit que
$$
H^0(\mathcal{V}, \mathcal{G}) = H^0(\mathcal{V}', \mathcal{G})
$$
(notations de Sites, section \ref{sites-section-sheafification};
comparer avec
Cohomologie \'etale, lemme \ref{etale-cohomology-lemma-cech-complex}).
On peut donc supposer que le recouvrement est donné par un unique morphisme, disons
$\{V' \to V\}$. Puisque $V$ est faiblement contractile, ce recouvrement
peut être raffiné par le recouvrement trivial $\{V \to V\}$.
Il en résulte que la valeur de $\mathcal{G}^+ = i^{-1}\mathcal{F}$
en $V$ est simplement $\mathcal{F}(v(V))$, ce qui démontre (2).

\medskip\noindent
Démonstration de (3). Tout objet de $Z_{app}$ admet un recouvrement par des objets
contractiles (lemme \ref{lemma-proetale-enough-w-contractible}).
D'après ce qui précède, on aurait $i^{Sh}_!h_V = h_{v(V)}$ pour $V$
faiblement contractile si $i^{Sh}_!$ existait. L'existence de
$i^{Sh}_!$ découle alors de
Sites, lemme \ref{sites-lemma-existence-lower-shriek}.

\medskip\noindent
Démonstration de (4). L'existence de $i_!$ s'obtient de la même manière en posant
$i_!\mathbf{Z}_V = \mathbf{Z}_{v(V)}$ pour $V$ faiblement contractile dans $Z_{app}$,
en procédant de même pour les sommes directes, et en appliquant le
lemme d'Homologie \ref{homology-lemma-partially-defined-adjoint}.
Nous omettons les détails.

\medskip\noindent
Démonstration de (5). Soit $V$ un objet contractile de $Z_{app}$.
Alors $i^{-1}i^{Sh}_!h_V = i^{-1}h_{v(V)} = h_{u(v(V))} = h_V$.
(C'est un fait général que $i^{-1}h_U = h_{u(U)}$.) Puisque les
faisceaux $h_V$, pour $V$ contractile, engendrent $\Sh(Z_{app})$
(Sites, lemme \ref{sites-lemma-sheaf-coequalizer-representable}),
on conclut que $\text{id} \to i^{-1}i^{Sh}_!$ est un isomorphisme.
Il en va de même du morphisme $\text{id} \to i^{-1}i_!$. Ensuite,
$(i^{-1}i_*\mathcal{H})(V) = i_*\mathcal{H}(v(V)) =
\mathcal{H}(u(v(V))) = \mathcal{H}(V)$, et l'on trouve que
$i^{-1}i_* \to \text{id}$ est un isomorphisme.

\medskip\noindent
Les assertions de pleine fidélité de (6) résultent alors de
Catégories, lemme \ref{categories-lemma-adjoint-fully-faithful}.
\end{proof}

\begin{lemma}
\label{lemma-closed-immersion}
Soit $i : Z \to X$ une immersion fermée de schémas. Alors
\begin{enumerate}
\item $i_\proetale^{-1}$ commute aux limites,
\item $i_{\proetale, *}$ est pleinement fidèle, et
\item $i_\proetale^{-1}i_{\proetale, *} \cong \text{id}_{\Sh(Z_\proetale)}$.
\end{enumerate}
\end{lemma}

\begin{proof}
Les assertions (2) et (3) sont équivalentes d'après le
lemme de Sites \ref{sites-lemma-exactness-properties}.
Les parties (1) et (3) sont locales sur $X$ pour la topologie de Zariski; on peut donc supposer que
$X$ est affine. Dans ce cas, le résultat découle du
lemme \ref{lemma-closed-immersion-affines-apply}.
\end{proof}

\begin{lemma}
\label{lemma-thickening}
Soit $i : Z \to X$ un morphisme de schémas entier, universellement injectif et surjectif.
Alors
$i_{\proetale, *}$ et $i_\proetale^{-1}$ définissent des équivalences de catégories
quasi inverses entre les topos pro-\'etales.
\end{lemma}

\begin{proof}
On se ramène immédiatement au cas où $X$ est affine.
Alors $Z$ est également affine. Posons $A = \mathcal{O}(X)$ et $B = \mathcal{O}(Z)$.
Les catégories des algèbres \'etales sur
$A$ et $B$ sont alors équivalentes; voir 
Cohomologie \'etale, théorème
\ref{etale-cohomology-theorem-topological-invariance} et la
remarque \ref{etale-cohomology-remark-affine-inside-equivalence}.
Les catégories des algèbres ind-\'etales sur $A$ et $B$ sont donc
équivalentes. Autrement dit, les catégories $X_{app}$ et $Z_{app}$
du lemme \ref{lemma-affine-alternative} sont équivalentes.
Nous omettons de vérifier
que cette équivalence transforme les recouvrements en recouvrements et réciproquement.
Le résultat découle alors du lemme \ref{lemma-affine-alternative},
qui identifie le topos pro-\'etale au topos des faisceaux sur $X_{app}$.
\end{proof}

\begin{lemma}
\label{lemma-compute-i-star}
Soit $i : Z \to X$ une immersion fermée de schémas.
Soit $U \to X$ un objet de $X_\proetale$ tel que
\begin{enumerate}
\item $U$ soit affine et faiblement contractile, et
\item tout point de $U$ se spécialise en un point de $U \times_X Z$.
\end{enumerate}
Alors $i_\proetale^{-1}\mathcal{F}(U \times_X Z) = \mathcal{F}(U)$
pour tout faisceau abélien sur $X_\proetale$.
\end{lemma}

\begin{proof}
Puisque l'image inverse commute à la restriction, on peut remplacer $X$ par $U$.
Nous pouvons donc supposer que $X$ est affine et faiblement contractile
et que tout point de $X$ se spécialise en un point de $Z$.
D'après le lemme \ref{lemma-closed-immersion-affines-apply}, partie (1),
il suffit de montrer que $v(Z) = X$ dans ce cas.
Nous devons donc montrer ceci : si $A$ est un anneau w-contractile, si $I \subset A$
est un idéal contenu dans le radical de Jacobson de $A$ et si $A \to B \to A/I$
est une factorisation telle que $A \to B$ soit ind-\'etale, alors il existe
une unique rétraction $B \to A$ compatible avec les morphismes vers $A/I$.
Observons que $B/IB = A/I \times R$ comme $A/I$-algèbres.
Après avoir remplacé $B$ par une localisation, nous pouvons supposer que $B/IB = A/I$.
Notons que $\Spec(B) \to \Spec(A)$ est surjectif, puisque son image
contient $V(I)$, donc tous les points fermés, et qu'elle est stable par
spécialisation. Comme $A$ est w-contractile, il existe une rétraction $B \to A$.
Puisque $B/IB = A/I$, cette rétraction est compatible avec le morphisme vers $A/I$.
Nous omettons la démonstration de l'unicité (indication : utiliser le fait que $A$ et $B$ ont
des anneaux locaux isomorphes en tout idéal maximal de $A$).
\end{proof}

\begin{lemma}
\label{lemma-closed-immersion-complement-retrocompact-exact}
Soit $i : Z \to X$ une immersion fermée de schémas.
Si $X \setminus i(Z)$ est un ouvert rétrocompact de $X$, alors
$i_{\proetale, *}$ est exact.
\end{lemma}

\begin{proof}
La question est locale sur $X$; nous pouvons donc supposer que $X$ est affine.
Écrivons $X = \Spec(A)$ et $Z = \Spec(A/I)$. Il existe
$f_1, \ldots, f_r \in I$ tels que $Z = V(f_1, \ldots, f_r)$
en tant qu'ensembles ; voir Algèbre, lemme \ref{algebra-lemma-qc-open}.
D'après le lemme \ref{lemma-thickening}, nous pouvons supposer que
$Z = \Spec(A/(f_1, \ldots, f_r))$. Dans ce cas,
le foncteur $i_{\proetale, *}$ est exact d'après le
lemme \ref{lemma-finite}.
\end{proof}





\section{Prolongement par zéro}
\label{section-extension-by-zero}

\noindent
Les résultats généraux de
Modules sur les sites, section \ref{sites-modules-section-localize},
nous permettent de poser la définition suivante.

\begin{definition}
\label{definition-extension-zero}
Soit $j : U \to X$ un morphisme de schémas faiblement \'etale.
\begin{enumerate}
\item Le foncteur de restriction
$j^{-1} : \Sh(X_\proetale) \to \Sh(U_\proetale)$
admet pour adjoint à gauche
$j_!^{Sh} : \Sh(X_\proetale) \to \Sh(U_\proetale)$.
\item Le foncteur de restriction
$j^{-1} : \textit{Ab}(X_\proetale) \to \textit{Ab}(U_\proetale)$
admet un adjoint à gauche, noté
$j_! : \textit{Ab}(U_\proetale) \to \textit{Ab}(X_\proetale)$
et appelé {\it prolongement par zéro}.
\item Soit $\Lambda$ un anneau. Le foncteur
$j^{-1} : \textit{Mod}(X_\proetale, \Lambda) \to
\textit{Mod}(U_\proetale, \Lambda)$
admet un adjoint à gauche, noté
$j_! : \textit{Mod}(U_\proetale, \Lambda) \to
\textit{Mod}(X_\proetale, \Lambda)$
et appelé {\it prolongement par zéro}.
\end{enumerate}
\end{definition}

\noindent
Comme d'habitude, comparons cette situation au cas \'etale.

\begin{lemma}
\label{lemma-jshriek-comparison}
Soit $j : U \to X$ un morphisme de schémas \'etale.
Soit $\mathcal{G}$ un faisceau abélien sur $U_\etale$.
Alors $\epsilon^{-1} j_!\mathcal{G} = j_!\epsilon^{-1}\mathcal{G}$
comme faisceaux sur $X_\proetale$.
\end{lemma}

\begin{proof}
Cela est vrai parce que les deux foncteurs sont adjoints à gauche de
$j_\proetale^{-1} \epsilon_* = \epsilon_* j_\etale^{-1}$.
L'égalité résulte de la discussion à la section \ref{section-morphism}.
\end{proof}

\begin{lemma}
\label{lemma-jshriek-zero}
Soit $j : U \to X$ un morphisme de schémas faiblement \'etale.
Soit $i : Z \to X$ une immersion fermée telle que $U \times_X Z = \emptyset$.
Soit $V \to X$ un objet affine de $X_\proetale$ tel que tout point
de $V$ se spécialise en un point de $V_Z = Z \times_X V$.
Alors $j_!\mathcal{F}(V) = 0$ pour tout faisceau abélien sur $U_\proetale$.
\end{lemma}

\begin{proof}
Soit $\{V_i \to V\}$ un recouvrement pro-\'etale. Il suffit de pouvoir
raffiner ce recouvrement en un recouvrement dont les membres n'admettent aucun
morphisme vers $U$ au-dessus de $X$ (voir la construction de $j_!$ dans
Modules sur les sites, section \ref{sites-modules-section-localize}).
Raffinons d'abord le recouvrement pour obtenir un recouvrement fini dont les $V_i$ sont affines.
Pour tout $i$, écrivons $V_i = \Spec(A_i)$ et notons $Z_i \subset V_i$ le
sous-schéma image inverse de $Z$.
Posons $W_i = \Spec(A_{i, Z_i}^\sim)$ avec les notations du
lemme \ref{lemma-localization}.
Alors $\coprod W_i \to V$ est faiblement \'etale et son image contient tous les
points de $V_Z$. Elle contient donc tous les points de $V$ d'après
notre hypothèse sur les spécialisations. Ainsi, $\{W_i \to V\}$ est un
recouvrement pro-\'etale qui raffine le recouvrement donné. Mais tout point de $W_i$
se spécialise en un point situé au-dessus de $Z$; il n'existe donc aucun morphisme
$W_i \to U$ au-dessus de $X$.
\end{proof}

\begin{lemma}
\label{lemma-open-immersion}
Soit $j : U \to X$ une immersion ouverte de schémas.
Alors $\text{id} \cong j^{-1}j_!$ et $j^{-1}j_* \cong \text{id}$,
et les foncteurs $j_!$ et $j_*$ sont pleinement fidèles.
\end{lemma}

\begin{proof}
Voir Modules sur les sites, lemme \ref{sites-modules-lemma-restrict-back}
(et Sites, lemme \ref{sites-lemma-restrict-back} pour le cas
des faisceaux d'ensembles), ainsi que
Catégories, lemme \ref{categories-lemma-adjoint-fully-faithful}.
\end{proof}

\noindent
Voici la relation entre le prolongement par zéro et la restriction
au sous-schéma fermé complémentaire.

\begin{lemma}
\label{lemma-ses-associated-to-open}
Soit $X$ un schéma. Soit $Z \subset X$ un sous-schéma fermé et soit
$U \subset X$ son complémentaire. Notons $i : Z \to X$ et $j : U \to X$
les morphismes d'inclusion. Supposons que $j$ soit quasi-compact.
Pour tout faisceau abélien sur $X_\proetale$, il existe une suite exacte courte
canonique
$$
0 \to j_!j^{-1}\mathcal{F} \to \mathcal{F} \to i_*i^{-1}\mathcal{F} \to 0
$$
sur $X_\proetale$, où tous les foncteurs se rapportent à la topologie pro-\'etale.
\end{lemma}

\begin{proof}
Les propriétés d'adjonction des foncteurs en jeu fournissent les morphismes.
Il suffit de montrer que $X_\proetale$ admet suffisamment d'objets
(Sites, définition \ref{sites-definition-w-contractible}) sur lesquels
la suite s'évalue en une suite exacte courte.
Soit $V = \Spec(A)$ un objet affine de $X_\proetale$
tel que $A$ soit w-contractile (il y a suffisamment d'objets de ce type).
Alors $V \times_X Z$ est défini par un idéal $I \subset A$.
L'hypothèse que $j$ est quasi-compact implique qu'il existe
$f_1, \ldots, f_r \in I$ tels que $V(I) = V(f_1, \ldots, f_r)$.
Nous obtenons un homomorphisme d'anneaux fidèlement plat et ind-Zariski
$$
A \longrightarrow A_{f_1} \times \ldots \times A_{f_r} \times
A_{V(I)}^\sim
$$
où $A_{V(I)}^\sim$ est défini comme dans le lemme \ref{lemma-localization}.
Puisque $V_i = \Spec(A_{f_i}) \to X$ se factorise par $U$, on a
$$
j_!j^{-1}\mathcal{F}(V_i) = \mathcal{F}(V_i)
\quad\text{et}\quad
i_*i^{-1}\mathcal{F}(V_i) = 0
$$
D'autre part, pour le schéma $V^\sim = \Spec(A_{V(I)}^\sim)$,
on a
$$
j_!j^{-1}\mathcal{F}(V^\sim) = 0
\quad\text{et}\quad
\mathcal{F}(V^\sim) = i_*i^{-1}\mathcal{F}(V^\sim)
$$
La première égalité résulte du lemme \ref{lemma-jshriek-zero},
et la seconde des
lemmes \ref{lemma-compute-i-star} et \ref{lemma-localization-w-contractible}.
La suite s'évalue donc en une suite exacte sur
$\Spec(A_{f_1} \times \ldots \times A_{f_r} \times A_{V(I)}^\sim)$,
ce qui démontre le lemme.
\end{proof}

\begin{lemma}
\label{lemma-j-shriek-limits}
Soit $j : U \to X$ une immersion ouverte quasi-compacte
de schémas. Le foncteur
$j_! : \textit{Ab}(U_\proetale) \to \textit{Ab}(X_\proetale)$
commute aux limites.
\end{lemma}

\begin{proof}
Puisque $j_!$ est exact, il suffit de montrer que $j_!$ commute aux produits.
La question est locale sur $X$; nous pouvons donc supposer que $X$ est affine.
Soit $\mathcal{G}$ un faisceau abélien sur $U_\proetale$.
On a $j^{-1}j_*\mathcal{G} = \mathcal{G}$. En appliquant alors
la suite exacte du lemme \ref{lemma-ses-associated-to-open},
on obtient
$$
0 \to j_!\mathcal{G} \to j_*\mathcal{G} \to i_*i^{-1}j_*\mathcal{G} \to 0
$$
où $i : Z \to X$ est l'inclusion du schéma muni de la structure réduite
induite sur le complément $Z = X \setminus U$.
Les foncteurs $j_*$ et $i_*$ commutent aux produits en tant qu'adjoints à droite.
Le foncteur $i^{-1}$ commute aux produits d'après le
lemme \ref{lemma-closed-immersion}.
Il en va donc de même pour $j_!$, car, sur le site pro-\'etale, les produits
sont exacts
(Cohomologie sur les sites, proposition
\ref{sites-cohomology-proposition-enough-weakly-contractibles}).
\end{proof}




\section{Faisceaux constructibles sur le site pro-\'etale}
\label{section-constructible}

\noindent
Nous nous limitons aux faisceaux constructibles de $\Lambda$-modules pour un
anneau noethérien. Nous envisageons d'étudier ultérieurement les faisceaux constructibles
d'ensembles, de groupes, etc.

\begin{definition}
\label{definition-constructible}
Soit $X$ un schéma.
Soit $\Lambda$ un anneau noethérien. Un faisceau de $\Lambda$-modules
sur $X_\proetale$ est {\it constructible} si, pour tout ouvert affine
$U \subset X$, il existe une décomposition finie
de $U$ en sous-schémas localement fermés constructibles
$U = \coprod_i U_i$ telle que
$\mathcal{F}|_{U_i}$ soit de type fini et localement constant pour tout $i$.
\end{definition}

\noindent
Là encore, cela ne donne rien de ``nouveau''.

\begin{lemma}
\label{lemma-compare-constructible}
Soit $X$ un schéma. Soit $\Lambda$ un anneau noethérien.
Le foncteur $\epsilon^{-1}$ définit une équivalence de catégories
$$
\left\{
\begin{matrix}
\text{faisceaux constructibles de}\\
\Lambda\text{-modules sur }X_\etale\\
\end{matrix}
\right\}
\longleftrightarrow
\left\{
\begin{matrix}
\text{faisceaux constructibles de}\\
\Lambda\text{-modules sur }X_\proetale\\
\end{matrix}
\right\}
$$
entre les faisceaux constructibles de $\Lambda$-modules sur $X_\etale$
et les faisceaux constructibles de $\Lambda$-modules sur $X_\proetale$.
\end{lemma}

\begin{proof}
D'après le lemme \ref{lemma-fully-faithful}, le foncteur $\epsilon^{-1}$
est pleinement fidèle et commute à l'image inverse, c'est-à-dire à la restriction aux
strates. Ainsi, l'image par $\epsilon^{-1}$ d'un faisceau \'etale
constructible est un faisceau pro-\'etale constructible. Pour achever la
démonstration, soit $\mathcal{F}$ un faisceau constructible de $\Lambda$-modules
sur $X_\proetale$ comme dans la définition \ref{definition-constructible}.
Il existe un morphisme canonique
$$
\epsilon^{-1}\epsilon_*\mathcal{F} \longrightarrow \mathcal{F}
$$
Nous allons montrer que ce morphisme est un isomorphisme. Cela prouvera que
$\mathcal{F}$ appartient à l'image essentielle de $\epsilon^{-1}$
et achèvera la démonstration (nous omettons les détails).

\medskip\noindent
Puisqu'il suffit de le démontrer localement sur $X$, on peut supposer que $X$ est
quasi-compact et quasi-séparé et que l'on dispose
d'une partition finie $X = \coprod_{i = 1, \ldots, n} X_i$
en strates localement fermées constructibles telles que $\mathcal{F}|_{X_i}$
soit localement constant de type fini. Nous raisonnons par récurrence sur $n$.
Le cas initial $n = 1$ résulte du lemme \ref{lemma-compare-locally-constant}.
Prenons un point $x \in X$; alors $x \in X_k$ pour un certain $k$.
Il suffit de montrer que le morphisme affiché est un isomorphisme dans un
voisinage ouvert de $x$. On peut donc supposer que $X_k$ est fermé dans $X$.
Posons $Z = X_k$ et notons $U \subset X$ le complément de $Z$.
Observons que l'hypothèse de récurrence s'applique à la restriction de
$\mathcal{F}$ à $Z$ et à $U$. D'après le lemme \ref{lemma-ses-associated-to-open},
on a une suite exacte courte
$$
0 \to j_!j^{-1}\mathcal{F} \to \mathcal{F} \to i_*i^{-1}\mathcal{F} \to 0
$$
sur $X_\proetale$. La fonctorialité donne un diagramme commutatif
$$
\xymatrix{
0 \ar[r] &
\epsilon^{-1}\epsilon_*j_!j^{-1}\mathcal{F} \ar[r] \ar[d] &
\epsilon^{-1}\epsilon_*\mathcal{F} \ar[r] \ar[d] &
\epsilon^{-1}\epsilon_*i_*i^{-1}\mathcal{F} \ar[r] \ar[d] & 0 \\
0 \ar[r] &
j_!j^{-1}\mathcal{F} \ar[r] &
\mathcal{F} \ar[r] &
i_*i^{-1}\mathcal{F} \ar[r] & 0
}
$$
Par récurrence, on sait, d'une part, que
$\epsilon^{-1}\epsilon_*i^{-1}\mathcal{F} \to i^{-1}\mathcal{F}$
et
$\epsilon^{-1}\epsilon_*j^{-1}\mathcal{F} \to j^{-1}\mathcal{F}$
sont des isomorphismes et, d'autre part, que
$i^{-1}\mathcal{F} = \epsilon^{-1}\mathcal{A}$
et
$j^{-1}\mathcal{F} = \epsilon^{-1}\mathcal{B}$
pour certains faisceaux constructibles de $\Lambda$-modules
$\mathcal{A}$ sur $Z_\etale$ et $\mathcal{B}$ sur $U_\etale$.
Alors
$$
\epsilon^{-1}\epsilon_*j_!j^{-1}\mathcal{F} =
\epsilon^{-1}\epsilon_*j_!\epsilon^{-1}\mathcal{B} =
\epsilon^{-1}\epsilon_*\epsilon^{-1}j_!\mathcal{B} =
\epsilon^{-1}j_!\mathcal{B} =
j_!\epsilon^{-1}\mathcal{B} =
j_!j^{-1}\mathcal{F}
$$
où la deuxième égalité résulte du lemme \ref{lemma-jshriek-comparison},
la troisième du lemme \ref{lemma-fully-faithful}, et la
quatrième, à nouveau, du lemme \ref{lemma-jshriek-comparison}.
De même, on a
$$
\epsilon^{-1}\epsilon_*i_*i^{-1}\mathcal{F} =
\epsilon^{-1}\epsilon_*i_*\epsilon^{-1}\mathcal{A} =
\epsilon^{-1}\epsilon_*\epsilon^{-1}i_*\mathcal{A} =
\epsilon^{-1}i_*\mathcal{A} =
i_*\epsilon^{-1}\mathcal{A} =
i_*i^{-1}\mathcal{F}
$$
Cette fois, on utilise le lemme \ref{lemma-morphism-comparison}.
Le lemme des cinq permet de conclure que le
morphisme vertical médian du grand diagramme est un isomorphisme.
\end{proof}

\begin{lemma}
\label{lemma-constructible-serre}
Soit $X$ un schéma. Soit $\Lambda$ un anneau noethérien.
La catégorie des faisceaux constructibles de $\Lambda$-modules sur $X_\proetale$
est une sous-catégorie de Serre faible de $\textit{Mod}(X_\proetale, \Lambda)$.
\end{lemma}

\begin{proof}
C'est une conséquence formelle des
lemmes \ref{lemma-compare-constructible} et \ref{lemma-compare-derived}
et du résultat correspondant pour le site \'etale
(Cohomologie \'etale, lemme \ref{etale-cohomology-lemma-constructible-abelian}).
\end{proof}

\begin{lemma}
\label{lemma-compare-constructible-derived}
Soit $X$ un schéma. Soit $\Lambda$ un anneau noethérien.
Soient $D_c(X_\etale, \Lambda)$, resp.\ $D_c(X_\proetale, \Lambda)$,
les sous-catégories pleines de
$D(X_\etale, \Lambda)$, resp.\ $D(X_\proetale, \Lambda)$,
formées des complexes dont les faisceaux de cohomologie sont des
faisceaux constructibles de $\Lambda$-modules. Alors
$$
\epsilon^{-1} :
D_c^+(X_\etale, \Lambda)
\longrightarrow
D_c^+(X_\proetale, \Lambda)
$$
est une équivalence de catégories.
\end{lemma}

\begin{proof}
Les catégories $D_c(X_\etale, \Lambda)$ et $D_c(X_\proetale, \Lambda)$
sont des sous-catégories strictement pleines, saturées et triangulées de
$D(X_\etale, \Lambda)$ et $D(X_\proetale, \Lambda)$ d'après
Cohomologie \'etale, lemme \ref{etale-cohomology-lemma-constructible-abelian}
et
le lemme \ref{lemma-constructible-serre}
et
Catégories dérivées, section \ref{derived-section-triangulated-sub}.
L'énoncé du lemme s'obtient en combinant les
lemmes \ref{lemma-compare-derived} et
\ref{lemma-compare-constructible}.
\end{proof}

\begin{lemma}
\label{lemma-tensor-c}
Soit $X$ un schéma. Soit $\Lambda$ un anneau noethérien.
Soient $K, L \in D_c^-(X_\proetale, \Lambda)$. Alors
$K \otimes_\Lambda^\mathbf{L} L$ appartient à $D_c^-(X_\proetale, \Lambda)$.
\end{lemma}

\begin{proof}
Remarquons que $H^i(K \otimes_\Lambda^\mathbf{L} L)$ s'identifie à
$H^i(\tau_{\geq i - 1}K \otimes_\Lambda^\mathbf{L} \tau_{\geq i - 1}L)$.
On peut donc supposer que $K$ et $L$ sont bornés.
Dans ce cas, le lemme \ref{lemma-compare-constructible-derived} permet de
se ramener au cas du site \'etale, voir
Cohomologie \'etale, lemme \ref{etale-cohomology-lemma-tensor-c}.
\end{proof}

\begin{lemma}
\label{lemma-compare-truncations-constructible}
Soit $X$ un schéma. Soit $\Lambda$ un anneau noethérien.
Soit $I \subset \Lambda$ un idéal.
Soit $K$ un objet de $D(X_\proetale, \Lambda)$.
Posons $K_n = K \otimes_\Lambda^\mathbf{L} \underline{\Lambda/I^n}$.
Si $K_1$ appartient à $D^-_c(X_\proetale, \Lambda/I)$, alors
$K_n$ appartient à $D^-_c(X_\proetale, \Lambda/I^n)$ pour tout $n$.
\end{lemma}

\begin{proof}
Considérons les triangles distingués
$$
K \otimes_\Lambda^\mathbf{L} \underline{I^n/I^{n + 1}} \to
K_{n + 1} \to
K_n \to
K \otimes_\Lambda^\mathbf{L} \underline{I^n/I^{n + 1}}[1]
$$
et les isomorphismes
$$
K \otimes_\Lambda^\mathbf{L} \underline{I^n/I^{n + 1}} =
K_1 \otimes_{\Lambda/I}^\mathbf{L} \underline{I^n/I^{n + 1}}
$$
D'après le lemme \ref{lemma-tensor-c}, ce produit tensoriel possède des
faisceaux de cohomologie constructibles (et nuls lorsque la cohomologie de $K_1$
est nulle). Une récurrence sur $n$ utilisant le
lemme \ref{lemma-constructible-serre}
montre donc que chaque $K_n$ possède des faisceaux de cohomologie constructibles.
\end{proof}








\section{Faisceaux adiques constructibles}
\label{section-adic}

\noindent
Dans cette section, nous définissons la notion de faisceau
constructible de $\Lambda$-modules, ainsi que quelques variantes.

\begin{definition}
\label{definition-adic}
Soit $\Lambda$ un anneau noethérien et soit $I \subset \Lambda$ un idéal.
Soit $X$ un schéma. Soit $\mathcal{F}$ un faisceau de $\Lambda$-modules
sur $X_\proetale$.
\begin{enumerate}
\item Nous disons que $\mathcal{F}$ est un {\it faisceau constructible de $\Lambda$-modules}
si $\mathcal{F} = \lim \mathcal{F}/I^n\mathcal{F}$ et si chaque
$\mathcal{F}/I^n\mathcal{F}$ est un faisceau constructible de $\Lambda/I^n$-modules.
\item Si $\mathcal{F}$ est un faisceau constructible de $\Lambda$-modules, nous disons
que $\mathcal{F}$ est {\it lisse} si chaque $\mathcal{F}/I^n\mathcal{F}$ est
localement constant.
\item Nous disons que $\mathcal{F}$ est {\it lisse adique}\footnote{Cette notation n'est
peut-être pas standard.} s'il existe un module
complet pour la topologie $I$-adique, muni d'une structure de $\Lambda$-module, noté $M$, tel que $M/IM$ soit fini
et que $\mathcal{F}$ soit localement isomorphe à
$$
\underline{M}^\wedge = \lim \underline{M/I^nM}.
$$
\item Nous disons que $\mathcal{F}$ est
{\it adique constructible}\footnote{Cette notation n'est peut-être pas standard.}
si, pour tout ouvert affine $U \subset X$,
il existe une décomposition $U = \coprod U_i$ en
sous-schémas localement fermés constructibles tels que $\mathcal{F}|_{U_i}$
soit lisse adique.
\end{enumerate}
\end{definition}

\noindent
La définition d'un faisceau constructible de $\Lambda$-modules équivaut
à celle de \cite[Expos\'e VI, d\'efinition 1.1.1]{SGA5} lorsque
$\Lambda = \mathbf{Z}_\ell$ et $I = (\ell)$. Il est clair que
l'on a les implications
$$
\xymatrix{
\text{lisse adique} \ar@{=>}[r] \ar@{=>}[d] &
\text{adique constructible} \ar@{=>}[d] \\
\text{faisceau constructible lisse de }\Lambda\text{-modules} \ar@{=>}[r] &
\text{faisceau constructible de }\Lambda\text{-modules}
}
$$
Les flèches verticales peuvent être inversées dans certains cas
(voir les lemmes \ref{lemma-Noetherian-constructible} et
\ref{lemma-Noetherian-adic-constructible}). En général,
ni la catégorie des faisceaux adiques constructibles, ni
la catégorie des faisceaux constructibles de $\Lambda$-modules ne sont stables
par noyaux et conoyaux.

\medskip\noindent
En effet, soit $X$ un schéma affine dont l'espace topologique sous-jacent $|X|$
est homéomorphe à $\Lambda = \mathbf{Z}_\ell$; voir
l'exemple \ref{example-construct-space}. Notons
$f : |X| \to \mathbf{Z}_\ell = \Lambda$
un homéomorphisme. Nous pouvons considérer $f$ comme une section de
$\underline{\Lambda}^\wedge$ sur $X$; la multiplication par $f$
définit alors un complexe à deux termes
$$
\underline{\Lambda}^\wedge \xrightarrow{f} \underline{\Lambda}^\wedge
$$
sur $X_\proetale$. Le faisceau $\underline{\Lambda}^\wedge$ est lisse adique.
Cependant, le conoyau du morphisme ci-dessus n'est pas adique constructible, car
le type d'isomorphisme de ses fibres prend une infinité
de valeurs : $\mathbf{Z}/\ell^n\mathbf{Z}$ et $\mathbf{Z}_\ell$.
Le conoyau est un faisceau constructible de $\mathbf{Z}_\ell$-modules.
Cependant, le noyau n'est même pas un faisceau constructible de $\mathbf{Z}_\ell$-modules,
car il est nul sur un ouvert non quasi-compact, mais n'est pas nul.

\begin{lemma}
\label{lemma-Noetherian-constructible}
Soit $X$ un schéma noethérien. Soit $\Lambda$ un anneau noethérien et
soit $I \subset \Lambda$ un idéal. Soit $\mathcal{F}$ un
faisceau constructible de $\Lambda$-modules sur $X_\proetale$.
Il existe alors une partition finie $X = \coprod X_i$ en
sous-schémas localement fermés tels que la restriction $\mathcal{F}|_{X_i}$
soit lisse.
\end{lemma}

\begin{proof}
Posons $R = \bigoplus I^n/I^{n + 1}$. Observons que $R$ est un anneau noethérien.
Puisque chacun des faisceaux
$\mathcal{F}/I^n\mathcal{F}$ est un faisceau constructible de
$\Lambda/I^n\Lambda$-modules, $I^n\mathcal{F}/I^{n + 1}\mathcal{F}$
est également un faisceau constructible de $\Lambda/I$-modules, donc l'image inverse
d'un faisceau constructible $\mathcal{G}_n$ sur $X_\etale$, d'après le
lemme \ref{lemma-compare-constructible}.
Posons $\mathcal{G} = \bigoplus \mathcal{G}_n$. C'est un faisceau
de $R$-modules sur $X_\etale$, et le morphisme
$$
\mathcal{G}_0 \otimes_{\Lambda/I} \underline{R}
\longrightarrow
\mathcal{G}
$$
est surjectif parce que les morphismes
$$
\mathcal{F}/I\mathcal{F} \otimes \underline{I^n/I^{n + 1}} \to
I^n\mathcal{F}/I^{n + 1}\mathcal{F}
$$
sont surjectifs. Ainsi, $\mathcal{G}$ est un faisceau constructible de
$R$-modules d'après Cohomologie \'etale, proposition
\ref{etale-cohomology-proposition-constructible-over-noetherian}.
Choisissons une partition $X = \coprod X_i$ telle que
$\mathcal{G}|_{X_i}$ soit un faisceau localement constant de $R$-modules
de type fini (Cohomologie \'etale, lemme
\ref{etale-cohomology-lemma-constructible-quasi-compact-quasi-separated}).
Nous affirmons que cette partition possède la propriété de l'énoncé.
En effet, en remplaçant $X$ par $X_i$, nous pouvons supposer que $\mathcal{G}$ est localement
constant. Il s'ensuit que chacun des faisceaux
 $I^n\mathcal{F}/I^{n + 1}\mathcal{F}$
est localement constant. À l'aide des suites exactes courtes
$$
0 \to I^n\mathcal{F}/I^{n + 1}\mathcal{F} \to
\mathcal{F}/I^{n + 1}\mathcal{F} \to \mathcal{F}/I^n\mathcal{F} \to 0
$$
et, par récurrence, du lemme de Modules sur les sites
\ref{sites-modules-lemma-kernel-finite-locally-constant}
on obtient le résultat.
\end{proof}

\begin{lemma}
\label{lemma-weird}
Soit $X$ un schéma affine faiblement contractile. Soit $\Lambda$ un anneau noethérien,
soit $I \subset \Lambda$ un idéal et soit $\mathcal{F}$ un faisceau de
$\Lambda$-modules sur $X_\proetale$ tel que
\begin{enumerate}
\item $\mathcal{F} = \lim \mathcal{F}/I^n\mathcal{F}$,
\item $\mathcal{F}/I^n\mathcal{F}$ soit un faisceau constant de
$\Lambda/I^n$-modules,
\item $\mathcal{F}/I\mathcal{F}$ soit de type fini.
\end{enumerate}
Alors $\mathcal{F} \cong \underline{M}^\wedge$, où $M$ est
un module de type fini sur $\Lambda^\wedge$.
\end{lemma}

\begin{proof}
Choisissons un $\Lambda/I^n$-module $M_n$ tel que
$\mathcal{F}/I^n\mathcal{F} \cong \underline{M_n}$.
Puisque nous disposons des surjections
$\mathcal{F}/I^{n + 1}\mathcal{F} \to \mathcal{F}/I^n\mathcal{F}$,
nous concluons qu'il existe des
surjections $M_{n + 1} \to M_n$ induisant des isomorphismes
$M_{n + 1}/I^nM_{n + 1} \to M_n$. Fixons un choix de telles surjections
et posons $M = \lim M_n$. Alors $M$, complet pour la topologie $I$-adique,
est un $\Lambda$-module vérifiant $M/I^nM = M_n$; voir le
lemme Algèbre, \ref{algebra-lemma-limit-complete}.
Puisque $M_1$ est un $\Lambda$-module de type fini
(Modules sur les sites, lemme
\ref{sites-modules-lemma-locally-constant-finite-type}),
nous voyons que $M$ est un module de type fini sur $\Lambda^\wedge$.
Considérons les faisceaux
$$
\mathcal{I}_n = \mathit{Isom}(\underline{M_n}, \mathcal{F}/I^n\mathcal{F})
$$
sur $X_\proetale$. Le passage au quotient par $I^n$ définit un morphisme de transition
$$
\mathcal{I}_{n + 1} \longrightarrow \mathcal{I}_n
$$
D'après notre choix de $M_n$, le faisceau $\mathcal{I}_n$ est un torseur sous
$$
\mathit{Isom}(\underline{M_n}, \underline{M_n}) =
\underline{\text{Isom}_\Lambda(M_n, M_n)}
$$
(Modules sur les sites, lemme \ref{sites-modules-lemma-locally-constant}),
puisque $\mathcal{F}/I^n\mathcal{F}$ est localement isomorphe pour la topologie \'etale
à $\underline{M_n}$. D'après
Compléments d'algèbre, lemme \ref{more-algebra-lemma-hom-systems-ML},
que le système de faisceaux $(\mathcal{I}_n)$ vérifie la condition de Mittag-Leffler.
Pour tout $n$, soit $\mathcal{I}'_n \subset \mathcal{I}_n$ l'image
de $\mathcal{I}_N \to \mathcal{I}_n$ pour tout $N \gg n$.
Alors
$$
\ldots \to \mathcal{I}'_3 \to \mathcal{I}'_2 \to \mathcal{I}'_1 \to *
$$
est une suite de faisceaux d'ensembles sur $X_\proetale$ dont les morphismes de
transition sont surjectifs. Puisque $*(X)$ est réduit à un élément (et donc non vide)
et que l'évaluation en $X$ transforme les morphismes surjectifs de faisceaux d'ensembles
en surjections d'ensembles, nous pouvons choisir
$s \in \lim \mathcal{I}'_n(X)$. Ces sections définissent des isomorphismes
$\underline{M}^\wedge \to \lim \mathcal{F}/I^n\mathcal{F} = \mathcal{F}$,
ce qui achève la démonstration.
\end{proof}

\begin{lemma}
\label{lemma-connected-lisse}
Soit $X$ un schéma connexe. Soit $\Lambda$ un anneau noethérien et soit
$I \subset \Lambda$ un idéal. Si $\mathcal{F}$ est un faisceau constructible lisse
de $\Lambda$-modules sur $X_\proetale$, alors $\mathcal{F}$
est lisse adique.
\end{lemma}

\begin{proof}
D'après le lemme \ref{lemma-compare-locally-constant}, on a
$\mathcal{F}/I^n\mathcal{F} = \epsilon^{-1}\mathcal{G}_n$
pour un certain faisceau localement constant $\mathcal{G}_n$ de $\Lambda/I^n$-modules. D'après
le lemme de Cohomologie \'etale
\ref{etale-cohomology-lemma-connected-locally-constant},
il existe un module de type fini sur $\Lambda/I^n$, noté $M_n$, tel que
$\mathcal{G}_n$ soit localement isomorphe à $\underline{M_n}$.
Choisissons un recouvrement $\{W_t \to X\}_{t \in T}$ où chaque $W_t$ est
affine et faiblement contractile.
Alors $\mathcal{F}|_{W_t}$ satisfait aux hypothèses du
lemme \ref{lemma-weird},
et donc $\mathcal{F}|_{W_t} \cong \underline{N_t}^\wedge$
pour un module de type fini sur $\Lambda^\wedge$, noté $N_t$. Remarquons que
$N_t/I^nN_t \cong M_n$ pour tous $t$ et $n$. Ainsi,
$N_t \cong N_{t'}$ pour tous $t, t' \in T$; voir le
lemme de Compléments d'algèbre \ref{more-algebra-lemma-isomorphic-completions}.
Cela démontre que $\mathcal{F}$ est lisse adique.
\end{proof}

\begin{lemma}
\label{lemma-Noetherian-adic-constructible}
Soit $X$ un schéma noethérien. Soit $\Lambda$ un anneau noethérien et
soit $I \subset \Lambda$ un idéal. Soit $\mathcal{F}$ un
faisceau constructible de $\Lambda$-modules sur $X_\proetale$. Alors $\mathcal{F}$
est adique constructible.
\end{lemma}

\begin{proof}
C'est une conséquence des lemmes \ref{lemma-Noetherian-constructible} et
\ref{lemma-connected-lisse}, du fait qu'un schéma noethérien
est localement connexe
(Topologie, lemme \ref{topology-lemma-locally-Noetherian-locally-connected}),
et des définitions.
\end{proof}

\noindent
Il sera utile d'identifier les faisceaux constructibles de $\Lambda$-modules
dans la catégorie des faisceaux de $\Lambda$-modules complets au sens dérivé.
Il se trouve que l'analogue naïf de
Compléments d'algèbre, lemme \ref{more-algebra-lemma-derived-complete-finite},
est faux dans ce contexte. Cependant, voici
l'analogue de Compléments d'algèbre, lemme
\ref{more-algebra-lemma-derived-complete-zero}.

\begin{lemma}
\label{lemma-derived-complete-zero}
Soit $X$ un schéma. Soit $\Lambda$ un anneau et soit
$I \subset \Lambda$ un idéal de type fini.
Soit $\mathcal{F}$ un faisceau de $\Lambda$-modules sur $X_\proetale$.
Si $\mathcal{F}$ est complet au sens dérivé et si $\mathcal{F}/I\mathcal{F} = 0$,
alors $\mathcal{F} = 0$.
\end{lemma}

\begin{proof}
Supposons que $\mathcal{F}/I\mathcal{F}$ soit nul.
Écrivons $I = (f_1, \ldots, f_r)$. Soit $i < r$ le plus grand
entier tel que $\mathcal{G} = \mathcal{F}/(f_1, \ldots, f_i)\mathcal{F}$
soit non nul. Si un tel $i$ n'existe pas, alors $\mathcal{F} = 0$, ce que nous
voulions montrer. Alors $\mathcal{G}$ est complet au sens dérivé, comme conoyau
d'un morphisme entre modules complets au sens dérivé; voir la
proposition \ref{proposition-enough-weakly-contractibles}.
D'après le choix de $i$, le morphisme $f_{i + 1} : \mathcal{G} \to \mathcal{G}$
est surjectif. Ainsi,
$$
\lim (\ldots \to \mathcal{G} \xrightarrow{f_{i + 1}} \mathcal{G}
\xrightarrow{f_{i + 1}} \mathcal{G})
$$
est non nul, ce qui contredit le fait que $\mathcal{G}$ soit complet au sens dérivé.
\end{proof}

\begin{lemma}
\label{lemma-derived-complete-limit}
Soit $X$ un schéma affine faiblement contractile.
Soit $\Lambda$ un anneau noethérien et soit $I \subset \Lambda$ un idéal.
Soit $\mathcal{F}$ un faisceau de $\Lambda$-modules complet au sens dérivé
sur $X_\proetale$, tel que $\mathcal{F}/I\mathcal{F}$ soit un faisceau localement
constant de $\Lambda/I$-modules de type fini.
Alors il existe un entier $t$ et un morphisme surjectif
$$
(\underline{\Lambda}^\wedge)^{\oplus t} \to \mathcal{F}
$$
\end{lemma}

\begin{proof}
Puisque $X$ est faiblement contractile, il existe un recouvrement fini par des ouverts disjoints
$X = \coprod U_i$ tel que $\mathcal{F}/I\mathcal{F}|_{U_i}$
soit isomorphe au faisceau constant associé à un module de type fini sur $\Lambda/I$
$M_i$. Choisissons un nombre fini de générateurs $m_{ij}$ de $M_i$. Nous
pouvons trouver des sections $s_{ij} \in \mathcal{F}(X)$ dont les restrictions sont les
$m_{ij}$ considérés comme sections de $\mathcal{F}/I\mathcal{F}$ sur $U_i$. 
Soit $t$ le nombre total des $s_{ij}$. Nous obtenons alors un morphisme
$$
\alpha : \underline{\Lambda}^{\oplus t} \longrightarrow \mathcal{F}
$$
qui est surjectif modulo $I$ par construction. D'après le
lemme \ref{lemma-naive-completion},
la complétion dérivée de $\underline{\Lambda}^{\oplus t}$ est le
faisceau $(\underline{\Lambda}^\wedge)^{\oplus t}$. Puisque $\mathcal{F}$
est complet au sens dérivé, on voit que $\alpha$ se factorise par un morphisme
$$
\alpha^\wedge :
(\underline{\Lambda}^\wedge)^{\oplus t}
\longrightarrow
\mathcal{F}
$$
Alors $\mathcal{Q} = \Coker(\alpha^\wedge)$ est un faisceau de
$\Lambda$-modules complet au sens dérivé d'après la
proposition \ref{proposition-enough-weakly-contractibles}.
Par construction, $\mathcal{Q}/I\mathcal{Q} = 0$. Il résulte du
lemme \ref{lemma-derived-complete-zero}
que $\mathcal{Q} = 0$, ce que nous voulions montrer.
\end{proof}








\section{Une catégorie dérivée appropriée}
\label{section-suitable-derived}

\noindent
Soit $X$ un schéma. Il s'avérera que, pour de nombreux schémas $X$, on peut obtenir une
catégorie dérivée appropriée de faisceaux $\ell$-adiques en considérant les
objets $K$ de $D(X_\proetale, \Lambda)$ complets au sens dérivé tels
que $K \otimes_\Lambda^\mathbf{L} \mathbf{F}_\ell$ soit borné et possède des
faisceaux de cohomologie constructibles. Voici la définition générale.

\begin{definition}
\label{definition-Dbc}
Soit $\Lambda$ un anneau noethérien et soit $I \subset \Lambda$ un idéal.
Soit $X$ un schéma. Un objet $K$ de $D(X_\proetale, \Lambda)$ est dit
{\it constructible} si
\begin{enumerate}
\item $K$ est complet au sens dérivé relativement à $I$,
\item $K \otimes_\Lambda^\mathbf{L} \underline{\Lambda/I}$
possède des faisceaux de cohomologie constructibles et a localement une dimension de Tor finie.
\end{enumerate}
Nous notons $D_{cons}(X, \Lambda)$ la sous-catégorie pleine formée des objets constructibles
$K$ de $D(X_\proetale, \Lambda)$.
\end{definition}

\noindent
Rappelons qu'avec nos conventions, un complexe de dimension de Tor finie
est borné (Cohomologie sur les sites, définition
\ref{sites-cohomology-definition-tor-amplitude}).
En fait, rassemblons dans un lemme tout ce que nous avons démontré jusqu'ici.

\begin{lemma}
\label{lemma-describe-constructible-complexes}
Dans la situation ci-dessus, supposons que $K$ appartienne à $D_{cons}(X, \Lambda)$
et que $X$ soit quasi-compact. Posons
$K_n = K \otimes_\Lambda^\mathbf{L} \underline{\Lambda/I^n}$.
Il existe $a, b$ tels que
\begin{enumerate}
\item $K = R\lim K_n$ et $H^i(K) = 0$ pour $i \not \in [a, b]$,
\item chaque $K_n$ est d'amplitude de Tor dans $[a, b]$,
\item chaque $K_n$ possède des faisceaux de cohomologie constructibles,
\item pour tout indice, $K_n = \epsilon^{-1}L_n$ pour un certain
$L_n \in D_{ctf}(X_\etale, \Lambda/I^n)$
(Cohomologie \'etale, définition \ref{etale-cohomology-definition-ctf}).
\end{enumerate}
\end{lemma}

\begin{proof}
Puisque $K$ est complet au sens dérivé, nous avons $K = R\lim K_n$ d'après le
lemme \ref{lemma-naive-completion}.
Par définition de la finitude locale de la dimension de Tor, nous pouvons trouver
$a, b$ tels que $K_1$ soit d'amplitude de Tor dans $[a, b]$.
La partie (2) résulte de
Cohomologie sur les sites, lemme \ref{sites-cohomology-lemma-bounded}.
Ensuite, (1) résulte du fait que $K$ est complet au sens dérivé, par la description
des limites dans
Cohomologie sur les sites, proposition
\ref{sites-cohomology-proposition-enough-weakly-contractibles},
et du fait que $H^b(K_{n + 1}) \to H^b(K_n)$ est surjective,
puisque $K_n = K_{n + 1} \otimes^\mathbf{L}_\Lambda \underline{\Lambda/I^n}$.
La partie (3) résulte du
lemme \ref{lemma-compare-truncations-constructible}.
La partie (4) résulte du
lemme \ref{lemma-compare-constructible-derived}
et du fait que $L_n$ est de dimension de Tor finie parce que $K_n$ l'est
(nous omettons un court argument).
\end{proof}

\begin{lemma}
\label{lemma-local-structure-constructible-complex}
Soit $X$ un schéma affine faiblement contractile. Soit $\Lambda$ un
anneau noethérien et soit $I \subset \Lambda$ un idéal. Soit $K$ un objet de
$D_{cons}(X, \Lambda)$ tel que les faisceaux de cohomologie de
$K \otimes_\Lambda^\mathbf{L} \underline{\Lambda/I}$ soient localement
constants. Alors il existe un recouvrement fini par des ouverts disjoints
$X = \coprod U_i$ et, pour chaque $i$, une famille finie de
modules projectifs de type fini sur $\Lambda^\wedge$, notés $M^a, \ldots, M^b$,
telle que $K|_{U_i}$ soit représenté par un complexe
$$
(\underline{M^a})^\wedge \to \ldots \to (\underline{M^b})^\wedge
$$
dans $D(U_{i, \proetale}, \Lambda)$, pour certains morphismes de faisceaux de
$\Lambda$-modules $(\underline{M^i})^\wedge \to (\underline{M^{i + 1}})^\wedge$.
\end{lemma}

\begin{proof}
Nous utilisons librement les résultats du
lemme \ref{lemma-describe-constructible-complexes}.
Choisissons $a, b$ comme dans ce lemme. Nous allons démontrer le lemme par
récurrence sur $b - a$. Posons $\mathcal{F} = H^b(K)$.
Remarquons que $\mathcal{F}$ est un faisceau de
$\Lambda$-modules complet au sens dérivé, d'après la
proposition \ref{proposition-enough-weakly-contractibles}.
De plus, $\mathcal{F}/I\mathcal{F}$ est un faisceau localement
constant de $\Lambda/I$-modules de type fini.
Appliquons le lemme \ref{lemma-derived-complete-limit} pour obtenir une surjection
$\rho : (\underline{\Lambda}^\wedge)^{\oplus t} \to \mathcal{F}$.

\medskip\noindent
Si $a = b$, alors $K = \mathcal{F}[-b]$. Dans ce cas, nous voyons que
$$
\mathcal{F} \otimes_\Lambda^\mathbf{L} \underline{\Lambda/I} =
\mathcal{F}/I\mathcal{F}
$$
Comme $X$ est faiblement contractile et que $\mathcal{F}/I\mathcal{F}$ est
localement constant, nous pouvons trouver une décomposition finie
$X = \coprod U_i$ en ouverts affines disjoints $U_i$
et des $\Lambda/I$-modules $\overline{M}_i$
tels que $\mathcal{F}/I\mathcal{F}$ se restreigne à
$\underline{\overline{M}_i}$ sur $U_i$. Après avoir raffiné le recouvrement,
nous pouvons supposer que le morphisme
$$
\rho|_{U_i} \bmod I :
\underline{\Lambda/I}^{\oplus t}
\longrightarrow
\underline{\overline{M}_i}
$$
est égal à $\underline{\alpha_i}$ pour un certain morphisme surjectif de modules
$\alpha_i : \Lambda/I^{\oplus t} \to \overline{M}_i$ ; voir
Modules sur les sites, lemme \ref{sites-modules-lemma-morphism-locally-constant}.
Remarquons que chaque $\overline{M}_i$ est un module de type fini sur $\Lambda/I$.
Puisque $\mathcal{F}/I\mathcal{F}$ est d'amplitude de Tor dans $[0, 0]$,
nous en concluons que $\overline{M}_i$ est un $\Lambda/I$-module plat.
Ainsi, $\overline{M}_i$ est projectif de type fini
(Algèbre, lemme \ref{algebra-lemma-finite-projective}).
Nous pouvons donc trouver un projecteur
$\overline{p}_i : (\Lambda/I)^{\oplus t} \to (\Lambda/I)^{\oplus t}$
dont l'image est envoyée isomorphiquement sur $\overline{M}_i$ par le morphisme $\alpha_i$.
Nous pouvons relever $\overline{p}_i$ en un projecteur
$p_i : (\Lambda^\wedge)^{\oplus t} \to
(\Lambda^\wedge)^{\oplus t}$\footnote{Démonstration : d'après
Algèbre, lemme \ref{algebra-lemma-lift-idempotents-noncommutative},
nous pouvons relever $\overline{p}_i$ en un système compatible de
projecteurs $p_{i, n} : (\Lambda/I^n)^{\oplus t} \to (\Lambda/I^n)^{\oplus t}$,
puis poser $p_i = \lim p_{i, n}$, ce qui convient parce que
$\Lambda^\wedge = \lim \Lambda/I^n$.}.
Alors $M_i = \Im(p_i)$ est complet pour la topologie $I$-adique et de type fini
comme $\Lambda^\wedge$-module ; de plus, $M_i/IM_i = \overline{M}_i$.
Sur $U_i$, considérons les morphismes
$$
\underline{M_i}^\wedge \to
(\underline{\Lambda}^\wedge)^{\oplus t} \to
\mathcal{F}|_{U_i}
$$
Par construction, leur composé induit un isomorphisme modulo $I$.
La source et le but sont complets au sens dérivé ; il en va donc de même du conoyau
$\mathcal{Q}$ et du noyau $\mathcal{K}$. Nous avons
$\mathcal{Q}/I\mathcal{Q} = 0$ par construction ; ainsi, $\mathcal{Q}$
est nul d'après le lemme \ref{lemma-derived-complete-zero}.
Alors
$$
0 \to \mathcal{K}/I\mathcal{K} \to
\underline{\overline{M}_i}
\to \mathcal{F}/I\mathcal{F} \to 0
$$
est exacte grâce à l'annulation de $\text{Tor}_1$ rappelée au début de ce
paragraphe. On utilise aussi l'expression
$\underline{\Lambda}^\wedge/I\overline{\Lambda}^\wedge$,
ainsi que Modules sur les sites, lemme \ref{sites-modules-lemma-completion-flat},
pour voir que
$\underline{M_i}^\wedge/I\underline{M_i}^\wedge = \underline{\overline{M}_i}$.
Ainsi, $\mathcal{K}/I\mathcal{K} = 0$ par construction, et nous concluons
que $\mathcal{K} = 0$ comme précédemment. Cela démontre le résultat lorsque $a = b$.

\medskip\noindent
Si $b > a$, nous relevons le morphisme $\rho$ en un morphisme
$$
\tilde \rho : (\underline{\Lambda}^\wedge)^{\oplus t}[-b] \longrightarrow K
$$
dans $D(X_\proetale, \Lambda)$. Cela est possible car nous pouvons considérer
$K$ comme un complexe de $\underline{\Lambda}^\wedge$-modules, d'après la
discussion dans l'introduction de la
section \ref{section-derived-completion-noetherian},
et parce que $X_\proetale$ est faiblement contractile ;
il n'y a donc aucune obstruction à relever les éléments
$\rho(e_s) \in H^0(X, \mathcal{F})$ en des éléments de $H^b(X, K)$.
En insérant $\tilde \rho$ dans un triangle distingué
$$
(\underline{\Lambda}^\wedge)^{\oplus t}[-b] \to K \to L \to
(\underline{\Lambda}^\wedge)^{\oplus t}[-b + 1]
$$
nous voyons que $L$ est un objet de $D_{cons}(X, \Lambda)$ tel
que $L \otimes_\Lambda^\mathbf{L} \underline{\Lambda/I}$
soit d'amplitude de Tor contenue dans $[a, b - 1]$ (nous omettons les détails).
Par récurrence, nous pouvons décrire $L$ localement comme dans l'énoncé du lemme ; disons que
$L$ est isomorphe à
$$
(\underline{M^a})^\wedge \to \ldots \to (\underline{M^{b - 1}})^\wedge
$$
Le morphisme
$L \to (\underline{\Lambda}^\wedge)^{\oplus t}[-b + 1]$
correspond à un morphisme
$(\underline{M^{b - 1}})^\wedge \to (\underline{\Lambda}^\wedge)^{\oplus t}$,
qui nous permet de prolonger le complexe d'un terme. Le complexe correspondant
est isomorphe à $K$ dans la catégorie dérivée, par les propriétés
des catégories triangulées. Cela achève la démonstration.
\end{proof}

\noindent
Motivés par ce qui se passe pour les faisceaux constructibles de $\Lambda$-modules,
nous introduisons la notion suivante.

\begin{definition}
\label{definition-adic-constructible}
Soit $X$ un schéma. Soit $\Lambda$ un anneau noethérien et soit
$I \subset \Lambda$ un idéal. Soit $K \in D(X_\proetale, \Lambda)$.
\begin{enumerate}
\item Nous disons que $K$ est {\it lisse adique}\footnote{Cette notation n'est peut-être
pas standard.} s'il existe un complexe borné de modules
projectifs de type fini sur $\Lambda^\wedge$, noté $M^\bullet$, tel que
$K$ soit localement isomorphe à
$$
\underline{M^a}^\wedge \to \ldots \to \underline{M^b}^\wedge
$$
\item Nous disons que $K$ est {\it adique constructible}\footnote{Cette notation n'est peut-être
pas standard.} si, pour tout ouvert affine $U \subset X$,
il existe une décomposition $U = \coprod U_i$ en
sous-schémas localement fermés constructibles telle que $K|_{U_i}$
soit lisse adique.
\end{enumerate}
\end{definition}

\noindent
La différence entre la structure locale obtenue dans le
lemme \ref{lemma-local-structure-constructible-complex}
et la structure d'un complexe lisse adique tient à ce que
les morphismes $\underline{M^i}^\wedge \to \underline{M^{i + 1}}^\wedge$ du
lemme \ref{lemma-local-structure-constructible-complex}
ne sont pas nécessairement constants, tandis que, dans la définition ci-dessus, ils
doivent être constants.

\begin{lemma}
\label{lemma-weakly-contractible-locally-constant-ML}
Soit $X$ un schéma affine faiblement contractile. Soit $\Lambda$ un
anneau noethérien et soit $I \subset \Lambda$ un idéal. Soit $K$ un objet de
$D_{cons}(X, \Lambda)$ tel que
$K \otimes_\Lambda^\mathbf{L} \underline{\Lambda/I^n}$
soit isomorphe dans $D(X_\proetale, \Lambda/I^n)$ à un
complexe de faisceaux constants de $\Lambda/I^n$-modules. Alors
$$
H^0(X, K \otimes_\Lambda^\mathbf{L} \Lambda/I^n)
$$
satisfait à la condition de Mittag-Leffler.
\end{lemma}

\begin{proof}
Supposons que $K \otimes_\Lambda^\mathbf{L} \underline{\Lambda/I^n}$ soit isomorphe
à $\underline{E_n}$ pour un objet $E_n$ de $D(\Lambda/I^n)$.
Puisque $K \otimes_\Lambda^\mathbf{L} \underline{\Lambda/I}$ a une
dimension de Tor finie et possède des faisceaux de cohomologie de type fini,
on voit que $E_1$ est parfait (voir
Compléments d'algèbre, lemme \ref{more-algebra-lemma-perfect}). Les morphismes de transition
$$
K \otimes_\Lambda^\mathbf{L} \underline{\Lambda/I^{n + 1}}
\to
K \otimes_\Lambda^\mathbf{L} \underline{\Lambda/I^n}
$$
proviennent localement de morphismes de complexes (éventuellement nombreux et distincts)
$E_{n + 1} \to E_n$ dans $D(\Lambda/I^{n + 1})$ ; voir
Cohomologie sur les sites, lemme \ref{sites-cohomology-lemma-locally-constant-map}.
Pour chaque $n$, choisissons un tel morphisme et observons qu'il induit
un isomorphisme
$E_{n + 1} \otimes_{\Lambda/I^{n + 1}}^\mathbf{L} \Lambda/I^n \to E_n$
dans $D(\Lambda/I^n)$. D'après
Compléments d'algèbre, lemme \ref{more-algebra-lemma-Rlim-perfect-gives-complete}
il existe un complexe borné $M^\bullet$ de modules projectifs de type fini sur
$\Lambda^\wedge$ et des isomorphismes $M^\bullet/I^nM^\bullet \to E_n$
dans $D(\Lambda/I^n)$ compatibles avec les morphismes de transition.

\medskip\noindent
Observons maintenant que, pour toute famille finie d'indices
$n > m > k$, le triplet de morphismes
$$
H^0(X, K \otimes_\Lambda^\mathbf{L} \Lambda/I^n)
\to
H^0(X, K \otimes_\Lambda^\mathbf{L} \Lambda/I^m)
\to
H^0(X, K \otimes_\Lambda^\mathbf{L} \Lambda/I^k)
$$
est isomorphe à
$$
H^0(X, \underline{M^\bullet/I^nM^\bullet})
\to
H^0(X, \underline{M^\bullet/I^mM^\bullet})
\to
H^0(X, \underline{M^\bullet/I^kM^\bullet})
$$
En effet, tout isomorphisme
$$
\underline{M^\bullet/I^nM^\bullet} \to
K \otimes_\Lambda^\mathbf{L} \Lambda/I^n
$$
induit des isomorphismes analogues modulo $I^m$ et $I^k$; l'assertion est donc
vraie. Ainsi, pour démontrer le lemme, il suffit de montrer que
le système
$H^0(X, \underline{M^\bullet/I^nM^\bullet})$ vérifie la condition de Mittag-Leffler.
Comme la prise des sections sur $X$ est exacte, il suffit de montrer que
le système de $\Lambda$-modules
$$
H^0(M^\bullet/I^nM^\bullet)
$$
vérifie la condition de Mittag-Leffler. Posons $A = \Lambda^\wedge$ et considérons la suite
spectrale
$$
\text{Tor}_{-p}^A(H^q(M^\bullet), A/I^nA) \Rightarrow
H^{p + q}(M^\bullet/I^nM^\bullet)
$$
D'après Compléments d'algèbre, lemme \ref{more-algebra-lemma-tor-strictly-pro-zero}
les pro-systèmes $\{\text{Tor}_{-p}^A(H^q(M^\bullet), A/I^nA)\}$ sont nuls
pour $p > 0$. Le pro-système $\{H^0(M^\bullet/I^nM^\bullet)\}$
est donc égal au pro-système $\{H^0(M^\bullet)/I^nH^0(M^\bullet)\}$
et le lemme est démontré.
\end{proof}

\begin{lemma}
\label{lemma-connected-adic-lisse}
Soit $X$ un schéma connexe. Soit $\Lambda$ un anneau noethérien et soit
$I \subset \Lambda$ un idéal. Si $K$ appartient à $D_{cons}(X, \Lambda)$
et si $K \otimes_\Lambda \underline{\Lambda/I}$
a des faisceaux de cohomologie localement constants, alors $K$ est lisse adique
(définition \ref{definition-adic-constructible}).
\end{lemma}

\begin{proof}
Posons $K_n = K \otimes_\Lambda^\mathbf{L} \underline{\Lambda/I^n}$.
Nous utiliserons les résultats du lemme \ref{lemma-describe-constructible-complexes}
sans autre mention. D'après Cohomologie sur les sites, lemme
\ref{sites-cohomology-lemma-locally-constant-bounded}
on voit que $K_n$ a des faisceaux de cohomologie localement constants pour tout $n$.
On a $K_n = \epsilon^{-1}L_n$ pour un certain $L_n$ appartenant à
$D_{ctf}(X_\etale, \Lambda/I^n)$ à faisceaux de cohomologie localement constants.
D'après Cohomologie étale, lemme
\ref{etale-cohomology-lemma-connected-ctf-locally-constant}
il existe des complexes parfaits $M_n \in D(\Lambda/I^n)$
tels que $L_n$ soit localement isomorphe à $\underline{M_n}$ pour la topologie \'etale.
Les morphismes $L_{n + 1} \to L_n$ correspondant à $K_{n + 1} \to K_n$
induisent des isomorphismes
$L_{n + 1} \otimes_{\Lambda/I^{n + 1}}^\mathbf{L} \underline{\Lambda/I^n}
\to L_n$. En raisonnant localement sur $X$, on conclut qu'il
existe des morphismes $M_{n + 1} \to M_n$ dans $D(\Lambda/I^{n + 1})$
induisant des isomorphismes
$M_{n + 1} \otimes_{\Lambda/I^{n + 1}} \Lambda/I^n \to M_n$ ; voir
Cohomologie sur les sites, lemme \ref{sites-cohomology-lemma-locally-constant-map}.
Fixons un choix de tels morphismes. D'après
Compléments d'algèbre, lemme \ref{more-algebra-lemma-Rlim-perfect-gives-complete}
il existe un complexe borné $M^\bullet$ de modules projectifs de type fini sur
$\Lambda^\wedge$ et des isomorphismes $M^\bullet/I^nM^\bullet \to M_n$
dans $D(\Lambda/I^n)$ compatibles avec les morphismes de transition.
Pour achever la démonstration, nous allons montrer que $K$ est localement isomorphe
à
$$
\underline{M^\bullet}^\wedge =
\lim \underline{M^\bullet/I^nM^\bullet} =
R\lim \underline{M^\bullet/I^nM^\bullet}
$$
Soit $E^\bullet$ le complexe dual de $M^\bullet$ ; voir
Compléments d'algèbre, lemme \ref{more-algebra-lemma-dual-perfect-complex}
ainsi que sa démonstration. Considérons les objets
$$
H_n =
R\SheafHom_{\Lambda/I^n}(\underline{M^\bullet/I^nM^\bullet}, K_n) =
\underline{E^\bullet/I^nE^\bullet} \otimes_{\Lambda/I^n}^\mathbf{L} K_n
$$
de $D(X_\proetale, \Lambda/I^n)$. La réduction modulo $I^n$ définit un
morphisme de transition $H_{n + 1} \to H_n$. Posons $H = R\lim H_n$. Alors $H$ est un
objet de $D_{cons}(X, \Lambda)$ (nous omettons les détails) tel que
$H \otimes_\Lambda^\mathbf{L} \underline{\Lambda/I^n} = H_n$.
Choisissons un recouvrement $\{W_t \to X\}_{t \in T}$ où chaque $W_t$
est affine et faiblement contractile. Le choix de $M^\bullet$
montre que
\begin{align*}
H_n|_{W_t} & \cong 
R\SheafHom_{\Lambda/I^n}(\underline{M^\bullet/I^nM^\bullet},
\underline{M^\bullet/I^nM^\bullet}) \\
& =
\underline{
\text{Tot}(E^\bullet/I^nE^\bullet \otimes_{\Lambda/I^n} M^\bullet/I^nM^\bullet)
}
\end{align*}
Nous pouvons donc appliquer le lemme \ref{lemma-weakly-contractible-locally-constant-ML}
à $H = R\lim H_n$. On conclut que le système $H^0(W_t, H_n)$ vérifie la condition de
Mittag-Leffler. Puisque, pour tout $n \gg 1$, il existe un élément de $H^0(W_t, H_n)$
dont l'image est un isomorphisme dans
$$
H^0(W_t, H_1) = \Hom(\underline{M^\bullet/IM^\bullet}, K_1)
$$
on obtient dans la limite inverse un élément $(\varphi_{t, n})$
qui fournit un isomorphisme modulo $I$. Alors
$$
R\lim \varphi_{t, n} :
\underline{M^\bullet}^\wedge|_{W_t} =
R\lim \underline{M^\bullet/I^nM^\bullet}|_{W_t}
\longrightarrow
R\lim K_n|_{W_t} = K|_{W_t}
$$
est un isomorphisme. Cela achève la démonstration.
\end{proof}

\begin{proposition}
\label{proposition-Noetherian-adic-constructible}
Soit $X$ un schéma noethérien. Soit $\Lambda$ un anneau noethérien et
soit $I \subset \Lambda$ un idéal. Soit $K$ un objet de
$D_{cons}(X, \Lambda)$. Alors $K$ est adique constructible
(définition \ref{definition-adic-constructible}).
\end{proposition}

\begin{proof}
C'est une conséquence du lemme \ref{lemma-connected-adic-lisse}
ainsi que du fait qu'un schéma noethérien est localement connexe
(Topologie, lemme \ref{topology-lemma-locally-Noetherian-locally-connected}),
et des définitions.
\end{proof}







\section{Changement de base propre}
\label{section-proper-base-change}

\noindent
Dans cette section, nous expliquons comment démontrer le théorème de changement de base propre
pour les objets complets au sens dérivé sur le site pro-\'etale, à l'aide du théorème de
changement de base propre en cohomologie \'etale, conformément au principe général
selon lequel nous n'utilisons la topologie pro-\'etale que pour traiter les « problèmes de limite »
et faisons appel aux résultats établis pour la topologie \'etale afin de traiter tout
le reste.

\begin{theorem}
\label{theorem-proper-base-change}
Supposons que $f : X \to Y$ soit un morphisme propre de schémas et que $g : Y' \to Y$ soit
un morphisme de schémas donnant lieu au diagramme de changement de base
$$
\xymatrix{
X' \ar[r]_{g'} \ar[d]_{f'} & X \ar[d]^f \\
Y' \ar[r]^g & Y
}
$$
Soit $\Lambda$ un anneau noethérien et soit $I \subset \Lambda$ un idéal
tel que $\Lambda/I$ soit de torsion. Soit $K$ un objet
de $D(X_\proetale)$ tel que
\begin{enumerate}
\item $K$ est complet au sens dérivé ;
\item $K \otimes_\Lambda^\mathbf{L} \underline{\Lambda/I^n}$ est
borné inférieurement, à faisceaux de cohomologie provenant de $X_\etale$ ;
\item $\Lambda/I^n$ est un $\Lambda$-module parfait\footnote{Cette hypothèse
peut être supprimée si $K$ est un complexe constructible, voir \cite{BS}.}.
\end{enumerate}
Alors le morphisme de changement de base
$$
Lg_{comp}^*Rf_*K \longrightarrow Rf'_*L(g')^*_{comp}K
$$
est un isomorphisme.
\end{theorem}

\begin{proof}
Nous omettons la construction du morphisme de changement de base (elle n'utilise que
les propriétés formelles de l'image directe dérivée et de l'image inverse dérivée complétée ;
comparer avec
Cohomologie sur les sites, remarque \ref{sites-cohomology-remark-base-change}).
Posons $K_n = K \otimes^\mathbf{L}_\Lambda \underline{\Lambda/I^n}$.
D'après le lemme \ref{lemma-naive-completion}, on a $K = R\lim K_n$
parce que $K$ est complet au sens dérivé.
Les lemmes \ref{lemma-pushforward-Noetherian-case} et
\ref{lemma-naive-completion} permettent de réécrire le membre de gauche :
$$
Lg_{comp}^* Rf_* K =
R\lim Lg^*(Rf_*K)\otimes^\mathbf{L}_\Lambda \underline{\Lambda/I^n} =
R\lim Lg^* Rf_* K_n
$$
La dernière égalité résulte du fait que $\Lambda/I^n$ est un module parfait et de
la formule de projection (Cohomologie sur les sites, lemme
\ref{sites-cohomology-lemma-projection-formula}).
Le lemme \ref{lemma-pushforward-Noetherian-case} permet de réécrire le membre de
droite :
$$
Rf'_* L(g')^*_{comp} K =
Rf'_* R\lim L(g')^* K_n  =
R\lim Rf'_* L(g')^* K_n
$$
La dernière égalité résulte du fait que $Rf'_*$ commute à $R\lim$
(Cohomologie sur les sites, lemme
\ref{sites-cohomology-lemma-Rf-commutes-with-Rlim}).
Il suffit donc de montrer que les morphismes
$$
Lg^* Rf_* K_n \longrightarrow Rf'_* L(g')^* K_n
$$
sont des isomorphismes. D'après le lemme \ref{lemma-compare-derived} et la deuxième
condition, on peut écrire $K_n = \epsilon^{-1}L_n$ pour un certain
$L_n \in D^+(X_\etale, \Lambda/I^n)$. D'après le lemme \ref{lemma-morphism-comparison}
et le fait que $\epsilon^{-1}$ commute aux images inverses,
on obtient
$$
Lg^* Rf_* K_n =
Lg^* Rf_* \epsilon^*L_n =
Lg^* \epsilon^{-1} Rf_* L_n =
\epsilon^{-1} Lg^* Rf_* L_n
$$
et
$$
Rf'_* L(g')^* K_n =
Rf'_* L(g')^* \epsilon^{-1} L_n =
Rf'_* \epsilon^{-1} L(g')^* L_n =
\epsilon^{-1} Rf'_* L(g')^* L_n
$$
(on utilise également ici que $L_n$ est borné inférieurement).
Enfin, d'après le théorème de changement de base propre en cohomologie \'etale
(Cohomologie étale, théorème
\ref{etale-cohomology-theorem-proper-base-change}), on a
$$
Lg^* Rf_* L_n = Rf'_* L(g')^* L_n
$$
(en utilisant de nouveau que $L_n$ est borné inférieurement),
ce qui démontre le théorème.
\end{proof}






\section{Changement d'univers partiel}
\label{section-change-universe}

\noindent
Nous conseillons au lecteur de passer cette section : nous montrons ici que la cohomologie
des faisceaux pour la topologie pro-\'etale est indépendante du
choix de l'univers partiel. Plus précisément, le foncteur $g_*$ du
lemme \ref{lemma-proetale-cohomology-independent-partial-universe} ci-dessous
est un plongement de petits topos pro-\'etales qui ne modifie pas la cohomologie.
Pour les grands sites pro-\'etales, les lemmes \ref{lemma-change-alpha} et
\ref{lemma-cohomology-enlarge-partial-universe} disent essentiellement la
même chose.

\medskip\noindent
Mais d'abord, 
comme promis dans la section \ref{section-proetale}, démontrons que la topologie d'un
grand site pro-\'etale $\Sch_\proetale$ est, en un certain sens, induite par
la topologie pro-\'etale sur la catégorie de tous les schémas.

\begin{lemma}
\label{lemma-proetale-induced}
Soit $\Sch_\proetale$ un grand site pro-\'etale comme dans la
définition \ref{definition-big-proetale-site}.
Soit $T \in \Ob(\Sch_\proetale)$.
Soit $\{T_i \to T\}_{i \in I}$ un recouvrement pro-\'etale quelconque de $T$.
Il existe un recouvrement $\{U_j \to T\}_{j \in J}$ de $T$ dans le site
$\Sch_\proetale$ qui raffine $\{T_i \to T\}_{i \in I}$.
\end{lemma}

\begin{proof}
Prenons d'abord un recouvrement $\{V_k \to T\}$ comme dans le
lemme \ref{lemma-get-many-weakly-contractible}.
Alors les recouvrements pro-\'etales $\{T_i \times_T V_k \to V_k\}$
peuvent chacun être raffinés par un recouvrement fini par des ouverts disjoints
$V_k = V_{k, 1} \amalg \ldots \amalg V_{k, n_k}$, voir le
lemme \ref{lemma-w-contractible-proetale-cover}.
Alors $\{V_{k, i} \to T\}$ constitue, dans $\Sch_\proetale$, un recouvrement de l'objet considéré
qui raffine $\{T_i \to T\}_{i \in I}$.
\end{proof}

\noindent
Énonçons et démontrons d'abord la comparaison pour les petits sites
pro-\'etales. Notons que nous n'affirmons pas que le petit topos pro-\'etale
d'un schéma soit indépendant du choix de l'univers partiel ; cela est faux,
contrairement au cas du petit topos \'etale
(Cohomologie étale, lemme
\ref{etale-cohomology-lemma-etale-topos-independent-partial-universe}).

\begin{lemma}
\label{lemma-proetale-cohomology-independent-partial-universe}
Soit $S$ un schéma. Soient $S_\proetale \subset S_\proetale'$ deux
petits sites pro-\'etales de $S$ construits dans la
définition \ref{definition-big-small-proetale}. Alors le foncteur d'inclusion
vérifie les hypothèses de
Sites, lemme \ref{sites-lemma-bigger-site}.
Il existe donc des morphismes de topos
$$
\xymatrix{
\Sh(S_\proetale) \ar[r]^g &
\Sh(S_\proetale') \ar[r]^f &
\Sh(S_\proetale)
}
$$
dont la composée est isomorphe à l'identité et qui vérifient $f_* = g^{-1}$.
De plus,
\begin{enumerate}
\item pour $\mathcal{F}' \in \textit{Ab}(S_\proetale')$, on a
$H^p(S_\proetale', \mathcal{F}') = H^p(S_\proetale, g^{-1}\mathcal{F}')$,
\item pour $\mathcal{F} \in \textit{Ab}(S_\proetale)$, on a
$$
H^p(S_\proetale, \mathcal{F}) =
H^p(S_\proetale', g_*\mathcal{F}) =
H^p(S_\proetale', f^{-1}\mathcal{F}).
$$
\end{enumerate}
\end{lemma}

\begin{proof}
Le foncteur d'inclusion est pleinement fidèle et continu.
Nous avons vu que $S_\proetale$ et $S_\proetale'$ admettent des produits fibrés
et des objets finaux, et que notre foncteur commute à ces constructions
(lemme \ref{lemma-fibre-products-proetale}).
Il résulte du lemme \ref{lemma-proetale-induced}
que le foncteur d'inclusion est cocontinu.
L'existence de $f$ et $g$ résulte donc de
Sites, lemme \ref{sites-lemma-bigger-site}.
L'égalité de (1) résulte de
Cohomologie sur les sites, lemme \ref{sites-cohomology-lemma-cohomology-bigger-site}.
L'assertion (2) découle de (1), car
$\mathcal{F} = g^{-1}g_*\mathcal{F} = g^{-1}f^{-1}\mathcal{F}$.
\end{proof}

\noindent
Nous démontrons ensuite un résultat analogue pour les grands topos pro-\'etales.

\begin{lemma}
\label{lemma-change-alpha}
Supposons donnés des grands sites $\Sch_\proetale$ et $\Sch'_\proetale$ comme dans la
définition \ref{definition-big-proetale-site}.
Supposons que $\Sch_\proetale$ soit contenu dans $\Sch'_\proetale$.
Le foncteur d'inclusion $\Sch_\proetale \to \Sch'_\proetale$ vérifie
les hypothèses de Sites, lemme \ref{sites-lemma-bigger-site}.
Il existe des morphismes de topos
\begin{eqnarray*}
g : \Sh(\Sch_\proetale) &
\longrightarrow &
\Sh(\Sch'_\proetale) \\
f : \Sh(\Sch'_\proetale) &
\longrightarrow &
\Sh(\Sch_\proetale)
\end{eqnarray*}
tels que $f \circ g \cong \text{id}$. Pour tout objet $S$
de $\Sch_\proetale$, le foncteur d'inclusion
$(\Sch/S)_\proetale \to (\Sch'/S)_\proetale$ vérifie
les hypothèses de Sites, lemme \ref{sites-lemma-bigger-site}.
On obtient donc de même des morphismes
\begin{eqnarray*}
g : \Sh((\Sch/S)_\proetale) &
\longrightarrow &
\Sh((\Sch'/S)_\proetale) \\
f : \Sh((\Sch'/S)_\proetale) &
\longrightarrow &
\Sh((\Sch/S)_\proetale)
\end{eqnarray*}
avec $f \circ g \cong \text{id}$.
\end{lemma}

\begin{proof}
Les hypothèses (b), (c) et (e) du
lemme de Sites \ref{sites-lemma-bigger-site}
sont immédiates pour les foncteurs
$\Sch_\proetale \to \Sch'_\proetale$ et
$(\Sch/S)_\proetale \to (\Sch'/S)_\proetale$. La propriété (a) résulte du
lemme \ref{lemma-proetale-induced}.
La propriété (d) est satisfaite parce que
les produits fibrés existent dans les catégories $\Sch_\proetale$, $\Sch'_\proetale$
et sont compatibles avec les produits fibrés dans la catégorie des schémas.
\end{proof}

\begin{lemma}
\label{lemma-cohomology-enlarge-partial-universe}
Soit $S$ un schéma. Soient $(\Sch/S)_\proetale$ et $(\Sch'/S)_\proetale$ deux
grands sites pro-\'etales de $S$ comme dans la
définition \ref{definition-big-small-proetale}.
Supposons que le premier soit contenu dans
le second. Alors
\begin{enumerate}
\item pour tout faisceau abélien $\mathcal{F}'$ défini sur $(\Sch'/S)_\proetale$
et tout objet $U$ de $(\Sch/S)_\proetale$, on a
$$
H^p(U, \mathcal{F}'|_{(\Sch/S)_\proetale}) =
H^p(U, \mathcal{F}')
$$
Autrement dit, la cohomologie de $\mathcal{F}'$ sur $U$ calculée dans le site le plus grand
coïncide avec la cohomologie de la restriction de $\mathcal{F}'$ au site le plus petit
sur $U$.
\item pour tout faisceau abélien $\mathcal{F}$ sur $(\Sch/S)_\proetale$, il existe un
faisceau abélien $\mathcal{F}'$ sur $(\Sch/S)_\proetale'$ dont la restriction à
$(\Sch/S)_\proetale$ est isomorphe à $\mathcal{F}$.
\end{enumerate}
\end{lemma}

\begin{proof}
D'après le lemme \ref{lemma-change-alpha}, le foncteur d'inclusion
$(\Sch/S)_\proetale \to (\Sch'/S)_\proetale$ vérifie les hypothèses de
Sites, lemme \ref{sites-lemma-bigger-site}. Cela implique (2), tandis que (1)
résulte de
Cohomologie sur les sites, lemme \ref{sites-cohomology-lemma-cohomology-bigger-site}.
\end{proof}










\begin{multicols}{2}[\section{Autres chapitres}]
\noindent
Préliminaires
\begin{enumerate}
\item \hyperref[introduction-section-phantom]{Introduction}
\item \hyperref[conventions-section-phantom]{Conventions}
\item \hyperref[sets-section-phantom]{Théorie des ensembles}
\item \hyperref[categories-section-phantom]{Catégories}
\item \hyperref[topology-section-phantom]{Topologie}
\item \hyperref[sheaves-section-phantom]{Faisceaux sur les espaces}
\item \hyperref[sites-section-phantom]{Sites et faisceaux}
\item \hyperref[stacks-section-phantom]{Champs}
\item \hyperref[fields-section-phantom]{Corps}
\item \hyperref[algebra-section-phantom]{Algèbre commutative}
\item \hyperref[brauer-section-phantom]{Groupes de Brauer}
\item \hyperref[homology-section-phantom]{Algèbre homologique}
\item \hyperref[derived-section-phantom]{Catégories dérivées}
\item \hyperref[simplicial-section-phantom]{Méthodes simpliciales}
\item \hyperref[more-algebra-section-phantom]{Compléments d'algèbre}
\item \hyperref[smoothing-section-phantom]{Lissage des morphismes d'anneaux}
\item \hyperref[modules-section-phantom]{Faisceaux de modules}
\item \hyperref[sites-modules-section-phantom]{Modules sur les sites}
\item \hyperref[injectives-section-phantom]{Objets injectifs}
\item \hyperref[cohomology-section-phantom]{Cohomologie des faisceaux}
\item \hyperref[sites-cohomology-section-phantom]{Cohomologie sur les sites}
\item \hyperref[dga-section-phantom]{Algèbre différentielle graduée}
\item \hyperref[dpa-section-phantom]{Algèbre à puissances divisées}
\item \hyperref[sdga-section-phantom]{Faisceaux différentiels gradués}
\item \hyperref[hypercovering-section-phantom]{Hyperrecouvrements}
\end{enumerate}
Schémas
\begin{enumerate}
\setcounter{enumi}{25}
\item \hyperref[schemes-section-phantom]{Schémas}
\item \hyperref[constructions-section-phantom]{Constructions de schémas}
\item \hyperref[properties-section-phantom]{Propriétés des schémas}
\item \hyperref[morphisms-section-phantom]{Morphismes de schémas}
\item \hyperref[coherent-section-phantom]{Cohomologie des schémas}
\item \hyperref[divisors-section-phantom]{Diviseurs}
\item \hyperref[limits-section-phantom]{Limites de schémas}
\item \hyperref[varieties-section-phantom]{Variétés}
\item \hyperref[topologies-section-phantom]{Topologies sur les schémas}
\item \hyperref[descent-section-phantom]{Descente}
\item \hyperref[perfect-section-phantom]{Catégories dérivées de schémas}
\item \hyperref[more-morphisms-section-phantom]{Compléments sur les morphismes}
\item \hyperref[flat-section-phantom]{Compléments sur la platitude}
\item \hyperref[groupoids-section-phantom]{Schémas en groupoïdes}
\item \hyperref[more-groupoids-section-phantom]{Compléments sur les schémas en groupoïdes}
\item \hyperref[etale-section-phantom]{Morphismes étales de schémas}
\end{enumerate}
Thèmes de théorie des schémas
\begin{enumerate}
\setcounter{enumi}{41}
\item \hyperref[chow-section-phantom]{Homologie de Chow}
\item \hyperref[intersection-section-phantom]{Théorie de l'intersection}
\item \hyperref[pic-section-phantom]{Schémas de Picard des courbes}
\item \hyperref[weil-section-phantom]{Théories cohomologiques de Weil}
\item \hyperref[adequate-section-phantom]{Modules adéquats}
\item \hyperref[dualizing-section-phantom]{Complexes dualisants}
\item \hyperref[duality-section-phantom]{Dualité pour les schémas}
\item \hyperref[discriminant-section-phantom]{Discriminants et différentes}
\item \hyperref[derham-section-phantom]{Cohomologie de de Rham}
\item \hyperref[local-cohomology-section-phantom]{Cohomologie locale}
\item \hyperref[algebraization-section-phantom]{Géométrie algébrique et formelle}
\item \hyperref[curves-section-phantom]{Courbes algébriques}
\item \hyperref[resolve-section-phantom]{Résolution des surfaces}
\item \hyperref[models-section-phantom]{Réduction semi-stable}
\item \hyperref[functors-section-phantom]{Foncteurs et morphismes}
\item \hyperref[equiv-section-phantom]{Catégories dérivées de variétés}
\item \hyperref[pione-section-phantom]{Groupes fondamentaux de schémas}
\item \hyperref[etale-cohomology-section-phantom]{Cohomologie étale}
\item \hyperref[crystalline-section-phantom]{Cohomologie cristalline}
\item \hyperref[proetale-section-phantom]{Cohomologie pro-étale}
\item \hyperref[relative-cycles-section-phantom]{Cycles relatifs}
\item \hyperref[more-etale-section-phantom]{Compléments de cohomologie étale}
\item \hyperref[trace-section-phantom]{La formule des traces}
\end{enumerate}
Espaces algébriques
\begin{enumerate}
\setcounter{enumi}{64}
\item \hyperref[spaces-section-phantom]{Espaces algébriques}
\item \hyperref[spaces-properties-section-phantom]{Propriétés des espaces algébriques}
\item \hyperref[spaces-morphisms-section-phantom]{Morphismes d'espaces algébriques}
\item \hyperref[decent-spaces-section-phantom]{Espaces algébriques décents}
\item \hyperref[spaces-cohomology-section-phantom]{Cohomologie des espaces algébriques}
\item \hyperref[spaces-limits-section-phantom]{Limites d'espaces algébriques}
\item \hyperref[spaces-divisors-section-phantom]{Diviseurs sur les espaces algébriques}
\item \hyperref[spaces-over-fields-section-phantom]{Espaces algébriques sur des corps}
\item \hyperref[spaces-topologies-section-phantom]{Topologies sur les espaces algébriques}
\item \hyperref[spaces-descent-section-phantom]{Descente et espaces algébriques}
\item \hyperref[spaces-perfect-section-phantom]{Catégories dérivées d'espaces}
\item \hyperref[spaces-more-morphisms-section-phantom]{Compléments sur les morphismes d'espaces}
\item \hyperref[spaces-flat-section-phantom]{Platitude sur les espaces algébriques}
\item \hyperref[spaces-groupoids-section-phantom]{Groupoïdes dans les espaces algébriques}
\item \hyperref[spaces-more-groupoids-section-phantom]{Compléments sur les groupoïdes dans les espaces}
\item \hyperref[bootstrap-section-phantom]{Amorçage}
\item \hyperref[spaces-pushouts-section-phantom]{Sommes amalgamées d'espaces algébriques}
\end{enumerate}
Thèmes de géométrie
\begin{enumerate}
\setcounter{enumi}{81}
\item \hyperref[spaces-chow-section-phantom]{Groupes de Chow des espaces}
\item \hyperref[groupoids-quotients-section-phantom]{Quotients de groupoïdes}
\item \hyperref[spaces-more-cohomology-section-phantom]{Compléments sur la cohomologie des espaces}
\item \hyperref[spaces-simplicial-section-phantom]{Espaces simpliciaux}
\item \hyperref[spaces-duality-section-phantom]{Dualité pour les espaces}
\item \hyperref[formal-spaces-section-phantom]{Espaces algébriques formels}
\item \hyperref[restricted-section-phantom]{Algébrisation des espaces formels}
\item \hyperref[spaces-resolve-section-phantom]{Résolution des surfaces revisitée}
\end{enumerate}
Théorie des déformations
\begin{enumerate}
\setcounter{enumi}{89}
\item \hyperref[formal-defos-section-phantom]{Théorie formelle des déformations}
\item \hyperref[defos-section-phantom]{Théorie des déformations}
\item \hyperref[cotangent-section-phantom]{Le complexe cotangent}
\item \hyperref[examples-defos-section-phantom]{Problèmes de déformation}
\end{enumerate}
Champs algébriques
\begin{enumerate}
\setcounter{enumi}{93}
\item \hyperref[algebraic-section-phantom]{Champs algébriques}
\item \hyperref[examples-stacks-section-phantom]{Exemples de champs}
\item \hyperref[stacks-sheaves-section-phantom]{Faisceaux sur les champs algébriques}
\item \hyperref[criteria-section-phantom]{Critères de représentabilité}
\item \hyperref[artin-section-phantom]{Axiomes d'Artin}
\item \hyperref[quot-section-phantom]{Espaces de Quot et de Hilbert}
\item \hyperref[stacks-properties-section-phantom]{Propriétés des champs algébriques}
\item \hyperref[stacks-morphisms-section-phantom]{Morphismes de champs algébriques}
\item \hyperref[stacks-limits-section-phantom]{Limites de champs algébriques}
\item \hyperref[stacks-cohomology-section-phantom]{Cohomologie des champs algébriques}
\item \hyperref[stacks-perfect-section-phantom]{Catégories dérivées de champs}
\item \hyperref[stacks-introduction-section-phantom]{Introduction aux champs algébriques}
\item \hyperref[stacks-more-morphisms-section-phantom]{Compléments sur les morphismes de champs}
\item \hyperref[stacks-geometry-section-phantom]{La géométrie des champs}
\end{enumerate}
Thèmes de théorie des espaces de modules
\begin{enumerate}
\setcounter{enumi}{107}
\item \hyperref[moduli-section-phantom]{Champs de modules}
\item \hyperref[moduli-curves-section-phantom]{Espaces de modules de courbes}
\end{enumerate}
Divers
\begin{enumerate}
\setcounter{enumi}{109}
\item \hyperref[examples-section-phantom]{Exemples}
\item \hyperref[exercises-section-phantom]{Exercices}
\item \hyperref[guide-section-phantom]{Guide bibliographique}
\item \hyperref[desirables-section-phantom]{Desiderata}
\item \hyperref[coding-section-phantom]{Style de programmation}
\item \hyperref[obsolete-section-phantom]{Obsolète}
\item \hyperref[fdl-section-phantom]{Licence de documentation libre GNU}
\item \hyperref[index-section-phantom]{Index généré automatiquement}
\end{enumerate}
\end{multicols}


\bibliography{my}
\bibliographystyle{amsalpha}

\end{document}
